% Équations, inéquations polynômiales et irrationnelles — Mathématiques 1re Tech — Cours signé OnBuch+
% Programme officiel MINESEC. Compiler avec : tectonic N07-equations-inequations-polynomiales-irrationnelles.tex
\newcommand{\DOCMATIERE}{Mathématiques}
\newcommand{\DOCNIVEAU}{1re Tech}
\newcommand{\DOCMODULE}{Mathématiques – classe de Première technique}
\newcommand{\DOCLECON}{N07}
\newcommand{\DOCTITRE}{Équations, inéquations polynômiales et irrationnelles}
% OnBuch+ — préambule « cours signés » ENRICHI (compilé avec tectonic / XeLaTeX)
% Système visuel « Pop » d'OnBuch+ : fond crème (jamais blanc), bordures encre
% épaisses, ombres « dures » décalées, titres Archivo Black en capitales.
% Version MATHÉMATIQUES : graphes (pgfplots/tikz), boîtes pédagogiques
% (définition, propriété, méthode, attention, le savais-tu, exercice résolu,
% exercice, TP, bilan), page de garde et signature OnBuch+ ancrée en bas.
%
% Macros attendues dans main.tex AVANT \input{preamble} :
%   \DOCMATIERE  \DOCNIVEAU  \DOCMODULE  \DOCLECON (numéro)  \DOCTITRE
\documentclass[11pt]{article}

\usepackage{fontspec}
\usepackage[french]{babel}
\usepackage[a4paper,margin=2cm,top=3.4cm,bottom=2.8cm,headheight=1.4cm,footskip=1.3cm]{geometry}
\usepackage{amsmath,amssymb,amsfonts}
\usepackage{mathtools} % \xmapsto, \coloneqq... souvent utilisés par les modèles
\usepackage{cancel} % simplifications barrées
% \abs{z} (module, valeur absolue) et \norm{u} (norme), utilisés par les modèles
\providecommand{\abs}{}\let\abs\relax\DeclarePairedDelimiter{\abs}{\lvert}{\rvert}
\providecommand{\norm}{}\let\norm\relax\DeclarePairedDelimiter{\norm}{\lVert}{\rVert}
\usepackage[table]{xcolor} % [table] : fournit aussi \rowcolors (alternance de lignes)
\usepackage{pagecolor}
\usepackage{tikz}
\usetikzlibrary{arrows.meta,positioning,calc,shapes.geometric,shapes.misc,shapes.arrows,decorations.pathmorphing,decorations.pathreplacing,decorations.markings,patterns,shadows,mindmap,trees,fit,backgrounds,matrix,intersections,angles,quotes}
\usepackage{pgfplots}
\usepackage{pgf-pie} % diagrammes circulaires \pie (statistiques)
\pgfplotsset{compat=1.18}
\usepgfplotslibrary{fillbetween,groupplots} % aires entre courbes (calcul intégral)
\usepackage{enumitem}
\usepackage{fancyhdr}
% [most] charge toutes les bibliothèques tcolorbox (listings, minted, raster,
% documentation...) et gonfle inutilement la mémoire TeX sur un document long
% (~300 boîtes) : on ne charge que ce qui est réellement utilisé.
\usepackage[skins,breakable]{tcolorbox}
\usepackage{graphicx}
\usepackage{colortbl}
\usepackage{booktabs}
\usepackage{tabularx}
\usepackage{diagbox} % case d'en-tête diagonale des tableaux à double entrée
% Colonnes extensibles centrée / gauche / droite (C, L, R), souvent utilisées par les modèles
\newcolumntype{C}{>{\centering\arraybackslash}X}
\newcolumntype{L}{>{\raggedright\arraybackslash}X}
\newcolumntype{R}{>{\raggedleft\arraybackslash}X}
\usepackage{array}
\usepackage{multicol}
\usepackage{siunitx}
% parse-numbers=false : accepte aussi \SI{2,5 \times 10^{-3}}{...}
\sisetup{locale=FR,output-decimal-marker={,},group-minimum-digits=4,parse-numbers=false}
\DeclareSIUnit\ohm{\ensuremath{\symup{\Omega}}} % U+2126 absent de la police
\usepackage{qrcode}
% \degree : commande fréquemment utilisée par les modèles pour un angle ou une
% température en dehors de \SI{}{} (non fournie par les paquets chargés).
\providecommand{\degree}{^{\circ}}
% \qed (fin de démonstration), utilisé par les modèles sans amsthm
\providecommand{\qed}{\hfill$\square$}
% \barre : double barre des valeurs interdites dans un tableau de variations (tkz-tab)
\providecommand{\barre}{\ensuremath{\|}}
% environnement proof (amsthm non chargé) utilisé par les modèles
\ifdefined\proof\else\newenvironment{proof}[1][Démonstration]{\par\noindent\textit{#1.}\ }{\qed\par}\fi
\usepackage{lastpage}
\usepackage{hyperref}
\hypersetup{hidelinks}

% ============================================================
% Palette « Pop » OnBuch+ — mêmes hex que la classe P (plus_ui.dart)
% ============================================================
\definecolor{poporange}{HTML}{F59321}
\definecolor{popdark}{HTML}{E07A0C}
\definecolor{popcream}{HTML}{FFF6E8}
\definecolor{popink}{HTML}{1C1714}
\definecolor{poppurple}{HTML}{7A5AE0}
\definecolor{popgreen}{HTML}{2E9E6A}
\definecolor{poppink}{HTML}{E0457B}
\definecolor{popblue}{HTML}{2F7DE1}
\definecolor{popgold}{HTML}{C98A12}
\definecolor{popmuted}{HTML}{8A7F76}
\definecolor{popline}{HTML}{EADFD2}
\definecolor{popgray}{HTML}{8A7F76}
\colorlet{popgrey}{popgray}
% Alias tolérants (le modèle nomme parfois les couleurs différemment)
\colorlet{popwhite}{white}
\colorlet{popblack}{popink}
\colorlet{popred}{poppink}
\colorlet{popyellow}{popgold}
% Teintes claires pour fonds de tableaux / zones
\colorlet{popblueL}{popblue!10!white}
\colorlet{poporangeL}{poporange!14!white}
\colorlet{popgreenL}{popgreen!12!white}
\colorlet{poppurpleL}{poppurple!12!white}
\colorlet{poppinkL}{poppink!12!white}
\colorlet{popgoldL}{popgold!14!white}

\pagecolor{popcream}

% ============================================================
% Typographie
% ============================================================
\defaultfontfeatures{Ligatures=TeX}
\setmainfont{PlusJakartaSans-Medium.ttf}[Path=../../../fonts/,BoldFont=PlusJakartaSans-Bold.ttf]
% Maths en sans-serif (Fira Math) avec chiffres et lettres droites de Plus
% Jakarta, pour que les formules se fondent dans le texte.
\usepackage[math-style=ISO,bold-style=ISO]{unicode-math}
\setmathfont{FiraMath-Regular.otf}[Path=../../../fonts/]
\setmathfont{PlusJakartaSans-Medium.ttf}[Path=../../../fonts/,range={up/num,up/{Latin,latin},\mathup}]
\usepackage{icomma} % virgule décimale sans espace en mode math
% Caractères mathématiques Unicode absents des polices de texte (Plus Jakarta,
% Archivo Black) : rendus par leur équivalent mathématique, en texte comme en
% formule — sinon ils s'affichent en carré vide (ex. « Arithmétique dans ℤ »).
\usepackage{newunicodechar}
\newunicodechar{ℝ}{\ensuremath{\mathbb{R}}}
\newunicodechar{ℤ}{\ensuremath{\mathbb{Z}}}
\newunicodechar{ℂ}{\ensuremath{\mathbb{C}}}
\newunicodechar{ℕ}{\ensuremath{\mathbb{N}}}
\newunicodechar{ℚ}{\ensuremath{\mathbb{Q}}}
\newunicodechar{𝔻}{\ensuremath{\mathbb{D}}}
\newunicodechar{α}{\ensuremath{\alpha}}
\newunicodechar{β}{\ensuremath{\beta}}
\newunicodechar{γ}{\ensuremath{\gamma}}
\newunicodechar{δ}{\ensuremath{\delta}}
\newunicodechar{ε}{\ensuremath{\varepsilon}}
\newunicodechar{θ}{\ensuremath{\theta}}
\newunicodechar{λ}{\ensuremath{\lambda}}
\newunicodechar{μ}{\ensuremath{\mu}}
\newunicodechar{σ}{\ensuremath{\sigma}}
\newunicodechar{τ}{\ensuremath{\tau}}
\newunicodechar{φ}{\ensuremath{\varphi}}
\newunicodechar{ψ}{\ensuremath{\psi}}
\newunicodechar{ω}{\ensuremath{\omega}}
\newunicodechar{Σ}{\ensuremath{\Sigma}}
% majuscules produites par \MakeUppercase dans les titres (α-aminés -> Α)
\newunicodechar{Α}{\ensuremath{\alpha}}
\newunicodechar{Β}{\ensuremath{\beta}}
\newunicodechar{Γ}{\ensuremath{\Gamma}}
\newunicodechar{Δ}{\ensuremath{\Delta}}
\newunicodechar{Ε}{\ensuremath{\varepsilon}}
\newunicodechar{Θ}{\ensuremath{\Theta}}
\newunicodechar{Λ}{\ensuremath{\Lambda}}
\newunicodechar{Μ}{\ensuremath{\mu}}
\newunicodechar{Τ}{\ensuremath{\tau}}
\newunicodechar{Φ}{\ensuremath{\Phi}}
\newunicodechar{Ψ}{\ensuremath{\Psi}}
\newunicodechar{Ω}{\ensuremath{\Omega}}
\newunicodechar{↦}{\ensuremath{\mapsto}}
\newunicodechar{∈}{\ensuremath{\in}}
\newunicodechar{∉}{\ensuremath{\notin}}
\newunicodechar{∧}{\ensuremath{\wedge}}
\newunicodechar{⇔}{\ensuremath{\Leftrightarrow}}
\newunicodechar{⟺}{\ensuremath{\Longleftrightarrow}}
\newunicodechar{⇒}{\ensuremath{\Rightarrow}}
\newunicodechar{⟹}{\ensuremath{\Longrightarrow}}
\newunicodechar{∘}{\ensuremath{\circ}}
\newunicodechar{∪}{\ensuremath{\cup}}
\newunicodechar{∩}{\ensuremath{\cap}}
\newunicodechar{≡}{\ensuremath{\equiv}}
\newunicodechar{⊂}{\ensuremath{\subset}}
\newunicodechar{⊥}{\ensuremath{\perp}}
\newunicodechar{∃}{\ensuremath{\exists}}
\newunicodechar{∀}{\ensuremath{\forall}}
\newunicodechar{∛}{\ensuremath{\sqrt[3]{\,}}}
\newunicodechar{∜}{\ensuremath{\sqrt[4]{\,}}}
\newunicodechar{⃗}{\ensuremath{{}^{\to}}}
\newunicodechar{ⁿ}{\ifmmode^{n}\else\textsuperscript{\ensuremath{n}}\fi}
\newunicodechar{ˣ}{\ifmmode^{x}\else\textsuperscript{\ensuremath{x}}\fi}
\newunicodechar{ᵇ}{\ifmmode^{b}\else\textsuperscript{\ensuremath{b}}\fi}
\newunicodechar{ʳ}{\ifmmode^{r}\else\textsuperscript{\ensuremath{r}}\fi}
\newunicodechar{ᵘ}{\ifmmode^{u}\else\textsuperscript{\ensuremath{u}}\fi}
\newunicodechar{ʸ}{\ifmmode^{y}\else\textsuperscript{\ensuremath{y}}\fi}
\newunicodechar{ᵃ}{\ifmmode^{a}\else\textsuperscript{\ensuremath{a}}\fi}
\newunicodechar{ᵉ}{\ifmmode^{e}\else\textsuperscript{\ensuremath{e}}\fi}
\newunicodechar{ᵏ}{\ifmmode^{k}\else\textsuperscript{\ensuremath{k}}\fi}
\newunicodechar{ᵖ}{\ifmmode^{p}\else\textsuperscript{\ensuremath{p}}\fi}
\newunicodechar{ᴺ}{\ifmmode^{N}\else\textsuperscript{\ensuremath{N}}\fi}
\newunicodechar{ᶜ}{\ifmmode^{c}\else\textsuperscript{\ensuremath{c}}\fi}
\newunicodechar{ʰ}{\ifmmode^{h}\else\textsuperscript{\ensuremath{h}}\fi}
\newunicodechar{ᵠ}{\ifmmode^{q}\else\textsuperscript{\ensuremath{q}}\fi}
\newunicodechar{ₙ}{\ifmmode_{n}\else\textsubscript{\ensuremath{n}}\fi}
\newunicodechar{ₐ}{\ifmmode_{a}\else\textsubscript{\ensuremath{a}}\fi}
\newunicodechar{ᵢ}{\ifmmode_{i}\else\textsubscript{\ensuremath{i}}\fi}
\newunicodechar{ᵣ}{\ifmmode_{r}\else\textsubscript{\ensuremath{r}}\fi}
\newunicodechar{ₖ}{\ifmmode_{k}\else\textsubscript{\ensuremath{k}}\fi}
\newunicodechar{ᵦ}{\ifmmode_{\beta}\else\textsubscript{\ensuremath{\beta}}\fi}
\newunicodechar{ₓ}{\ifmmode_{x}\else\textsubscript{\ensuremath{x}}\fi}
\newfontfamily\popdisplay{ArchivoBlack-Regular.ttf}[Path=../../../fonts/]
\newfontfamily\poplogo{SpaceGrotesk-Bold.ttf}[Path=../../../fonts/]
\newfontfamily\popsemi{PlusJakartaSans-SemiBold.ttf}[Path=../../../fonts/]
\newfontfamily\popxbold{PlusJakartaSans-ExtraBold.ttf}[Path=../../../fonts/]
\newfontfamily\popxboldls{PlusJakartaSans-ExtraBold.ttf}[Path=../../../fonts/,LetterSpace=3.0]

\setlist[enumerate]{leftmargin=1.7em,itemsep=5pt,topsep=4pt}
\setlist[itemize]{leftmargin=1.7em,itemsep=4pt,topsep=4pt,label={\color{poporange}\textbullet}}
\setlength{\parindent}{0pt}
\setlength{\parskip}{5pt}
\renewcommand{\arraystretch}{1.25}

% pgfplots : style Pop par défaut
\pgfplotsset{
  every axis/.append style={axis line style={popink,line width=1pt},tick style={popink},
    grid style={popline,line width=0.5pt},label style={font=\small},tick label style={font=\footnotesize}},
  popcurve/.style={poporange,line width=1.8pt,no markers,smooth},
}
% Même style disponible dans un \draw TikZ simple (hors pgfplots)
\tikzset{popcurve/.style={poporange,line width=1.8pt}}

% ============================================================
% Pill / squiggle
% ============================================================
\newcommand{\poppill}[3]{%
  \tikz[baseline=-0.55ex]\node[draw=popink,line width=1.2pt,rounded corners=2.6pt,%
    fill=#2,inner xsep=6.5pt,inner ysep=3pt]{%
    \popxboldls\fontsize{7}{7}\selectfont\color{#3}\MakeUppercase{#1}};%
}
\newcommand{\popsquiggle}[1]{%
  \tikz[baseline=-0.4ex]{
    \draw[#1,line width=2.6pt,line cap=round]
      plot [smooth] coordinates {(0,0.09) (0.34,0.01) (0.68,0.21) (1.02,0.01) (1.36,0.10)};
  }%
}
% Mot-clé surligné (marqueur orange sous le texte)
\newcommand{\cle}[1]{{\popxbold\color{popdark}#1}}

% ============================================================
% En-tête / pied
% ============================================================
\pagestyle{fancy}
\fancyhf{}
\renewcommand{\headrulewidth}{0pt}
\renewcommand{\footrulewidth}{0pt}
\lhead{\raisebox{-0.15ex}{\poplogo\fontsize{13}{13}\selectfont\color{popink}OnBuch\color{poporange}+}}
\rhead{\poppill{\DOCMATIERE\ \textbullet\ \DOCNIVEAU\ \textbullet\ Leçon \DOCLECON}{white}{popink}}
\lfoot{\popxbold\fontsize{7.6}{7.6}\selectfont\color{popmuted}\MakeUppercase{Cours signé OnBuch+}\ \textbullet\ {\popsemi\DOCTITRE}}
\rfoot{\tikz[baseline=-0.6ex]\node[rounded corners=8.5pt,fill=popink,minimum height=17pt,minimum width=17pt,inner xsep=5pt,inner sep=0]{\popxbold\fontsize{8}{8}\selectfont\color{popcream}\thepage\,/\,\pageref*{LastPage}};}

% ============================================================
% Titres
% ============================================================
\newcommand{\chaptitre}[2]{%
  \begin{tcolorbox}[enhanced,colback=poporange,colframe=popink,boxrule=2.6pt,arc=11pt,
      drop shadow=popink,left=15pt,right=15pt,top=12pt,bottom=11pt,before skip=3pt,after skip=18pt]
    \poppill{#2}{popink}{popcream}\par\vspace{9pt}
    {\raggedright\hyphenpenalty=10000\exhyphenpenalty=10000\popdisplay\fontsize{22}{24}\selectfont\color{popcream}\MakeUppercase{#1}\par}
  \end{tcolorbox}
}
% Section numérotée : pastille orange avec le numéro + titre Archivo
\newcounter{popsec}
\newcommand{\coursec}[1]{%
  \stepcounter{popsec}\setcounter{popsub}{0}%
  \par\vspace{16pt}\noindent
  \begin{minipage}{\linewidth}
  \tikz[baseline=-0.7ex]\node[draw=popink,line width=1.6pt,fill=poporange,rounded corners=4pt,minimum size=22pt,inner sep=0]{\popdisplay\fontsize{12}{12}\selectfont\color{popcream}\thepopsec};%
  \hspace{8pt}{\popdisplay\fontsize{15}{17}\selectfont\color{popink}\MakeUppercase{#1}}\par
  \vspace{4pt}\noindent{\color{poporange}\rule{36pt}{3.6pt}}
  \end{minipage}\par\nopagebreak\vspace{8pt}
}
\newcounter{popsub}
\newcommand{\courssub}[1]{%
  \stepcounter{popsub}%
  \par\vspace{9pt}\noindent{\popxbold\fontsize{11.5}{13}\selectfont\color{popink}{\color{poporange}\thepopsec.\thepopsub}\hspace{6pt}#1}\par\nopagebreak\vspace{4pt}
}

% ============================================================
% Boîtes pédagogiques — même recette : fond blanc, bordure épaisse colorée,
% ombre dure encre, badge plein. Titre optionnel [#1].
% ============================================================
\tcbset{popbase/.style={breakable,enhanced,colback=white,boxrule=2.3pt,arc=9pt,
  shadow={3pt}{-3pt}{0pt}{popink},left=14pt,right=14pt,top=11pt,bottom=11pt,before skip=11pt,after skip=15pt,
  fontupper=\popsemi\fontsize{10.6}{14.5}\selectfont\color{popink},
  fontlower=\popsemi\fontsize{10.6}{14.5}\selectfont\color{popink}}}
% Coche dessinée (le glyphe ✓ n'existe pas dans les polices du document)
\newcommand{\popcheck}[1]{\tikz[baseline=-0.5ex]\draw[#1,line width=1.6pt,line cap=round,line join=round](-0.13,0.0)--(-0.04,-0.09)--(0.13,0.1);}
\newcommand{\popboxhead}[3]{\poppill{#1}{#2}{white}\ifx\relax#3\relax\else\hspace{6pt}{\popxbold\fontsize{10.6}{12}\selectfont\color{popink}#3}\fi\par\vspace{7pt}}

\newtcolorbox{aretenir}[1][]{popbase,colframe=popgreen,before upper={\popboxhead{À retenir}{popgreen}{#1}}}
\newtcolorbox{exemplebox}[1][]{popbase,colframe=poporange,before upper={\popboxhead{Exemple}{poporange}{#1}}}
\newtcolorbox{definition}[1][]{popbase,colframe=popblue,before upper={\popboxhead{Définition}{popblue}{#1}}}
\newtcolorbox{propriete}[1][]{popbase,colframe=poppurple,before upper={\popboxhead{Propriété}{poppurple}{#1}}}
\newtcolorbox{methode}[1][]{popbase,colframe=popgold,before upper={\popboxhead{Méthode}{popgold}{#1}}}
\newtcolorbox{attention}[1][]{popbase,colframe=poppink,before upper={\popboxhead{Attention}{poppink}{#1}}}
\newtcolorbox{experience}[1][]{popbase,colframe=popblue,colback=popblueL,before upper={\popboxhead{Expérience}{popblue}{#1}}}
% « Le savais-tu ? » : carte violette pleine (contexte, culture, Cameroun)
\newtcolorbox{savaistu}[1][]{popbase,colback=poppurple,colframe=popink,
  fontupper=\popsemi\fontsize{10.4}{14.2}\selectfont\color{white},
  before upper={\poppill{Le savais-tu\,?}{white}{poppurple}\ifx\relax#1\relax\else\hspace{6pt}{\popxbold\color{white}#1}\fi\par\vspace{7pt}}}
% Exercice résolu : énoncé en haut, \tcblower puis la solution sur fond crème
\newcounter{popexo}\newcounter{popexr}
\newtcolorbox{exoresolu}[1][]{popbase,colframe=poporange,colbacklower=popcream,
  segmentation style={popink,line width=1.2pt,dashed},
  before upper={\stepcounter{popexr}\popboxhead{Exercice résolu \thepopexr}{poporange}{#1}},
  before lower={\poppill{Solution}{popink}{popcream}\par\vspace{6pt}}}
% Exercice (non résolu en place) — la correction est dans la partie Corrigés
\newtcolorbox{exercice}[1][]{popbase,colframe=popink,
  before upper={\stepcounter{popexo}\popboxhead{Exercice \thepopexo}{popink}{#1}}}
% Corrigé
\newtcolorbox{corrige}[1][]{popbase,colframe=popgreen,colback=popgreenL,
  before upper={\popboxhead{Corrigé}{popgreen}{#1}}}
% Objectifs / prérequis / bilan
\newtcolorbox{objectifs}{popbase,colframe=popink,colback=poporangeL,
  before upper={\popboxhead{Objectifs de la leçon}{popink}{}}}
\newtcolorbox{prerequis}{popbase,colframe=popink,colback=popblueL,
  before upper={\popboxhead{Prérequis}{popblue}{}}}
\newtcolorbox{bilan}{popbase,colframe=popink,colback=white,boxrule=2.8pt,
  before upper={\popboxhead{Fiche bilan}{poporange}{L'essentiel à connaître pour l'examen}}}
% Figure encadrée avec légende
\newtcolorbox{popfigure}[1][]{enhanced,breakable=false,colback=white,colframe=popline,boxrule=1.4pt,arc=8pt,
  left=10pt,right=10pt,top=10pt,bottom=8pt,before skip=10pt,after skip=12pt,halign=center,
  lower separated=false,halign lower=center,
  fontlower=\popsemi\fontsize{9}{11.5}\selectfont\color{popmuted},
  #1}
\newcommand{\legende}[1]{\par\vspace{6pt}{\popsemi\fontsize{9}{11.5}\selectfont\color{popmuted}#1}}

% ============================================================
% Signature OnBuch+
% ============================================================
\newcommand{\poplogoinline}[1]{{\poplogo\fontsize{#1}{#1}\selectfont\color{popink}OnBuch\color{poporange}+}}
\newcommand{\signatureonbuch}{%
  \begin{tcolorbox}[enhanced,colback=popink,colframe=popink,boxrule=2.2pt,arc=10pt,
      drop shadow=poporange,left=14pt,right=14pt,top=10pt,bottom=10pt,before skip=0pt,after skip=0pt]
    \tikz[baseline=-0.7ex]\node[circle,fill=popgreen,draw=popcream,line width=1.3pt,minimum size=22pt,inner sep=0]{\popcheck{white}};%
    \hspace{9pt}\begin{minipage}[c]{0.62\linewidth}
      {\popxbold\fontsize{12}{13}\selectfont\color{popcream}Signé\ \textbullet\ {\poplogo OnBuch\color{poporange}+}}\par\vspace{2pt}
      {\raggedright\popsemi\fontsize{8}{10.5}\selectfont\color{popline}\MakeUppercase{Équipe pédagogique OnBuch+}\\\MakeUppercase{Conforme au programme officiel MINESEC}\par}
    \end{minipage}\hfill
    {\poplogo\fontsize{22}{22}\selectfont\color{popcream}OnBuch\color{poporange}+}
  \end{tcolorbox}%
}

% ============================================================
% Page de garde
% #1 titre · #2 matière · #3 niveau · #4 module · #5 n° leçon · #6 durée ·
% #7 description / objectifs (texte ou itemize)
% ============================================================
\newcommand{\pagegarde}[7]{%
  \thispagestyle{empty}
  \newgeometry{a4paper,margin=1.7cm,top=1.6cm,bottom=1.4cm}%
  \begin{tikzpicture}[remember picture,overlay]
    \fill[poporange!18] ([xshift=-3.2cm,yshift=-3.6cm]current page.north east) circle (4.6cm);
    \fill[poppurple!14] ([xshift=2.4cm,yshift=7.6cm]current page.south west) circle (3.8cm);
  \end{tikzpicture}%
  {\poplogo\fontsize{30}{30}\selectfont\color{popink}OnBuch\color{poporange}+}\hfill
  \raisebox{0.6ex}{\poppill{Cours signé \textbullet\ Premium}{popink}{popcream}}\par
  \vspace{1pt}\popsquiggle{poppurple}\par
  \vspace{8pt}
  \begin{tcolorbox}[enhanced,colback=poporange,colframe=popink,boxrule=2.8pt,arc=13pt,
      drop shadow=popink,left=18pt,right=18pt,top=12pt,bottom=13pt,after skip=10pt]
    \poppill{#2 \textbullet\ #3}{popink}{popcream}\hspace{5pt}\poppill{#4}{white}{popink}\par\vspace{8pt}
    {\popdisplay\fontsize{14}{14}\selectfont\color{popink}LEÇON #5}\par\vspace{4pt}
    {\raggedright\hyphenpenalty=10000\exhyphenpenalty=10000\popdisplay\fontsize{25}{27}\selectfont\color{popcream}\MakeUppercase{#1}\par}
    \vspace{8pt}
    {\popxbold\fontsize{9.5}{9.5}\selectfont\color{popink}\MakeUppercase{Durée conseillée : #6}}
  \end{tcolorbox}
  \begin{tcolorbox}[enhanced,colback=white,colframe=popink,boxrule=2.2pt,arc=10pt,
      drop shadow=popink,left=16pt,right=16pt,top=10pt,bottom=10pt,after skip=8pt]
    \poppill{Dans cette leçon}{poppurple}{white}\par\vspace{5pt}
    \popsemi\fontsize{9.3}{12.6}\selectfont\color{popink}#7
  \end{tcolorbox}
  \noindent\begin{minipage}[t]{0.487\linewidth}
    \begin{tcolorbox}[enhanced,colback=popgreen,colframe=popink,boxrule=2.2pt,arc=10pt,
        drop shadow=popink,left=11pt,right=11pt,top=10pt,bottom=10pt,height=3.1cm,valign=center]
      \tcbox[colback=white,colframe=popink,boxrule=1.3pt,arc=5pt,boxsep=0pt,nobeforeafter,
        left=3pt,right=3pt,top=3pt,bottom=3pt,valign=center]{\qrcode[height=1.9cm]{https://play.google.com/store/apps/details?id=com.luvvix.onbuch&hl=fr}}%
      \hspace{8pt}\begin{minipage}[c]{0.5\linewidth}
        {\popdisplay\fontsize{11}{12.5}\selectfont\color{white}\MakeUppercase{Télécharge OnBuch}}\par\vspace{4pt}
        {\raggedright\popsemi\fontsize{8.4}{11}\selectfont\color{white}Gratuit \textbullet\ Google Play\\Tuteur IA, annales, concours, résultats\par}
      \end{minipage}
    \end{tcolorbox}
  \end{minipage}\hfill
  \begin{minipage}[t]{0.487\linewidth}
    \begin{tcolorbox}[enhanced,colback=poppurple,colframe=popink,boxrule=2.2pt,arc=10pt,
        drop shadow=popink,left=11pt,right=11pt,top=10pt,bottom=10pt,height=3.1cm,valign=center]
      \tikz[baseline=-0.7ex]\node[circle,fill=white,draw=popink,line width=1.3pt,minimum size=1.6cm,inner sep=0]{\popdisplay\fontsize{24}{24}\selectfont\color{poppurple}+};%
      \hspace{8pt}\begin{minipage}[c]{0.6\linewidth}
        {\popdisplay\fontsize{11}{12.5}\selectfont\color{white}\MakeUppercase{Passe à OnBuch+}}\par\vspace{4pt}
        {\raggedright\popsemi\fontsize{8.4}{11}\selectfont\color{white}Tous les cours signés, Tuteur IA illimité, fiches PDF — onglet OnBuch+ de l'app\par}
      \end{minipage}
    \end{tcolorbox}
  \end{minipage}
  \vspace{8pt}
  \popcontenu
  \vfill
  \signatureonbuch
  \restoregeometry
  \newpage
}

% Rangée « Ce cours contient » (page de garde)
\newcommand{\poptile}[3]{% #1 couleur #2 gros texte #3 légende
  \begin{tcolorbox}[enhanced,colback=#1,colframe=popink,boxrule=2pt,arc=9pt,shadow={2.5pt}{-2.5pt}{0pt}{popink},
      left=6pt,right=6pt,top=9pt,bottom=9pt,halign=center,valign=center,height=1.75cm,width=\linewidth,nobeforeafter]
    {\popdisplay\fontsize{12.5}{13}\selectfont\color{white}\MakeUppercase{#2}}\par\vspace{3pt}
    {\centering\popsemi\fontsize{7.8}{9.5}\selectfont\color{white}#3\par}
  \end{tcolorbox}}
\newcommand{\popcontenu}{%
  \par\noindent{\popxboldls\fontsize{7.5}{7.5}\selectfont\color{popmuted}CE COURS SIGNÉ CONTIENT}\par\vspace{7pt}
  \noindent\begin{minipage}[t]{0.235\linewidth}\poptile{popblue}{Cours}{complet, illustré, pas à pas}\end{minipage}\hfill
  \begin{minipage}[t]{0.235\linewidth}\poptile{poporange}{Méthodes}{et exercices résolus}\end{minipage}\hfill
  \begin{minipage}[t]{0.235\linewidth}\poptile{poppink}{Exercices}{type examen + corrigés}\end{minipage}\hfill
  \begin{minipage}[t]{0.235\linewidth}\poptile{popgreen}{Fiche bilan}{l'essentiel à retenir}\end{minipage}\par}

% Sommaire
\newtcolorbox{sommaire}{enhanced,colback=white,colframe=popink,boxrule=2.2pt,arc=10pt,
  drop shadow=popink,left=18pt,right=18pt,top=13pt,bottom=13pt,before skip=6pt,after skip=20pt,
  before upper={\poppill{Sommaire}{poppurple}{white}\par\vspace{9pt}\popsemi\fontsize{10.6}{14.5}\selectfont\color{popink}}}

% Signature de fin de cours — poussée tout en bas de la dernière page
\newcommand{\findecours}{%
  \par\vfill\nopagebreak
  \begin{center}\popsquiggle{poporange}\end{center}\vspace{4pt}
  \signatureonbuch
}
\AtBeginDocument{\def\llbracket{\mathopen{[\mkern-3.5mu[}}\def\rrbracket{\mathclose{]\mkern-3.5mu]}}} % intervalles d'entiers (glyphe ⟦ absent de la police)
% Glyphes absents de Fira Math : redéfinis à partir de glyphes disponibles
\AtBeginDocument{%
  \def\setminus{\mathbin{\backslash}}\let\smallsetminus\setminus
  \def\vdots{\mathord{\vbox{\baselineskip3.2pt\lineskiplimit0pt\kern4pt\hbox{.}\hbox{.}\hbox{.}}}}%
  \def\ddots{\mathinner{\mkern1mu\raise6.5pt\hbox{.}\mkern2mu\raise3.6pt\hbox{.}\mkern2mu\raise0.7pt\hbox{.}\mkern1mu}}%
  \def\triangle{\mathord{\scalebox{0.9}{$\symup{\Delta}$}}}%
  \def\nabla{\mathord{\raisebox{\depth}{\scalebox{1}[-1]{$\symup{\Delta}$}}}}%
  \def\checkmark{\popcheck{popdark}}%
  \def\star{\mathbin{\ast}}\def\bigstar{\mathord{\ast}}%
  \def\diamond{\mathbin{\scalebox{0.6}{$\Diamond$}}}%
  \def\bigcup{\mathop{\mathchoice{\raisebox{-0.3em}{\scalebox{1.7}{$\cup$}}}{\scalebox{1.25}{$\cup$}}{\cup}{\cup}}\displaylimits}%
  \def\bigcap{\mathop{\mathchoice{\raisebox{-0.3em}{\scalebox{1.7}{$\cap$}}}{\scalebox{1.25}{$\cap$}}{\cap}{\cap}}\displaylimits}%
  \def\lesseqqgtr{\mathrel{\lessgtr}}\def\bigoplus{\mathop{\oplus}\displaylimits}%
}
\providecommand{\ssi}{si et seulement si} % abréviation fréquente
\AtBeginDocument{\providecommand{\wideparen}{\overparen}} % arc de cercle

%% FIGFIX-BEGIN v5 — correctifs globaux des figures (docs/onbuch-generation/tools/figfix)
\usepackage{adjustbox}\usepackage{etoolbox}\tcbuselibrary{raster}
% toute figure TikZ/pgfplots plus large que la ligne est réduite à la largeur disponible
\BeforeBeginEnvironment{tikzpicture}{\begin{adjustbox}{max width=\linewidth}}
\AfterEndEnvironment{tikzpicture}{\end{adjustbox}}
% légendes pgfplots lisibles par-dessus les courbes
\pgfplotsset{every axis/.append style={legend style={fill=white,fill opacity=0.92,text opacity=1,draw=black!15,inner sep=3pt}}}
% étiquettes d'arêtes (syntaxe edge["..."]) avec fond blanc
\tikzset{every edge quotes/.append style={fill=white,inner sep=1.2pt}}
%% FIGFIX-END
\begin{document}

\pagegarde{Équations, inéquations polynômiales et irrationnelles}{Mathématiques}{1re Tech}{Mathématiques – classe de Première technique}{N07}{2 séances (fiche de progression harmonisée)}{Cette leçon poursuit l'étude des équations et inéquations polynômiales : systèmes se ramenant au second degré, polynômes de degré 3 (racines, factorisation par x − a, signe), puis équations et inéquations irrationnelles. Elle fournit les outils algébriques indispensables pour modéliser et résoudre des problèmes concrets de la vie courante et préparer l'épreuve du Probatoire technique.
\vspace{4pt}
\begin{multicols}{2}\small
\begin{itemize}
\item À la fin de cette leçon, je sais résoudre dans ℝ² un système de deux équations à deux inconnues se ramenant à une équation du second degré (somme-produit, substitution)
\item À la fin de cette leçon, je sais déterminer une racine évidente d'un polynôme de degré 3 et factoriser ce polynôme par (x − a)
\item À la fin de cette leçon, je sais effectuer la division euclidienne d'un polynôme de degré 3 par (x − a) ou utiliser l'identification des coefficients
\item À la fin de cette leçon, je sais dresser le tableau de signe d'un polynôme de degré 3 factorisé
\item À la fin de cette leçon, je sais résoudre une équation irrationnelle en posant les conditions d'existence et en vérifiant les solutions obtenues
\item À la fin de cette leçon, je sais résoudre une inéquation irrationnelle des types √A ≤ B et √A ≥ B en discutant selon le signe de B
\end{itemize}\end{multicols}}

\begin{sommaire}
\begin{multicols}{2}
\begin{enumerate}
\item Rappels : équations et inéquations du second degré
\item Systèmes se ramenant à une équation du second degré
\item Polynômes de degré 3 : racines et factorisation
\item Signe d'un polynôme de degré 3 et inéquations
\item Équations irrationnelles
\item Inéquations irrationnelles et applications
\item Activité d'intégration
\item Méthodes et exercices résolus
\item Exercices
\item Corrigés des exercices
\item Fiche bilan
\end{enumerate}
\end{multicols}
\end{sommaire}

\begin{objectifs}
\begin{itemize}
\item À la fin de cette leçon, je sais résoudre dans ℝ² un système de deux équations à deux inconnues se ramenant à une équation du second degré (somme-produit, substitution)
\item À la fin de cette leçon, je sais déterminer une racine évidente d'un polynôme de degré 3 et factoriser ce polynôme par (x − a)
\item À la fin de cette leçon, je sais effectuer la division euclidienne d'un polynôme de degré 3 par (x − a) ou utiliser l'identification des coefficients
\item À la fin de cette leçon, je sais dresser le tableau de signe d'un polynôme de degré 3 factorisé
\item À la fin de cette leçon, je sais résoudre une équation irrationnelle en posant les conditions d'existence et en vérifiant les solutions obtenues
\item À la fin de cette leçon, je sais résoudre une inéquation irrationnelle des types √A ≤ B et √A ≥ B en discutant selon le signe de B
\item À la fin de cette leçon, je sais modéliser une situation concrète de la vie courante par une équation ou un système et interpréter les solutions
\item À la fin de cette leçon, je sais rédiger et justifier une démarche complète de résolution (conditions, équivalence, vérification, conclusion)
\end{itemize}
\end{objectifs}

\begin{prerequis}
\begin{itemize}
\item Résolution des équations et inéquations du second degré dans ℝ (discriminant, factorisation, tableau de signes)
\item Règles de calcul sur les racines carrées et condition d'existence de √A (A ≥ 0)
\item Développement, factorisation et identités remarquables (classe de Seconde)
\item Résolution des systèmes linéaires de deux équations à deux inconnues (substitution, combinaison)
\item Notion de polynôme : degré, coefficients, égalité de deux polynômes
\end{itemize}
\end{prerequis}

\newpage

\coursec{Rappels : équations et inéquations du second degré}

\courssub{Situation d'accroche : la plaque de contreplaqué de Monsieur Nkoulou}

Monsieur Nkoulou est menuisier à Yaoundé. Un client lui commande une table dont le plateau sera découpé dans une plaque de contreplaqué rectangulaire. Le client est exigeant : il veut un plateau d'\textbf{aire exactement égale à \SI{2}{m^2}} et de \textbf{périmètre exactement égal à \SI{6}{m}}. Monsieur Nkoulou se demande : quelles longueur et quelle largeur dois-je choisir ?

Notons $x$ la longueur et $y$ la largeur (en mètres). Les conditions du client s'écrivent :
\[
\begin{cases} xy = 2 \\ 2x + 2y = 6 \end{cases}
\qquad\text{c'est-à-dire}\qquad
\begin{cases} xy = 2 \\ x + y = 3 \end{cases}
\]
En remplaçant $y = 3 - x$ dans la première équation, on obtient $x(3 - x) = 2$, soit :
\[ x^2 - 3x + 2 = 0. \]
Voilà une \cle{équation du second degré} ! Plus tard, pour dimensionner un canal d'irrigation, Monsieur Nkoulou devra même résoudre une équation contenant une \emph{racine carrée}, où certaines « solutions » trouvées au brouillon devront être rejetées : ce sont les fameuses \cle{solutions parasites}.

\textbf{Problématique :} comment résoudre rigoureusement une équation ou une inéquation du second degré, puis une équation ou une inéquation contenant des racines carrées, sans tomber dans les pièges ? C'est tout l'objet de cette leçon. Nous commençons par consolider nos outils : le trinôme du second degré.

\courssub{Forme canonique du trinôme}

\begin{definition}[Trinôme du second degré]
On appelle \cle{trinôme du second degré} toute expression de la forme
\[ f(x) = ax^2 + bx + c \quad \text{avec } a \neq 0, \]
où $a$, $b$ et $c$ sont des réels fixés et $x$ un réel variable. Sa représentation graphique est une \cle{parabole}.
\end{definition}

L'idée de la \cle{forme canonique} est de réécrire le trinôme en faisant apparaître un carré parfait, comme on le fait au collège avec les identités remarquables. Partons de $ax^2 + bx + c$ et factorisons par $a$ :
\[
ax^2 + bx + c = a\left(x^2 + \frac{b}{a}x\right) + c.
\]
Or $x^2 + \dfrac{b}{a}x$ est le début du carré $\left(x + \dfrac{b}{2a}\right)^2 = x^2 + \dfrac{b}{a}x + \dfrac{b^2}{4a^2}$. En ajoutant et en retirant $\dfrac{b^2}{4a^2}$, on obtient :

\begin{propriete}[Forme canonique]
Tout trinôme $ax^2 + bx + c$ ($a \neq 0$) s'écrit de manière unique sous la \cle{forme canonique} :
\[
ax^2 + bx + c = a\left(x + \frac{b}{2a}\right)^2 - \frac{b^2 - 4ac}{4a}
= a\left(x - \alpha\right)^2 + \beta,
\]
avec $\alpha = -\dfrac{b}{2a}$ et $\beta = -\dfrac{b^2 - 4ac}{4a} = f(\alpha)$. Le point $S(\alpha\,;\,\beta)$ est le \cle{sommet} de la parabole, et la droite d'équation $x = \alpha$ est son axe de symétrie.
\end{propriete}

\begin{exemplebox}[Mise sous forme canonique]
Écrivons $f(x) = 2x^2 - 5x + 3$ sous forme canonique. Ici $a = 2$, $b = -5$, $c = 3$ :
\[
\alpha = -\frac{b}{2a} = \frac{5}{4}, \qquad
\beta = f\left(\frac{5}{4}\right) = 2\times\frac{25}{16} - 5\times\frac{5}{4} + 3 = \frac{25}{8} - \frac{25}{4} + 3 = -\frac{1}{8}.
\]
Donc $f(x) = 2\left(x - \dfrac{5}{4}\right)^2 - \dfrac{1}{8}$. Le sommet de la parabole est $S\left(\dfrac{5}{4}\,;\,-\dfrac{1}{8}\right)$ et, comme $a = 2 > 0$, la parabole est tournée vers le haut.
\end{exemplebox}

\courssub{Discriminant et résolution de l'équation $ax^2 + bx + c = 0$}

Pour résoudre $ax^2 + bx + c = 0$, on utilise la forme canonique : l'équation équivaut à
\[
a\left(x + \frac{b}{2a}\right)^2 = \frac{b^2 - 4ac}{4a}
\quad\Longleftrightarrow\quad
\left(x + \frac{b}{2a}\right)^2 = \frac{b^2 - 4ac}{4a^2}.
\]
Tout dépend donc du signe du nombre $b^2 - 4ac$, qui porte un nom :

\begin{definition}[Discriminant]
Le \cle{discriminant} du trinôme $ax^2 + bx + c$ est le réel
\[ \Delta = b^2 - 4ac. \]
\end{definition}

\begin{propriete}[Nombre de solutions selon le signe de $\Delta$ — admise]
Soit l'équation $ax^2 + bx + c = 0$ avec $a \neq 0$ et $\Delta = b^2 - 4ac$.
\begin{itemize}
\item Si $\Delta > 0$ : l'équation admet \textbf{deux solutions réelles distinctes} :
\[ x_1 = \frac{-b - \sqrt{\Delta}}{2a} \qquad \text{et} \qquad x_2 = \frac{-b + \sqrt{\Delta}}{2a}. \]
\item Si $\Delta = 0$ : l'équation admet \textbf{une unique solution} (dite \emph{racine double}) : $x_0 = -\dfrac{b}{2a}$.
\item Si $\Delta < 0$ : l'équation \textbf{n'admet aucune solution réelle}.
\end{itemize}
\end{propriete}

Graphiquement, ces trois cas correspondent aux trois positions possibles de la parabole par rapport à l'axe des abscisses : elle le coupe en deux points, le touche en un seul point, ou ne le rencontre pas.

\begin{popfigure}
\begin{tikzpicture}
\begin{axis}[width=\linewidth, height=5.6cm, axis lines=middle, xlabel={$x$}, ylabel={$y$},
xmin=-3.2, xmax=3.2, ymin=-2.2, ymax=3.2, xtick=\empty, ytick=\empty,
legend style={at={(0.02,0.98)}, anchor=north west, font=\small, draw=popline}]
\addplot[popcurve,domain=-2.6:2.6,samples=200]{x^2 - 1};
\addlegendentry{$\Delta > 0$ : 2 racines}
\addplot[poporange,very thick,domain=-2.6:2.6,samples=200]{x^2};
\addlegendentry{$\Delta = 0$ : racine double}
\addplot[popgreen,very thick,domain=-2.6:2.6,samples=200]{x^2 + 1};
\addlegendentry{$\Delta < 0$ : aucune racine}
\end{axis}
\end{tikzpicture}
\legende{Les trois positions d'une parabole par rapport à l'axe des abscisses selon le signe de $\Delta$.}
\end{popfigure}

\begin{exoresolu}[Résolution de $2x^2 - 5x + 3 = 0$]
Résoudre dans $\mathbb{R}$ l'équation $2x^2 - 5x + 3 = 0$.
\tcblower
\textbf{Étape 1 : calcul du discriminant.} Avec $a = 2$, $b = -5$, $c = 3$ :
\[ \Delta = b^2 - 4ac = (-5)^2 - 4 \times 2 \times 3 = 25 - 24 = 1. \]
\textbf{Étape 2 : discussion.} $\Delta = 1 > 0$ : il y a deux solutions distinctes.
\textbf{Étape 3 : calcul des racines.}
\[ x_1 = \frac{-b - \sqrt{\Delta}}{2a} = \frac{5 - 1}{4} = 1, \qquad
x_2 = \frac{-b + \sqrt{\Delta}}{2a} = \frac{5 + 1}{4} = \frac{3}{2}. \]
\textbf{Conclusion :} $S = \left\{1\,;\,\dfrac{3}{2}\right\}$. Vérification : $2(1)^2 - 5(1) + 3 = 0$ et $2\left(\tfrac32\right)^2 - 5\left(\tfrac32\right) + 3 = \tfrac92 - \tfrac{15}{2} + 3 = 0$. 
\end{exoresolu}

\begin{attention}[Bien identifier $a$, $b$, $c$ avec leurs signes]
Avant de calculer $\Delta$, écris toujours l'équation sous la forme $ax^2 + bx + c = 0$ et identifie les coefficients \emph{avec leur signe}. Pour $2x^2 - 5x + 3 = 0$, on a $b = -5$ (et non $5$) : oublier ce signe fausse le calcul de $-b$ dans les formules des racines. Autre piège : si l'équation est donnée sous la forme $2x^2 - 5x = -3$, il faut d'abord tout ramener dans le membre de gauche.
\end{attention}

\courssub{Factorisation du trinôme}

\begin{propriete}[Forme factorisée]
Soit $ax^2 + bx + c$ un trinôme ($a \neq 0$) de discriminant $\Delta$.
\begin{itemize}
\item Si $\Delta > 0$, de racines $x_1$ et $x_2$ : \quad $ax^2 + bx + c = a(x - x_1)(x - x_2)$.
\item Si $\Delta = 0$, de racine double $x_0$ : \quad $ax^2 + bx + c = a(x - x_0)^2$.
\item Si $\Delta < 0$ : le trinôme \textbf{ne se factorise pas} dans $\mathbb{R}$.
\end{itemize}
\end{propriete}

\begin{exemplebox}[Factorisation]
Comme $2x^2 - 5x + 3$ a pour racines $1$ et $\dfrac{3}{2}$, on a :
\[ 2x^2 - 5x + 3 = 2(x - 1)\left(x - \frac{3}{2}\right) = (x - 1)(2x - 3). \]
Cette factorisation sera précieuse pour étudier le signe du trinôme.
\end{exemplebox}

\courssub{Somme et produit des racines}

Quand on connaît une racine « évidente » d'un trinôme, on peut trouver l'autre très rapidement, sans calculer $\Delta$, grâce à deux relations remarquables.

\begin{propriete}[Relations somme-produit]
Si le trinôme $ax^2 + bx + c$ ($a \neq 0$) admet deux racines $x_1$ et $x_2$ (éventuellement égales), alors :
\[ S = x_1 + x_2 = -\frac{b}{a} \qquad \text{et} \qquad P = x_1 \times x_2 = \frac{c}{a}. \]
\end{propriete}

\begin{experience}[Démonstration guidée des relations somme-produit]
Supposons $\Delta \geq 0$, de sorte que le trinôme se factorise : $ax^2 + bx + c = a(x - x_1)(x - x_2)$.
\begin{enumerate}
\item Développons le membre de droite : $(x - x_1)(x - x_2) = x^2 - (x_1 + x_2)x + x_1x_2$, donc
\[ ax^2 + bx + c = a\,x^2 - a(x_1 + x_2)\,x + a\,x_1x_2. \]
\item Deux polynômes égaux pour tout $x$ ont les mêmes coefficients. Identifions :
\begin{itemize}
\item coefficient de $x$ : \quad $b = -a(x_1 + x_2)$, d'où $x_1 + x_2 = -\dfrac{b}{a}$ ;
\item coefficient constant : \quad $c = a\,x_1x_2$, d'où $x_1x_2 = \dfrac{c}{a}$.
\end{itemize}
\end{enumerate}
Les deux relations sont démontrées.
\end{experience}

\begin{exemplebox}[Application : trouver une racine sans calculer $\Delta$]
Pour l'équation $2x^2 - 5x + 3 = 0$, on remarque que $x_1 = 1$ est racine évidente (car $2 - 5 + 3 = 0$). Alors le produit donne immédiatement : $1 \times x_2 = \dfrac{c}{a} = \dfrac{3}{2}$, donc $x_2 = \dfrac{3}{2}$. Contrôle avec la somme : $1 + \dfrac{3}{2} = \dfrac{5}{2} = -\dfrac{b}{a}$. C'est cohérent.
\end{exemplebox}

\begin{savaistu}[Retour à la plaque de Monsieur Nkoulou]
Chercher deux nombres dont on connaît la somme $S$ et le produit $P$, c'est exactement résoudre l'équation $x^2 - Sx + P = 0$. Pour la plaque de contreplaqué, $S = 3$ et $P = 2$, donc les dimensions sont les solutions de $x^2 - 3x + 2 = 0$, à savoir $1$ et $2$. La plaque mesure donc \SI{2}{m} de long sur \SI{1}{m} de large : Monsieur Nkoulou peut passer commande à son fournisseur de Douala ! Nous étudierons ces systèmes en détail dans la section suivante.
\end{savaistu}

\courssub{Signe du trinôme et inéquations du second degré}

Résoudre une inéquation comme $2x^2 - 5x + 3 \leq 0$ revient à déterminer \emph{le signe} du trinôme selon les valeurs de $x$. La factorisation $a(x - x_1)(x - x_2)$ rend cette étude possible par un tableau de signes.

\begin{propriete}[Signe du trinôme]
Soit $ax^2 + bx + c$ un trinôme ($a \neq 0$) de discriminant $\Delta$.
\begin{itemize}
\item Si $\Delta > 0$ (racines $x_1 < x_2$) : le trinôme est \textbf{du signe de $a$ à l'extérieur des racines} et \textbf{du signe de $-a$ entre les racines}.
\item Si $\Delta = 0$ (racine double $x_0$) : le trinôme est \textbf{du signe de $a$} pour tout $x \neq x_0$, et s'annule en $x_0$.
\item Si $\Delta < 0$ : le trinôme est \textbf{du signe de $a$ pour tout réel $x$} (il ne s'annule jamais).
\end{itemize}
\end{propriete}

Voici le tableau de signes type dans le cas $\Delta > 0$ :

\begin{center}
\begin{tabular}{|c|ccccc|}\hline
$x$ & $-\infty$ & & $x_1$ & & $+\infty$ \\ \hline
$ax^2+bx+c$ & & signe de $a$ & $0$ & signe de $-a$ & \\ \hline
\end{tabular}
\end{center}

\begin{attention}[Le signe dépend de $a$, pas seulement des racines !]
L'erreur la plus fréquente au Probatoire est d'écrire mécaniquement « $+$ à l'extérieur, $-$ entre les racines ». Cela n'est vrai que si $a > 0$ ! Si $a < 0$, c'est l'inverse : le trinôme est \emph{négatif} à l'extérieur des racines et \emph{positif} entre elles. Avant tout tableau de signes, commence par regarder le signe de $a$. Exemple : $-x^2 + 3x - 2$ (avec $a = -1 < 0$, racines $1$ et $2$) est \textbf{positif} entre $1$ et $2$.
\end{attention}

\begin{methode}[Résoudre une inéquation du second degré en 4 étapes]
Pour résoudre une inéquation du type $ax^2 + bx + c \leq 0$ (ou $\geq$, $<$, $>$) :
\begin{enumerate}
\item \textbf{Racines} : ramener tout dans un membre, calculer $\Delta = b^2 - 4ac$ et déterminer les racines éventuelles du trinôme.
\item \textbf{Tableau} : construire le tableau de signes du trinôme en appliquant la règle « signe de $a$ à l'extérieur des racines, signe de $-a$ entre les racines » (attention au signe de $a$ !).
\item \textbf{Lecture} : repérer dans le tableau les intervalles où le trinôme a le signe demandé par l'inéquation.
\item \textbf{Conclusion en intervalles} : écrire l'ensemble des solutions sous forme d'intervalle ou de réunion d'intervalles, en \emph{fermant} les crochets si l'inégalité est large ($\leq$ ou $\geq$) et en les \emph{ouvrant} si elle est stricte ($<$ ou $>$).
\end{enumerate}
\end{methode}

\begin{exoresolu}[Résolution de $2x^2 - 5x + 3 \leq 0$]
Résoudre dans $\mathbb{R}$ l'inéquation $2x^2 - 5x + 3 \leq 0$.
\tcblower
\textbf{Étape 1 : racines.} Nous avons déjà calculé $\Delta = 1 > 0$ et les racines $x_1 = 1$ et $x_2 = \dfrac{3}{2}$.

\textbf{Étape 2 : tableau de signes.} Ici $a = 2 > 0$ : le trinôme est positif à l'extérieur des racines, négatif entre elles.
\begin{center}
\begin{tabular}{|c|ccccc|}\hline
$x$ & $-\infty$ & & $1$ & & $+\infty$ \\ \hline
$2x^2 - 5x + 3$ & & $+$ & $0$ & $-$ & \\ \hline
\end{tabular}
\end{center}
Plus précisément : signe $+$ sur $\left]-\infty\,;\,1\right[$, nul en $1$, signe $-$ sur $\left]1\,;\,\dfrac{3}{2}\right[$, nul en $\dfrac{3}{2}$, signe $+$ sur $\left]\dfrac{3}{2}\,;\,+\infty\right[$.

\textbf{Étape 3 : lecture.} On cherche où le trinôme est négatif ou nul : c'est entre les racines, racines comprises.

\textbf{Étape 4 : conclusion.} L'inégalité est large ($\leq$), donc on ferme les crochets :
\[ S = \left[1\,;\,\frac{3}{2}\right]. \]
\end{exoresolu}

\begin{exercice}[À toi de jouer]
\begin{enumerate}
\item Résoudre dans $\mathbb{R}$ l'équation $3x^2 + 2x - 1 = 0$ (on pourra chercher une racine évidente).
\item En déduire la factorisation de $3x^2 + 2x - 1$.
\item Résoudre dans $\mathbb{R}$ l'inéquation $3x^2 + 2x - 1 > 0$.
\end{enumerate}
\end{exercice}

\begin{corrige}[Exercice 1]
\begin{enumerate}
\item On remarque que $x_1 = -1$ est racine : $3(-1)^2 + 2(-1) - 1 = 3 - 2 - 1 = 0$. Le produit des racines vaut $\dfrac{c}{a} = -\dfrac{1}{3}$, donc $x_2 = \dfrac{1}{3}$. Contrôle : $\Delta = 4 + 12 = 16$, $\sqrt{\Delta} = 4$, $x_1 = \dfrac{-2-4}{6} = -1$, $x_2 = \dfrac{-2+4}{6} = \dfrac{1}{3}$. Donc $S = \left\{-1\,;\,\dfrac{1}{3}\right\}$.
\item $3x^2 + 2x - 1 = 3(x + 1)\left(x - \dfrac{1}{3}\right) = (x + 1)(3x - 1)$.
\item Ici $a = 3 > 0$ : le trinôme est strictement positif à l'extérieur des racines. L'inégalité est stricte, donc les crochets sont ouverts : $S = \left]-\infty\,;\,-1\right[ \cup \left]\dfrac{1}{3}\,;\,+\infty\right[$.
\end{enumerate}
\end{corrige}

\begin{aretenir}[L'essentiel de la section]
\begin{itemize}
\item Tout trinôme $ax^2 + bx + c$ ($a \neq 0$) s'écrit sous \cle{forme canonique} $a\left(x + \frac{b}{2a}\right)^2 - \frac{\Delta}{4a}$.
\item Le \cle{discriminant} $\Delta = b^2 - 4ac$ décide du nombre de solutions : deux si $\Delta > 0$, une si $\Delta = 0$, aucune si $\Delta < 0$.
\item Si $\Delta > 0$ : $x_1 = \dfrac{-b - \sqrt{\Delta}}{2a}$, $x_2 = \dfrac{-b + \sqrt{\Delta}}{2a}$, et $ax^2 + bx + c = a(x - x_1)(x - x_2)$.
\item Relations : $x_1 + x_2 = -\dfrac{b}{a}$ et $x_1x_2 = \dfrac{c}{a}$.
\item Signe du trinôme : \textbf{signe de $a$ à l'extérieur des racines, signe de $-a$ entre les racines} — toujours vérifier le signe de $a$ avant de conclure.
\end{itemize}
\end{aretenir}

\coursec{Systèmes se ramenant à une équation du second degré}

\courssub{Transition et mise en situation}

Dans la section précédente, nous avons rappelé les méthodes de résolution des équations et inéquations du second degré à une inconnue. Nous allons maintenant voir comment ces outils permettent de résoudre des \textbf{systèmes de deux équations à deux inconnues} qui, après une manipulation algébrique habile, se ramènent exactement à une équation du second degré.

\medskip

\begin{experience}[Problème d'accroche : dimensions d'un rectangle]
Un menuisier doit fabriquer une table rectangulaire. Le client impose un \textbf{périmètre de 6~m} et une \textbf{surface de 2~m²}. Quelles doivent être la longueur $L$ et la largeur $\ell$ de la table ?

\medskip
\noindent \textbf{Modélisation :}
\begin{itemize}
    \item Périmètre : $2(L + \ell) = 6 \quad\Leftrightarrow\quad L + \ell = 3$.
    \item Aire : $L \times \ell = 2$.
\end{itemize}
On cherche donc deux nombres réels $L$ et $\ell$ dont la \textbf{somme} $S = 3$ et le \textbf{produit} $P = 2$.
Ce système s'écrit :
\[
\begin{cases}
x + y = 3 \\
xy = 2
\end{cases}
\qquad\text{avec } x = L,\; y = \ell.
\]

\begin{popfigure}
\begin{tikzpicture}[scale=1.2, >=Stealth]
% Axes
\draw[->] (-0.5,0) -- (3.5,0) node[right] {$x$};
\draw[->] (0,-0.5) -- (0,3.5) node[above] {$y$};
% Grid
\draw[popmuted, very thin] (0,0) grid (3,3);
% Hyperbole xy = 2
\draw[popblue, thick, domain=0.6:3, samples=100, variable=\x] plot ({\x}, {2/\x}) node[right, pos=0.9] {$xy=2$};
% Droite x + y = 3
\draw[poporange, thick] (0,3) -- (3,0) node[below right, pos=0.9] {$x+y=3$};
% Points d'intersection
\fill[popred] (1,2) circle (2.5pt) node[above left] {$(1;2)$};
\fill[popred] (2,1) circle (2.5pt) node[below right] {$(2;1)$};
% Rectangle illustration (inset)
\begin{scope}[shift={(4,1.5)}, scale=0.6]
\draw[popdark, thick] (0,0) rectangle (2,1);
\node at (1,0.5) {$2\,\text{m}^2$};
\draw[<->] (0,-0.3) -- (2,-0.3) node[midway, below] {$L$};
\draw[<->] (-0.3,0) -- (-0.3,1) node[midway, left] {$\ell$};
\end{scope}
\end{tikzpicture}
\legende{Intersection de la droite $x+y=3$ (somme) et de l'hyperbole $xy=2$ (produit). Les deux points d'intersection correspondent aux deux couples solutions $(1;2)$ et $(2;1)$.}
\end{popfigure}

\noindent On voit graphiquement deux points d'intersection : $(1;2)$ et $(2;1)$. Algébriquement, comment les trouver systématiquement ? C'est l'objet de cette section.
\end{experience}

\courssub{Systèmes symétriques de type « somme-produit »}

Un système est dit \textbf{symétrique} si les équations ne changent pas quand on échange $x$ et $y$. La forme canonique est :
\[
\begin{cases}
x + y = S \\
xy = P
\end{cases}
\]
où $S$ (somme) et $P$ (produit) sont des nombres réels connus.

\begin{propriete}[Théorème fondamental : lien avec le trinôme du second degré]
Soient $x$ et $y$ deux nombres réels, $S = x+y$ et $P = xy$.
Alors $x$ et $y$ sont les \textbf{deux racines} (réelles ou complexes) du trinôme :
\[
X^2 - S X + P = 0.
\]
Réciproquement, si $x_1$ et $x_2$ sont les deux racines de $X^2 - S X + P = 0$, alors le couple $(x_1; x_2)$ (et $(x_2; x_1)$) est solution du système $\begin{cases} x+y=S \\ xy=P \end{cases}$.

\medskip
\noindent \textbf{Démonstration (exigible au Probatoire) :}
Si $x$ et $y$ sont solutions du système, on a $y = S - x$. En remplaçant dans $xy = P$ :
\[
x(S - x) = P \;\Leftrightarrow\; Sx - x^2 = P \;\Leftrightarrow\; x^2 - Sx + P = 0.
\]
Donc $x$ est racine de $X^2 - SX + P = 0$. Par symétrie, $y$ l'est aussi.
Réciproquement, si $x_1, x_2$ sont les racines, les relations de Viète donnent $x_1+x_2=S$ et $x_1x_2=P$.
\end{propriete}

\begin{aretenir}[Condition d'existence de solutions réelles]
Le système $\begin{cases} x+y=S \\ xy=P \end{cases}$ admet des solutions \textbf{réelles} si et seulement si le discriminant du trinôme associé est positif ou nul :
\[
\Delta = S^2 - 4P \ge 0.
\]
\begin{itemize}
    \item Si $\Delta > 0$ : deux couples solutions distincts $(x_1; x_2)$ et $(x_2; x_1)$.
    \item Si $\Delta = 0$ : un unique couple solution $(x_0; x_0)$ (racine double).
    \item Si $\Delta < 0$ : aucune solution réelle (les solutions sont complexes conjuguées).
\end{itemize}
\end{aretenir}

\courssub{Méthode de résolution : 3 étapes}

\begin{methode}[Résolution d'un système symétrique $x+y=S,\; xy=P$]
\begin{enumerate}
    \item \textbf{Identifier $S$ et $P$} (somme et produit) et écrire le trinôme $X^2 - SX + P = 0$.
    \item \textbf{Calculer le discriminant} $\Delta = S^2 - 4P$ et résoudre l'équation du second degré pour trouver les racines $x_1, x_2$.
    \item \textbf{Former les couples solutions} : $(x_1; x_2)$ et $(x_2; x_1)$. Vérifier éventuellement dans le système initial.
\end{enumerate}
\end{methode}

\begin{exoresolu}[Exemple type : $\{x+y=6;\; xy=8\}$]
Résoudre dans $\mathbb{R}^2$ le système :
\[
\begin{cases}
x + y = 6 \\
xy = 8
\end{cases}
\]

\medskip
\noindent \textbf{Solution.}
\begin{enumerate}
    \item $S = 6$, $P = 8$. Trinôme : $X^2 - 6X + 8 = 0$.
    \item $\Delta = 6^2 - 4\times 8 = 36 - 32 = 4 > 0$.
    Racines : $X_1 = \frac{6 - 2}{2} = 2$ et $X_2 = \frac{6 + 2}{2} = 4$.
    \item Les couples solutions sont $(2; 4)$ et $(4; 2)$.
\end{enumerate}
\noindent \textbf{Vérification :} $2+4=6,\; 2\times4=8$ et $4+2=6,\; 4\times2=8$. \checkmark
\end{exoresolu}

\begin{attention}[Piège classique : oublier le couple symétrique]
\emph{Erreur fréquente :} « Les solutions sont $x=2$ et $x=4$, donc le couple est $(2;4)$. »
\emph{Correction :} Un système symétrique admet \textbf{deux couples solutions} (sauf racine double). Il faut \textbf{toujours} écrire les deux couples : $(x_1; x_2)$ \textbf{et} $(x_2; x_1)$. Au Probatoire, l'oubli du second couple fait perdre des points.
\end{attention}

\courssub{Systèmes de type « différence-produit »}

Certains systèmes ne donnent pas directement la somme, mais la différence :
\[
\begin{cases}
x - y = d \\
xy = P
\end{cases}
\qquad (d \in \mathbb{R}).
\]

\begin{methode}[Transformation en système somme-produit]
On utilise l'identité remarquable : $(x+y)^2 = (x-y)^2 + 4xy$.
\begin{enumerate}
    \item Calculer $(x+y)^2 = d^2 + 4P$.
    \item En déduire les valeurs possibles de la somme $S = x+y = \pm\sqrt{d^2 + 4P}$ (si $d^2+4P \ge 0$).
    \item On retombe sur un (ou deux) système(s) symétrique(s) $\begin{cases} x+y=S \\ xy=P \end{cases}$ qu'on résout par la méthode précédente.
\end{enumerate}
\end{methode}

\begin{exoresolu}[Système différence-produit]
Résoudre : $\begin{cases} x - y = 2 \\ xy = 3 \end{cases}$.

\medskip
\noindent \textbf{Solution.}
\begin{enumerate}
    \item $d=2$, $P=3$. $(x+y)^2 = 2^2 + 4\times 3 = 4 + 12 = 16$.
    \item $x+y = \pm 4$. Deux cas :
    \begin{itemize}
        \item \textbf{Cas 1 :} $x+y=4$. Système $\begin{cases} x+y=4 \\ xy=3 \end{cases}$.
              $\Delta = 16 - 12 = 4$. Racines : $1$ et $3$. Couples : $(1;3)$ et $(3;1)$.
              Vérif différence : $3-1=2$ (OK), $1-3=-2$ (non). \textbf{Seul $(3;1)$ convient.}
        \item \textbf{Cas 2 :} $x+y=-4$. Système $\begin{cases} x+y=-4 \\ xy=3 \end{cases}$.
              $\Delta = 16 - 12 = 4$. Racines : $-1$ et $-3$. Couples : $(-1;-3)$ et $(-3;-1)$.
              Vérif différence : $-1 - (-3) = 2$ (OK). \textbf{Seul $(-1;-3)$ convient.}
    \end{itemize}
\end{enumerate}
\noindent \textbf{Solutions finales :} $(3; 1)$ et $(-1; -3)$.
\end{exoresolu}

\courssub{Méthode de substitution (cas général)}

La méthode « somme-produit » ne s'applique que si le système est symétrique (ou transformable en symétrique). La méthode \textbf{générale}, toujours applicable, est la \textbf{substitution}.

\begin{methode}[Résolution par substitution]
\begin{enumerate}
    \item Dans l'équation la plus simple (souvent linéaire), \textbf{exprimer une inconnue en fonction de l'autre} (ex: $y = \dots$ ou $x = \dots$).
    \item \textbf{Remplacer} cette expression dans la seconde équation. On obtient une équation à une inconnue.
    \item \textbf{Résoudre} cette équation (souvent du second degré).
    \item \textbf{Remonter} : pour chaque valeur trouvée, calculer la valeur de l'autre inconnue.
    \item \textbf{Rédiger} les couples solutions $(x; y)$.
\end{enumerate}
\end{methode}

\begin{exemplebox}[Substitution sur un système non symétrique]
Résoudre : $\begin{cases} 2x + y = 5 \\ x^2 + y^2 = 13 \end{cases}$.

\medskip
\noindent \textbf{Solution.}
\begin{enumerate}
    \item De la première : $y = 5 - 2x$.
    \item Substitution dans la seconde : $x^2 + (5-2x)^2 = 13$.
    \item Développement : $x^2 + 25 - 20x + 4x^2 = 13 \;\Leftrightarrow\; 5x^2 - 20x + 12 = 0$.
    \item $\Delta = 400 - 240 = 160 = 16\times 10$. $\sqrt{\Delta} = 4\sqrt{10}$.
    $x_1 = \frac{20 - 4\sqrt{10}}{10} = 2 - \frac{2\sqrt{10}}{5}$, $x_2 = 2 + \frac{2\sqrt{10}}{5}$.
    \item $y_1 = 5 - 2x_1 = 1 + \frac{4\sqrt{10}}{5}$, $y_2 = 1 - \frac{4\sqrt{10}}{5}$.
\end{enumerate}
\noindent Couples : $\left(2 - \frac{2\sqrt{10}}{5};\; 1 + \frac{4\sqrt{10}}{5}\right)$ et $\left(2 + \frac{2\sqrt{10}}{5};\; 1 - \frac{4\sqrt{10}}{5}\right)$.
\end{exemplebox}

\courssub{Interprétation géométrique : droite et hyperbole}

L'équation $xy = P$ (avec $P \neq 0$) représente une \textbf{hyperbole équilatère} de centre l'origine, asymptotes les axes coordonnés.
L'équation $x + y = S$ représente une \textbf{droite} de pente $-1$.

\begin{popfigure}
\begin{tikzpicture}[scale=0.9]
\begin{axis}[
    width=\linewidth, height=0.6\linewidth,
    axis lines=middle,
    xlabel={$x$}, ylabel={$y$},
    xmin=-5, xmax=5, ymin=-5, ymax=5,
    xtick={-4,-2,2,4}, ytick={-4,-2,2,4},
    grid=major, grid style={popmuted},
    clip=false,
    legend style={at={(0.98,0.98)}, anchor=north east, fill=white, draw=popline}
]
% Hyperbole xy = 8 (branche 1er quadrant)
\addplot[popblue, thick, domain=0.5:5, samples=200] {8/x};
\addplot[popblue, thick, domain=-5:-0.5, samples=200] {8/x} node[left, pos=0.1] {$xy=8$};
% Droite x + y = 6
\addplot[poporange, thick, domain=-1:7] {6-x} node[above right, pos=0.9] {$x+y=6$};
% Points d'intersection
\addplot[only marks, mark=*, mark size=2.5, popred] coordinates {(2,4) (4,2)};
\legend{$xy=8$, $x+y=6$, Solutions}
\end{axis}
\end{tikzpicture}
\legende{Intersection de la droite $x+y=6$ et de l'hyperbole $xy=8$. Deux points d'intersection dans le premier quadrant : $(2;4)$ et $(4;2)$.}
\end{popfigure}

\begin{propriete}[Nombre de solutions réelles et géométrie]
Pour le système $\begin{cases} x+y=S \\ xy=P \end{cases}$ avec $P \neq 0$ :
\begin{itemize}
    \item $\Delta > 0 \;\Leftrightarrow\; S^2 > 4P$ : la droite coupe l'hyperbole en \textbf{deux points} distincts (deux couples solutions).
    \item $\Delta = 0 \;\Leftrightarrow\; S^2 = 4P$ : la droite est \textbf{tangente} à l'hyperbole (un couple solution double, point de contact).
    \item $\Delta < 0 \;\Leftrightarrow\; S^2 < 4P$ : la droite ne rencontre \textbf{pas} l'hyperbole (aucune solution réelle).
\end{itemize}
\end{propriete}

\courssub{Applications concrètes : problèmes de dimensions}

Ces systèmes modélisent naturellement les problèmes où l'on connaît la \textbf{somme} (périmètre, demi-périmètre) et le \textbf{produit} (aire) de deux grandeurs positives (longueur, largeur, base, hauteur).

\begin{exoresolu}[Application : bassin rectangulaire]
Un entrepreneur doit creuser un bassin rectangulaire. Le périmètre doit être de $20\,\text{m}$ et la surface du fond de $21\,\text{m}^2$. Déterminer les dimensions du bassin.

\medskip
\noindent \textbf{Solution.}
Soient $x$ (longueur) et $y$ (largeur) en mètres. $x>0, y>0$.
\begin{itemize}
    \item Périmètre : $2(x+y)=20 \;\Rightarrow\; x+y=10$. ($S=10$)
    \item Aire : $xy=21$. ($P=21$)
\end{itemize}
Système : $\begin{cases} x+y=10 \\ xy=21 \end{cases}$.
Trinôme : $X^2 - 10X + 21 = 0$.
$\Delta = 100 - 84 = 16$. $\sqrt{\Delta}=4$.
$X_1 = \frac{10-4}{2}=3$, $X_2 = \frac{10+4}{2}=7$.
Couples : $(3;7)$ et $(7;3)$.
\emph{Interprétation :} Le bassin fait $3\,\text{m}$ sur $7\,\text{m}$ (ou $7\,\text{m}$ sur $3\,\text{m}$). Par convention, on dit souvent \textbf{longueur $=7\,\text{m}$, largeur $=3\,\text{m}$}.
\end{exoresolu}

\begin{savaistu}[Histoire : les scribes babyloniens]
Il y a près de \textbf{4000 ans} (vers 1800 av. J.-C.), les scribes de Mésopotamie résolvaient déjà des problèmes du type « \emph{J'ai additionné la surface et le côté de mon carré : 45} » ou « \emph{La somme de la longueur et la largeur est 50, leur produit est 600} ».
Sur la tablette \textbf{BM 13901} (British Museum), on trouve la résolution algorithmique de $x+y=S,\; xy=P$ par la méthode du « \emph{moitié de la somme, carré, soustraire le produit, racine carrée} », exactement équivalente à notre formule du discriminant. Ils n'avaient pas notre notation algébrique, mais leur méthode géométrique (complétion du carré) était rigoureuse et enseignée dans les écoles de scribes !
\end{savaistu}

\courssub{Exercices d'entraînement}

\begin{exercice}[Exercice 1 : Systèmes symétriques]
Résoudre dans $\mathbb{R}^2$ les systèmes suivants :
\begin{enumerate}
    \item $\begin{cases} x+y=5 \\ xy=6 \end{cases}$
    \item $\begin{cases} x+y=-4 \\ xy=4 \end{cases}$
    \item $\begin{cases} x+y=2 \\ xy=5 \end{cases}$
\end{enumerate}
\end{exercice}

\begin{exercice}[Exercice 2 : Système différence-produit]
Résoudre : $\begin{cases} x - y = 4 \\ xy = 12 \end{cases}$.
\end{exercice}

\begin{exercice}[Exercice 3 : Substitution]
Résoudre par substitution : $\begin{cases} y = 2x - 1 \\ x^2 + y^2 = 10 \end{cases}$.
\end{exercice}

\begin{exercice}[Exercice 4 : Problème concret]
Un électricien doit découper un fil de cuivre de $40\,\text{cm}$ en deux morceaux pour former deux carrés. La somme des aires des deux carrés doit être $104\,\text{cm}^2$. Quelle doit être la longueur de chaque morceau ?
\emph{Indication :} Si $x$ et $y$ sont les longueurs des morceaux, les côtés des carrés sont $x/4$ et $y/4$.
\end{exercice}

\begin{corrige}[Exercice 1]
\begin{enumerate}
    \item $S=5, P=6$. $\Delta=25-24=1$. Racines $2,3$. Couples : $(2;3)$ et $(3;2)$.
    \item $S=-4, P=4$. $\Delta=16-16=0$. Racine double $-2$. Couple unique : $(-2;-2)$.
    \item $S=2, P=5$. $\Delta=4-20=-16<0$. \textbf{Pas de solution réelle}.
\end{enumerate}
\end{corrige}

\begin{corrige}[Exercice 2]
$d=4, P=12$. $(x+y)^2 = 16 + 48 = 64 \Rightarrow x+y = \pm 8$.
\begin{itemize}
    \item Cas $x+y=8$ : système $\{x+y=8; xy=12\}$. $\Delta=64-48=16$. Racines $2,6$. Couples $(2;6),(6;2)$. Vérif $x-y$ : $6-2=4$ (OK), $2-6=-4$ (non). $\Rightarrow$ $(6;2)$.
    \item Cas $x+y=-8$ : système $\{x+y=-8; xy=12\}$. $\Delta=16$. Racines $-2,-6$. Couples $(-2;-6),(-6;-2)$. Vérif $x-y$ : $-2-(-6)=4$ (OK). $\Rightarrow$ $(-2;-6)$.
\end{itemize}
Solutions : $(6;2)$ et $(-2;-6)$.
\end{corrige}

\begin{corrige}[Exercice 3]
Substitution : $x^2 + (2x-1)^2 = 10 \Leftrightarrow x^2 + 4x^2 -4x +1 = 10 \Leftrightarrow 5x^2 -4x -9 = 0$.
$\Delta = 16 + 180 = 196 = 14^2$.
$x_1 = \frac{4-14}{10} = -1$, $x_2 = \frac{4+14}{10} = 1,8$.
$y_1 = 2(-1)-1 = -3$, $y_2 = 2(1,8)-1 = 2,6$.
Couples : $(-1;-3)$ et $(1,8; 2,6)$.
\end{corrige}

\begin{corrige}[Exercice 4]
$x+y=40$. Aire totale : $\frac{x^2}{16} + \frac{y^2}{16} = 104 \Leftrightarrow x^2+y^2 = 1664$.
On a $x+y=40 \Rightarrow (x+y)^2 = 1600 = x^2+y^2+2xy = 1664+2xy \Rightarrow 2xy = -64 \Rightarrow xy = -32$.
Système : $\begin{cases} x+y=40 \\ xy=-32 \end{cases}$.
$\Delta = 1600 + 128 = 1728 = 576\times 3 = (24\sqrt{3})^2$.
$x_1 = 20 - 12\sqrt{3} \approx 20 - 20,78 < 0$ (impossible, longueur négative).
$x_2 = 20 + 12\sqrt{3} \approx 40,78 > 40$ (impossible, somme dépasse 40).
\textbf{Conclusion :} Le problème n'a \textbf{pas de solution physique} (longueurs positives). L'énoncé contient une incohérence (somme des aires trop grande pour ce périmètre).
\emph{Note pédagogique :} Vérifier la compatibilité des données est crucial en technique !
\end{corrige}

\coursec{Polynômes de degré 3 : racines et factorisation}

Dans la section précédente, nous avons vu comment certains systèmes d'équations à deux inconnues se ramènent à une équation du second degré. Mais dans la vie professionnelle, toutes les équations ne sont pas du second degré. Un chaudronnier qui calcule le volume d'une cuve, un mécanicien qui étudie le déplacement d'un piston, un gestionnaire qui modélise un bénéfice en fonction de la quantité produite : tous rencontrent des expressions contenant $x^3$. Comment résoudre une équation comme $x^3 - 4x^2 + x + 6 = 0$ ? Il n'existe pas de « discriminant » simple comme pour le second degré. La stratégie, que nous allons construire pas à pas, consiste à \emph{factoriser} le polynôme pour se ramener à des équations que nous savons déjà résoudre : des équations du premier et du second degré.

\courssub{Polynômes de degré 3 : définition et vocabulaire}

\begin{definition}[Polynôme de degré 3]
On appelle \cle{polynôme de degré 3} (ou \cle{polynôme du troisième degré}) toute fonction $P$ définie sur $\mathbb{R}$ par
\[ P(x) = ax^3 + bx^2 + cx + d \]
où $a$, $b$, $c$, $d$ sont des réels avec $a \neq 0$. Les réels $a$, $b$, $c$, $d$ sont les \cle{coefficients} du polynôme ; $d$ est le \cle{terme constant}.
\end{definition}

\begin{exemplebox}[Reconnaître un polynôme de degré 3]
\begin{itemize}
\item $P(x) = 2x^3 - 5x^2 + x - 7$ est un polynôme de degré 3 ($a = 2$, $b = -5$, $c = 1$, $d = -7$).
\item $Q(x) = x^3 + 6$ est un polynôme de degré 3 ($a = 1$, $b = 0$, $c = 0$, $d = 6$) : des coefficients peuvent être nuls, sauf $a$.
\item $R(x) = x^4 - 2x^3 + 1$ n'est \textbf{pas} de degré 3 : il est de degré 4.
\end{itemize}
\end{exemplebox}

\begin{attention}[La condition $a \neq 0$]
Si $a = 0$, le terme en $x^3$ disparaît et le polynôme devient de degré 2 (ou moins). Le degré d'un polynôme est l'\emph{exposant le plus élevé dont le coefficient est non nul}. Avant d'annoncer « polynôme de degré 3 », vérifie toujours que le coefficient de $x^3$ est non nul.
\end{attention}

\courssub{Zéro (ou racine) d'un polynôme}

L'idée est toute simple : remplacer $x$ par un nombre dans $P(x)$ donne un résultat. Parfois, ce résultat vaut $0$. Ce nombre est alors précieux, comme nous le verrons.

\begin{definition}[Zéro d'un polynôme]
Soit $P$ un polynôme et $\alpha$ un nombre réel. On dit que $\alpha$ est un \cle{zéro} (ou une \cle{racine}) de $P$ lorsque
\[ P(\alpha) = 0. \]
Graphiquement, les zéros de $P$ sont les abscisses des points où la courbe de $P$ coupe l'axe des abscisses.
\end{definition}

\begin{exemplebox}[Vérifier qu'un nombre est un zéro]
Soit $P(x) = x^3 - 4x^2 + x + 6$. Testons $x = 2$ :
\[ P(2) = 2^3 - 4 \times 2^2 + 2 + 6 = 8 - 16 + 2 + 6 = 0. \]
Donc $2$ est un zéro de $P$. En revanche $P(1) = 1 - 4 + 1 + 6 = 4 \neq 0$ : $1$ n'est pas un zéro de $P$.
\end{exemplebox}

\begin{methode}[Chercher une racine évidente]
Pour un polynôme à coefficients entiers, on teste en priorité, dans l'ordre :
\begin{enumerate}
\item les valeurs $1$ et $-1$ (calculs rapides) ;
\item les valeurs $2$, $-2$, $3$, $-3$ ;
\item les autres \textbf{diviseurs du terme constant} $d$ (positifs et négatifs).
\end{enumerate}
Dès qu'une valeur $\alpha$ donne $P(\alpha) = 0$, on a trouvé une racine et on peut passer à la factorisation.
\end{methode}

\begin{attention}[Ne pas abandonner trop vite]
Erreur fréquente : tester seulement $1$ et $-1$, constater qu'aucun ne convient, et conclure « le polynôme n'a pas de racine ». C'est faux ! Il faut tester \textbf{tous les diviseurs du terme constant}. Par exemple, pour $P(x) = x^3 - 6x^2 + 11x - 6$, on a $P(1) = 0$... mais pour $Q(x) = x^3 + x - 10$, ni $1$ ni $-1$ ne conviennent : c'est $\alpha = 2$ qui est racine, car $Q(2) = 8 + 2 - 10 = 0$. Le terme constant $-10$ a pour diviseurs $\pm 1, \pm 2, \pm 5, \pm 10$ : il faut les essayer.
\end{attention}

\courssub{Le théorème fondamental de la factorisation}

Voici le résultat central de cette section. Il fait le pont entre « trouver une racine » et « factoriser ».

\begin{propriete}[Factorisation par $(x - \alpha)$]
Soit $P$ un polynôme de degré $n \geq 1$ et $\alpha$ un réel. Alors :
\[ P(\alpha) = 0 \iff P(x) \text{ se factorise par } (x - \alpha), \]
c'est-à-dire qu'il existe un polynôme $Q$ de degré $n - 1$ tel que, pour tout réel $x$ :
\[ P(x) = (x - \alpha)\,Q(x). \]
Pour un polynôme de degré 3, $Q$ est donc un polynôme de degré 2 : $P(x) = (x - \alpha)(a x^2 + b'x + c')$.
\end{propriete}

\begin{experience}[Démonstration guidée du sens direct]
On démontre le sens direct pour un polynôme de degré 3 : si $P(\alpha) = 0$, alors $P(x) = (x - \alpha)(ax^2 + b'x + c')$.

\textbf{Étape 1.} Posons $P(x) = ax^3 + bx^2 + cx + d$ et cherchons $b'$ et $c'$ tels que $P(x) = (x - \alpha)(ax^2 + b'x + c')$. Développons le membre de droite :
\[ (x - \alpha)(ax^2 + b'x + c') = ax^3 + (b' - a\alpha)x^2 + (c' - b'\alpha)x - c'\alpha. \]

\textbf{Étape 2.} Par identification des coefficients avec $ax^3 + bx^2 + cx + d$, on obtient :
\[ b' - a\alpha = b, \qquad c' - b'\alpha = c, \qquad -c'\alpha = d. \]

\textbf{Étape 3.} Les deux premières égalités donnent des candidats :
\[ b' = b + a\alpha, \qquad c' = c + b'\alpha = c + b\alpha + a\alpha^2. \]

\textbf{Étape 4.} Reste à vérifier la troisième condition : $-c'\alpha \stackrel{?}{=} d$. Calculons :
\[ -c'\alpha = -\alpha\left(c + b\alpha + a\alpha^2\right) = -\left(a\alpha^3 + b\alpha^2 + c\alpha\right). \]
Or $P(\alpha) = a\alpha^3 + b\alpha^2 + c\alpha + d = 0$, donc $a\alpha^3 + b\alpha^2 + c\alpha = -d$. Ainsi :
\[ -c'\alpha = -(-d) = d. \quad \checkmark \]

\textbf{Conclusion.} Les coefficients $b' = b + a\alpha$ et $c' = c + b\alpha + a\alpha^2$ conviennent : on a bien $P(x) = (x - \alpha)(ax^2 + b'x + c')$. Le sens réciproque est immédiat : si $P(x) = (x - \alpha)Q(x)$, alors $P(\alpha) = (\alpha - \alpha)Q(\alpha) = 0$.
\end{experience}

\begin{aretenir}[L'idée à retenir]
Trouver une racine $\alpha$ de $P$, c'est gagner le droit d'écrire $P(x) = (x - \alpha) \times (\text{trinôme})$. L'équation $P(x) = 0$ se ramène alors à $x - \alpha = 0$ ou à une équation du \textbf{second degré}, que l'on sait résoudre avec le discriminant.
\end{aretenir}

\courssub{Technique 1 : la division euclidienne par $(x - \alpha)$}

La division euclidienne des polynômes fonctionne comme la division des nombres entiers apprise au primaire : on divise les termes de plus haut degré, on multiplie, on soustrait, et on recommence jusqu'à obtenir un reste nul (puisque $\alpha$ est racine, le reste \emph{doit} être nul).

\begin{exemplebox}[Division de $x^3 - 4x^2 + x + 6$ par $(x - 2)$]
On sait que $P(2) = 0$, donc la division de $P(x) = x^3 - 4x^2 + x + 6$ par $x - 2$ tombe juste.
\begin{itemize}
\item $x^3 \div x = x^2$ ; on multiplie : $x^2(x - 2) = x^3 - 2x^2$ ; on soustrait : il reste $-2x^2 + x + 6$.
\item $-2x^2 \div x = -2x$ ; on multiplie : $-2x(x - 2) = -2x^2 + 4x$ ; on soustrait : il reste $-3x + 6$.
\item $-3x \div x = -3$ ; on multiplie : $-3(x - 2) = -3x + 6$ ; on soustrait : il reste $0$.
\end{itemize}
Le quotient est $x^2 - 2x - 3$ et le reste est nul : $P(x) = (x - 2)(x^2 - 2x - 3)$.
\end{exemplebox}

\begin{popfigure}
\begin{center}
\begin{tikzpicture}[scale=0.95, every node/.style={font=\small}]
% Dividende
\node[anchor=base] at (0,0) {$x^3 - 4x^2 + x + 6$};
% Diviseur avec cadre
\node[anchor=base west] (div) at (4.6,0) {$x - 2$};
\draw[popblue, thick] (4.4,-0.25) -- (4.4,1.6) -- (6.2,1.6);
% Quotient
\node[anchor=base west, popblue] at (4.6,1.0) {$x^2 - 2x - 3$};
% Étapes de la division
\node[anchor=base west, popmuted] at (0.35,-0.7) {$-\,(x^3 - 2x^2)$};
\draw[popline] (-0.1,-1.0) -- (4.3,-1.0);
\node[anchor=base west] at (1.15,-1.5) {$-2x^2 + x + 6$};
\node[anchor=base west, popmuted] at (1.5,-2.2) {$-\,(-2x^2 + 4x)$};
\draw[popline] (1.05,-2.5) -- (4.3,-2.5);
\node[anchor=base west] at (2.3,-3.0) {$-3x + 6$};
\node[anchor=base west, popmuted] at (2.65,-3.7) {$-\,(-3x + 6)$};
\draw[popline] (2.2,-4.0) -- (4.3,-4.0);
\node[anchor=base west, poporange] at (3.3,-4.5) {$0$};
\end{tikzpicture}
\end{center}
\legende{Disposition de la division euclidienne de $x^3 - 4x^2 + x + 6$ par $x - 2$ : le reste est nul, ce qui confirme que $2$ est racine.}
\end{popfigure}

\courssub{Technique 2 : l'identification des coefficients}

Plus rapide à rédiger, cette méthode consiste à écrire directement la forme factorisée avec des coefficients inconnus, à développer, puis à identifier.

\begin{methode}[Factoriser un polynôme de degré 3 en 4 étapes]
\begin{enumerate}
\item \textbf{Chercher une racine évidente} $\alpha$ parmi $\pm 1$, $\pm 2$, $\pm 3$ et les diviseurs du terme constant.
\item \textbf{Vérifier} que $P(\alpha) = 0$ par un calcul explicite (indispensable dans la rédaction).
\item \textbf{Factoriser} par $(x - \alpha)$ : écrire $P(x) = (x - \alpha)(ax^2 + b'x + c')$, développer et identifier les coefficients (ou poser la division euclidienne).
\item \textbf{Résoudre le trinôme} $ax^2 + b'x + c' = 0$ avec le discriminant $\Delta$ pour obtenir la factorisation complète et toutes les solutions.
\end{enumerate}
\end{methode}

\begin{exoresolu}[Factorisation complète de $P(x) = x^3 - 4x^2 + x + 6$]
\textbf{Énoncé.} Soit $P(x) = x^3 - 4x^2 + x + 6$.
\begin{enumerate}
\item Montrer que $2$ est une racine de $P$.
\item Factoriser $P(x)$ en un produit de trois facteurs du premier degré.
\item Résoudre dans $\mathbb{R}$ l'équation $P(x) = 0$.
\end{enumerate}
\tcblower
\textbf{Solution détaillée.}

\textbf{1.} Calculons : $P(2) = 2^3 - 4 \times 2^2 + 2 + 6 = 8 - 16 + 2 + 6 = 0$. Donc $2$ est bien une racine de $P$.

\textbf{2.} D'après le théorème de factorisation, il existe $b'$ et $c'$ tels que :
\[ P(x) = (x - 2)(x^2 + b'x + c'). \]
Développons : $(x - 2)(x^2 + b'x + c') = x^3 + (b' - 2)x^2 + (c' - 2b')x - 2c'$.

Identifions avec $x^3 - 4x^2 + x + 6$ :
\begin{align*}
b' - 2 &= -4 & c' - 2b' &= 1 & -2c' &= 6 \\
b' &= -2 & & & c' &= -3
\end{align*}
Vérifions la cohérence : $c' - 2b' = -3 - 2 \times (-2) = -3 + 4 = 1$. $\checkmark$

Donc $P(x) = (x - 2)(x^2 - 2x - 3)$.

Factorisons le trinôme $x^2 - 2x - 3$ : $\Delta = (-2)^2 - 4 \times 1 \times (-3) = 4 + 12 = 16 > 0$.
\[ x_1 = \frac{2 - 4}{2} = -1, \qquad x_2 = \frac{2 + 4}{2} = 3. \]
D'où la factorisation complète :
\[ \boxed{P(x) = (x - 2)(x + 1)(x - 3)} \]

\textbf{3.} $P(x) = 0 \iff (x + 1)(x - 2)(x - 3) = 0 \iff x = -1$ ou $x = 2$ ou $x = 3$.
\[ \mathcal{S} = \{-1 \,;\, 2 \,;\, 3\}. \]
\end{exoresolu}

\begin{popfigure}
\begin{center}
\begin{tikzpicture}
\begin{axis}[
    width=0.85\linewidth, height=6.5cm,
    axis lines=middle,
    xlabel={$x$}, ylabel={$y$},
    xmin=-2.2, xmax=4.2, ymin=-5, ymax=7,
    xtick={-1,1,2,3}, ytick={-4,-2,2,4,6},
    grid=both, grid style={popcream},
    tick label style={font=\small},
    label style={font=\small},
]
\addplot[popcurve, domain=-1.7:3.7, samples=200]{x^3 - 4*x^2 + x + 6};
\addplot[only marks, mark=*, mark size=2pt, poporange] coordinates {(-1,0) (2,0) (3,0)};
\node[poporange, font=\small, anchor=south] at (axis cs:-1,0.2) {$-1$};
\node[poporange, font=\small, anchor=north west] at (axis cs:2,-0.1) {$2$};
\node[poporange, font=\small, anchor=south west] at (axis cs:3,0.2) {$3$};
\end{axis}
\end{tikzpicture}
\end{center}
\legende{Courbe de $P(x) = x^3 - 4x^2 + x + 6$ : elle coupe l'axe des abscisses exactement aux trois racines $-1$, $2$ et $3$.}
\end{popfigure}

\begin{attention}[Pièges de rédaction]
\begin{itemize}
\item Ne jamais écrire la factorisation sans avoir \textbf{d'abord vérifié} que $P(\alpha) = 0$ : c'est cette vérification qui justifie l'existence du facteur $(x - \alpha)$.
\item Dans $(x - \alpha)$, le signe est un \textbf{moins} : si la racine est $\alpha = -1$, le facteur est $(x - (-1)) = (x + 1)$. Beaucoup d'élèves écrivent $(x - 1)$ pour la racine $-1$ : c'est faux.
\item Après identification, \textbf{toujours vérifier} en redéveloppant mentalement ou en testant une valeur : si $P(0) = 6$ et $(0 - 2)(0 + 1)(0 - 3) = 6$, la factorisation est cohérente.
\item Ne pas oublier la dernière étape : factoriser le trinôme du second degré. S'arrêter à $(x - 2)(x^2 - 2x - 3)$ ne donne pas toutes les solutions de l'équation.
\end{itemize}
\end{attention}

\begin{savaistu}[Pourquoi degré 3 au maximum trois racines ?]
Chaque racine $\alpha$ fournit un facteur $(x - \alpha)$ du premier degré. Un polynôme de degré 3 ne peut donc pas avoir plus de \textbf{trois} racines réelles distinctes : au-delà, le produit des facteurs serait de degré supérieur à 3. Au quotidien, les équations de degré 3 apparaissent dès qu'on manipule des volumes : par exemple, déterminer les dimensions d'un bac de \SI{6}{m^3} dont la longueur dépasse la largeur de \SI{1}{m} conduit précisément à une équation du type $x^3 + x^2 - 6 = 0$, que l'on résout par la méthode de cette section.
\end{savaistu}

\begin{exercice}[À toi de jouer]
Soit $P(x) = x^3 - 6x^2 + 11x - 6$.
\begin{enumerate}
\item Vérifier que $1$ est une racine de $P$.
\item Par identification, factoriser $P(x)$ sous la forme $(x - 1)(x^2 + b'x + c')$.
\item En déduire la factorisation complète de $P$ et résoudre $P(x) = 0$.
\end{enumerate}
\end{exercice}

\begin{corrige}[Exercice : factorisation de $x^3 - 6x^2 + 11x - 6$]
\textbf{1.} $P(1) = 1 - 6 + 11 - 6 = 0$ : $1$ est bien racine de $P$.

\textbf{2.} On écrit $P(x) = (x - 1)(x^2 + b'x + c') = x^3 + (b' - 1)x^2 + (c' - b')x - c'$. Par identification : $b' - 1 = -6$ donc $b' = -5$ ; $-c' = -6$ donc $c' = 6$. Vérification : $c' - b' = 6 - (-5) = 11$. $\checkmark$ D'où $P(x) = (x - 1)(x^2 - 5x + 6)$.

\textbf{3.} Pour $x^2 - 5x + 6$ : $\Delta = 25 - 24 = 1$, $x_1 = \frac{5 - 1}{2} = 2$, $x_2 = \frac{5 + 1}{2} = 3$. Donc $P(x) = (x - 1)(x - 2)(x - 3)$ et $\mathcal{S} = \{1 \,;\, 2 \,;\, 3\}$.
\end{corrige}

La factorisation complète d'un polynôme de degré 3 ne sert pas seulement à résoudre des équations : elle permet aussi d'étudier son \emph{signe} sur $\mathbb{R}$, ce qui est la clé des inéquations de degré 3 que nous aborderons dans la section suivante.

\coursec{Signe d'un polynôme de degré 3 et inéquations}

Dans la section précédente, nous avons appris à \cle{factoriser} un polynôme de degré 3 : grâce à une racine évidente $a$, on écrit $P(x) = (x-a)\,Q(x)$, puis on factorise le trinôme $Q(x)$. Nous savons donc maintenant écrire un polynôme de degré 3 sous la forme d'un \emph{produit} de facteurs du premier ou du second degré. À quoi cela sert-il concrètement ? Imagine un chef d'atelier qui sait que son bénéfice, en milliers de francs, est donné par $B(x) = x^3 - 4x^2 + x + 6$ selon le nombre $x$ de pièces produites. Pour savoir \emph{quand l'atelier est rentable}, il faut résoudre $B(x) \geq 0$ : c'est une \cle{inéquation de degré 3}. Toute cette section vise à maîtriser l'outil qui permet de la résoudre : le \cle{tableau de signes}.

\courssub{Le principe : le signe d'un produit}

\textbf{Intuition.}\par

Le signe d'un nombre factorisé ne dépend que des signes de ses facteurs : un produit est positif si ses facteurs non nuls sont \emph{en nombre pair de négatifs}, et négatif sinon. Par exemple $(-2)\times 3 \times (-5) = 30 > 0$ car il y a deux facteurs négatifs. Pour un polynôme factorisé comme $P(x) = (x+1)(x-2)(x-3)$, le signe de $P(x)$ dépend donc du signe de chacun des trois facteurs, qui eux-mêmes changent de signe en $-1$, $2$ et $3$. Il suffit d'étudier le signe de \emph{chaque facteur séparément}, puis de faire le « produit des signes » ligne par ligne dans un tableau.

\begin{propriete}[Signe d'un produit de facteurs]
Soit $P$ un polynôme écrit sous forme factorisée. Le signe de $P(x)$ s'obtient en appliquant la \cle{règle des signes} aux signes de ses facteurs :
\begin{itemize}
\item si un facteur s'annule, alors $P(x) = 0$ ;
\item sinon, $P(x) > 0$ lorsque le nombre de facteurs strictement négatifs est pair, et $P(x) < 0$ lorsqu'il est impair.
\end{itemize}
\end{propriete}

\courssub{Signe des facteurs de base : binôme et trinôme}

Avant d'assembler le tableau, rappelons le signe des deux types de facteurs que l'on rencontre dans un polynôme de degré 3 factorisé.

\begin{propriete}[Signe du binôme $x - a$]
Pour tout réel $x$ :
\begin{itemize}
\item $x - a < 0$ si $x < a$ ;
\item $x - a = 0$ si $x = a$ ;
\item $x - a > 0$ si $x > a$.
\end{itemize}
Le binôme $x-a$ est donc \textbf{négatif avant $a$, positif après $a$}. (De même, $x + a = x - (-a)$ s'annule en $-a$.)
\end{propriete}

\begin{propriete}[Signe du trinôme $ax^2+bx+c$ ($a \neq 0$)]
Soit $\Delta = b^2 - 4ac$ le discriminant.
\begin{itemize}
\item Si $\Delta > 0$ : le trinôme a deux racines $x_1 < x_2$ et est \textbf{du signe de $a$ à l'extérieur des racines}, du signe contraire entre les racines.
\item Si $\Delta = 0$ : le trinôme a une racine double $x_0$ et est \textbf{toujours du signe de $a$} (nul seulement en $x_0$).
\item Si $\Delta < 0$ : le trinôme est \textbf{toujours du signe de $a$}, pour tout réel $x$.
\end{itemize}
\end{propriete}

\begin{attention}[Le cas $\Delta < 0$ simplifie tout]
Quand un polynôme de degré 3 se factorise en $P(x) = (x-a)(x^2 + bx + c)$ avec un trinôme de discriminant strictement négatif et de coefficient dominant positif, le trinôme est \emph{toujours strictement positif} : le signe de $P(x)$ est alors exactement celui du binôme $x - a$. Inutile de chercher des racines qui n'existent pas !
\end{attention}

\courssub{Une propriété fondamentale (admise)}

\begin{propriete}[Factorisation toujours possible — admise]
Tout polynôme de degré 3 à coefficients réels admet \cle{au moins une racine réelle}. Par conséquent, tout polynôme de degré 3 se factorise \emph{toujours} au moins partiellement dans $\mathbb{R}$, sous l'une des formes :
\begin{itemize}
\item $P(x) = a(x - \alpha)(x - \beta)(x - \gamma)$ (trois racines réelles, éventuellement confondues) ;
\item $P(x) = a(x - \alpha)(x^2 + bx + c)$ avec $\Delta = b^2 - 4ac < 0$ (une seule racine réelle).
\end{itemize}
Cette propriété est admise au programme de Première technique ; elle garantit que la méthode du tableau de signes s'applique à \emph{tous} les polynômes de degré 3.
\end{propriete}

\begin{savaistu}[Et la formule magique pour le 3e degré ?]
Au XVI\textsuperscript{e} siècle, les mathématiciens italiens Tartaglia et \cle{Cardan} ont découvert une formule générale donnant les solutions de toute équation du troisième degré, la célèbre \emph{formule de Cardan}. Elle est spectaculaire... mais bien trop compliquée pour le programme du Probatoire ! Bonne nouvelle : au Probatoire technique, les polynômes proposés possèdent toujours une \cle{racine évidente} ($-2$, $-1$, $0$, $1$, $2$, $3$...) que l'on trouve en testant quelques valeurs. La méthode « racine évidente + factorisation par $x - a$ » suffit donc dans tous les cas rencontrés à l'examen.
\end{savaistu}

\courssub{Construction du tableau de signes : exemple complet}

Étudions le signe de $P(x) = (x+1)(x-2)(x-3)$. Les trois facteurs s'annulent respectivement en $-1$, $2$ et $3$. On range ces valeurs dans l'ordre croissant sur la première ligne, on met une ligne par facteur, puis une ligne « produit » pour $P(x)$.

\begin{center}
\begin{tabular}{|c|ccccccc|}
\hline
$x$ & $-\infty$ & & $-1$ & & $2$ & & $3$ \\
\hline
$x+1$ & & $-$ & $0$ & $+$ & $+$ & $+$ & $+$ \\
\hline
$x-2$ & & $-$ & $-$ & $-$ & $0$ & $+$ & $+$ \\
\hline
$x-3$ & & $-$ & $-$ & $-$ & $-$ & $-$ & $0$ \\
\hline
$P(x)$ & & $-$ & $0$ & $+$ & $0$ & $-$ & $0$ \\
\hline
\end{tabular}
\end{center}

\textbf{Lecture.} Sur $]-\infty\,;\,-1[$, les trois facteurs sont négatifs : produit négatif. Sur $]-1\,;\,2[$, un seul facteur ($x+1$) est positif... attention, deux facteurs restent négatifs : $(+)\times(-)\times(-) = (+)$, produit positif. Sur $]2\,;\,3[$, un seul facteur négatif : produit négatif. Sur $]3\,;\,+\infty[$, tout est positif.

\begin{popfigure}
\begin{center}
\begin{tikzpicture}
\begin{axis}[
  width=0.85\linewidth, height=6.2cm,
  axis lines=middle, xlabel={$x$}, ylabel={$y$},
  xmin=-2.2, xmax=4.4, ymin=-6, ymax=7,
  xtick={-1,2,3}, ytick=\empty,
  domain=-1.8:3.9, samples=200, clip=false]
  \addplot[popcurve]{(x+1)*(x-2)*(x-3)};
  \addplot[poporange, very thick, domain=-1:2, samples=100]{(x+1)*(x-2)*(x-3)};
  \addplot[poporange, very thick, domain=3:3.9, samples=60]{(x+1)*(x-2)*(x-3)};
  \addplot[only marks, mark=*, popdark, mark size=1.6pt] coordinates {(-1,0) (2,0) (3,0)};
  \node[poporange, font=\small] at (axis cs:0.5,4.2) {$P(x)\geq 0$};
  \node[poporange, font=\small] at (axis cs:3.55,4.2) {$P(x)\geq 0$};
  \node[popblue, font=\small] at (axis cs:2.5,-3.6) {$P(x)\leq 0$};
\end{axis}
\end{tikzpicture}
\end{center}
\legende{Courbe de $P(x) = (x+1)(x-2)(x-3)$. En orange, les morceaux où $P(x) \geq 0$ : la courbe est au-dessus de l'axe des abscisses exactement sur $[-1\,;\,2] \cup [3\,;\,+\infty[$.}
\end{popfigure}

\begin{aretenir}[Lecture graphique du signe]
Le signe d'un polynôme se lit sur sa courbe :
\begin{itemize}
\item $P(x) > 0$ là où la courbe est \textbf{au-dessus} de l'axe des abscisses ;
\item $P(x) < 0$ là où la courbe est \textbf{en dessous} ;
\item les racines de $P$ sont les abscisses des points où la courbe \textbf{coupe} l'axe (ou le touche, pour une racine double).
\end{itemize}
Graphique et tableau de signes disent exactement la même chose : l'un vérifie l'autre.
\end{aretenir}

\courssub{Le cas des racines multiples}

\begin{propriete}[Racine double : le signe ne change pas]
Si $a$ est une \cle{racine double} de $P$, le facteur $(x-a)$ apparaît au carré : $P(x) = (x-a)^2 \times (\text{autres facteurs})$. Or pour tout $x \neq a$, $(x-a)^2 > 0$. Donc :
\begin{itemize}
\item $P(a) = 0$ ;
\item le signe de $P(x)$ est le \textbf{même avant et après $a$} : le polynôme ne change pas de signe en $a$.
\end{itemize}
En revanche, à une racine \emph{simple} (ou triple), le signe change.
\end{propriete}

\begin{experience}[Découverte : pourquoi le signe ne change-t-il pas ?]
Considérons $P(x) = (x-1)^2(x-3)$.
\begin{enumerate}
\item Calcule $P(0)$, $P(2)$, $P(4)$. Que constates-tu pour le signe de part et d'autre de $1$ ? Et de part et d'autre de $3$ ?
\item Justifie : pour $x \neq 1$, quel est le signe de $(x-1)^2$ ? Le signe de $P(x)$ dépend donc uniquement de quel facteur ?
\item Trace la courbe de $P$ avec un logiciel ou une calculatrice : observe qu'en $x = 1$ la courbe \emph{touche} l'axe sans le traverser (elle « rebondit »), tandis qu'en $x = 3$ elle le traverse franchement.
\end{enumerate}
Conclusion : une racine double se reconnaît graphiquement à un \emph{contact} avec l'axe, une racine simple à un \emph{franchissement}.
\end{experience}

\begin{attention}[Erreur fréquente n°1 : changer le signe à une racine double]
Beaucoup d'élèves changent mécaniquement le signe à \emph{chaque} racine rencontrée dans le tableau. C'est faux pour une racine double ! Le facteur $(x-a)^2$ est un carré : il est \textbf{toujours positif ou nul}. Dans le tableau, la ligne de $(x-a)^2$ ne contient que des « $+$ » (et un $0$ en $a$). Par exemple, pour $P(x) = (x-1)^2(x-3)$ : $P(x) \leq 0$ sur $]-\infty\,;\,3]$ et $P(x) \geq 0$ sur $\{1\} \cup [3\,;\,+\infty[$, la racine $1$ ne jouant aucun rôle dans le changement de signe (seulement $P(1) = 0$).
\end{attention}

\courssub{Résolution d'inéquations de degré 3}

\begin{methode}[Résoudre $P(x) \geq 0$ (ou $>$, $\leq$, $<$) en 3 étapes]
\begin{enumerate}
\item \textbf{Factoriser complètement} $P(x)$ : chercher une racine évidente $a$ (tester $-2, -1, 0, 1, 2, 3$), écrire $P(x) = (x-a)Q(x)$, puis factoriser le trinôme $Q$ (discriminant).
\item \textbf{Construire le tableau de signes} : une ligne par facteur, les racines rangées dans l'ordre croissant, produit des signes sur la dernière ligne. Attention aux racines doubles !
\item \textbf{Conclure par une réunion d'intervalles} : lire les colonnes où le signe convient. Bornes \emph{fermées} (crochets) si l'inégalité est large ($\geq$ ou $\leq$), \emph{ouvertes} si elle est stricte ($>$ ou $<$).
\end{enumerate}
\end{methode}

\begin{exoresolu}[Résoudre $x^3 - 4x^2 + x + 6 \geq 0$]
\textbf{Énoncé.} Résoudre dans $\mathbb{R}$ l'inéquation $x^3 - 4x^2 + x + 6 \geq 0$.
\tcblower
\textbf{Solution détaillée.}

\emph{Étape 1 : factorisation.} On teste les valeurs simples dans $P(x) = x^3 - 4x^2 + x + 6$ :
\[ P(-1) = -1 - 4 - 1 + 6 = 0. \]
Donc $-1$ est racine et $P(x) = (x+1)(x^2 + bx + c)$. Par identification (ou division), on obtient :
\[ P(x) = (x+1)(x^2 - 5x + 6). \]
Pour le trinôme $x^2 - 5x + 6$ : $\Delta = 25 - 24 = 1 > 0$, d'où $x_1 = \dfrac{5-1}{2} = 2$ et $x_2 = \dfrac{5+1}{2} = 3$. Ainsi :
\[ P(x) = (x+1)(x-2)(x-3). \]

\emph{Étape 2 : tableau de signes.} C'est exactement le tableau construit plus haut : $P(x) \geq 0$ sur $[-1\,;\,2]$ et sur $[3\,;\,+\infty[$.

\emph{Étape 3 : conclusion.} L'inégalité est large, donc on ferme les crochets aux racines :
\[ \mathcal{S} = [-1\,;\,2] \cup [3\,;\,+\infty[. \]
\emph{Vérification graphique} : sur la figure ci-dessus, la courbe est bien au-dessus de l'axe (partie orange) précisément sur ces deux intervalles.
\end{exoresolu}

\begin{exemplebox}[Avec une racine double]
Résoudre $x^3 - 5x^2 + 7x - 3 < 0$. On trouve $P(1) = 0$, puis $P(x) = (x-1)(x^2 - 4x + 3) = (x-1)^2(x-3)$. Le carré $(x-1)^2$ est toujours positif (nul en $1$), donc $P(x)$ a le signe de $(x-3)$, sauf en $1$ où $P(1) = 0$. L'inégalité stricte donne :
\[ \mathcal{S} = \left]-\infty\,;\,1\right[ \cup \left]1\,;\,3\right[. \]
Remarque le piège : $1$ est exclu car $P(1) = 0$, et l'inégalité est stricte.
\end{exemplebox}

\courssub{Complément : inéquations quotients}

\begin{propriete}[Signe d'un quotient]
Le signe d'un quotient $\dfrac{A(x)}{B(x)}$ se détermine comme celui du produit $A(x)\times B(x)$, \textbf{à une différence près} : les valeurs qui annulent le dénominateur $B$ sont des \cle{valeurs interdites} — elles sont toujours exclues de l'ensemble solution, même si l'inégalité est large. On les signale par une \emph{double barre} dans le tableau de signes.
\end{propriete}

\begin{exemplebox}[Un quotient se ramène à un produit]
Résoudre $\dfrac{x-1}{x-3} \geq 0$. Valeur interdite : $x = 3$. Le quotient a le signe du produit $(x-1)(x-3)$, positif à l'extérieur des racines. Mais $3$ est interdit (division par zéro), donc :
\[ \mathcal{S} = \left]-\infty\,;\,1\right] \cup \left]3\,;\,+\infty\right[. \]
Le crochet est fermé en $1$ (le numérateur peut s'annuler) mais ouvert en $3$. Cette technique sera très utile en Terminale pour l'étude des fonctions rationnelles.
\end{exemplebox}

\begin{attention}[Les trois pièges classiques du tableau de signes]
\begin{enumerate}
\item \textbf{Oublier de factoriser d'abord} : on ne peut pas faire le signe de $x^3 - 4x^2 + x + 6$ directement ; il faut impérativement la forme factorisée.
\item \textbf{Inverser les signes du binôme} : $x - a$ est négatif \emph{avant} $a$ et positif \emph{après}. Écris toujours le facteur sous la forme $x - a$ pour repérer $a$ correctement : $x + 1$ s'annule en $-1$, pas en $1$ !
\item \textbf{Mal fermer les crochets} : inégalité large $\Rightarrow$ racines incluses ; inégalité stricte $\Rightarrow$ racines exclues ; valeur interdite (dénominateur) $\Rightarrow$ toujours exclue.
\end{enumerate}
\end{attention}

\begin{exercice}
\begin{enumerate}
\item Dresser le tableau de signes de $P(x) = (x+2)(x-1)(x-4)$, puis résoudre $P(x) > 0$.
\item Résoudre dans $\mathbb{R}$ : $x^3 - 2x^2 - 5x + 6 \leq 0$ (on vérifiera que $1$ est racine).
\item Résoudre $x^3 - 3x^2 + 4 > 0$ sachant que $x^3 - 3x^2 + 4 = (x+1)(x-2)^2$.
\item Un menuisier estime son bénéfice mensuel (en milliers de FCFA) par $B(x) = x^3 - 6x^2 + 9x$, où $x$ est le nombre de meubles vendus. Pour quelles valeurs de $x$ (entières, entre $0$ et $6$) l'atelier est-il bénéficiaire ?
\end{enumerate}
\end{exercice}

\begin{corrige}[Exercice]
\begin{enumerate}
\item Les racines sont $-2$, $1$, $4$. Produit des signes : $P(x) > 0$ sur $]-2\,;\,1[ \cup ]4\,;\,+\infty[$.
\item $P(1) = 1 - 2 - 5 + 6 = 0$, puis $P(x) = (x-1)(x^2 - x - 6) = (x-1)(x+2)(x-3)$. Signe négatif à l'extérieur : $\mathcal{S} = ]-\infty\,;\,-2] \cup [1\,;\,3]$.
\item $(x-2)^2 \geq 0$ toujours, nul en $2$ ; le signe est celui de $(x+1)$, sauf en $2$ où le polynôme vaut $0$. Donc $\mathcal{S} = ]-1\,;\,2[ \cup ]2\,;\,+\infty[$.
\item $B(x) = x(x^2 - 6x + 9) = x(x-3)^2$. Pour $x \geq 0$, $B(x) \geq 0$ toujours (carré), et $B(x) > 0$ pour $x > 0$ et $x \neq 3$. L'atelier est bénéficiaire pour $x \in \{1\,;\,2\,;\,4\,;\,5\,;\,6\}$ meubles vendus (en $x = 3$, le bénéfice est nul).
\end{enumerate}
\end{corrige}

\begin{aretenir}[L'essentiel de la section]
\begin{itemize}
\item Le signe d'un polynôme factorisé s'obtient par la \cle{règle des signes} appliquée à ses facteurs, dans un tableau de signes.
\item $x - a$ : négatif avant $a$, positif après. Un trinôme sans racine réelle garde un signe constant.
\item À une \cle{racine simple}, le signe change ; à une \cle{racine double}, il ne change pas (carré toujours positif).
\item Résoudre $P(x) \geq 0$ : factoriser, tableau de signes, conclure par une réunion d'intervalles avec les bons crochets.
\item Graphiquement, $P(x) \geq 0$ correspond aux portions de courbe situées au-dessus de l'axe des abscisses.
\end{itemize}
\end{aretenir}

\coursec{Équations irrationnelles}

Après avoir étudié le signe des polynômes de degré 3 et résolu les inéquations associées, nous abordons maintenant les équations où l'inconnue apparaît sous un radical. Ces équations, appelées \cle{équations irrationnelles}, nécessitent une attention particulière aux conditions d'existence et à la vérification des solutions, car l'élévation au carré peut introduire des solutions parasites.

\courssub{Définition et domaine de validité}

\begin{definition}[Équation irrationnelle et domaine de validité]
Une \textbf{équation irrationnelle} est une équation dans laquelle l'inconnue figure à l'intérieur d'un radical (généralement une racine carrée).\\
Le \textbf{domaine de validité} (ou ensemble de définition) est l'ensemble des valeurs de l'inconnue pour lesquelles toutes les expressions sous radicaux (les radicandes) sont définies dans $\mathbb{R}$, c'est-à-dire positives ou nulles.
\end{definition}

\emph{Exemple :} Dans $\sqrt{x+2}=x$, le radicande $x+2$ doit être $\ge 0$, donc $x\ge -2$. De plus, le second membre $x$ doit être $\ge 0$ car une racine carrée est toujours positive ou nulle. Le domaine de validité est donc $x\ge 0$.

\courssub{Résolution de $\sqrt{A}=B$}

\begin{propriete}[Équivalence pour $\sqrt{A}=B$]
Soient $A$ et $B$ deux expressions dépendant de $x$. Pour $B\ge 0$, on a l'équivalence :
\[
\sqrt{A}=B \;\Longleftrightarrow\; \begin{cases} A = B^2 \\ B \ge 0 \end{cases}
\]
\emph{Preuve.} 
\begin{itemize}
\item Si $\sqrt{A}=B$, alors par définition de la racine carrée, $B\ge 0$ et en élevant au carré on obtient $A=B^2$.
\item Réciproquement, si $A=B^2$ et $B\ge 0$, alors $\sqrt{A}=\sqrt{B^2}=|B|=B$ car $B\ge 0$.
\end{itemize}
La condition $B\ge 0$ est \textbf{indispensable} : sans elle, l'élévation au carré perd l'information du signe et peut créer des solutions qui ne satisfont pas l'équation originale (solutions parasites).
\end{propriete}

\begin{popfigure}
\begin{tikzpicture}[
    decision/.style={diamond, draw=popblue, thick, fill=popblueL, text width=3.5cm, align=center, inner sep=2pt},
    process/.style={rectangle, draw=poporange, thick, fill=poporangeL, text width=3.5cm, align=center, inner sep=2pt, rounded corners},
    startstop/.style={ellipse, draw=popgreen, thick, fill=popgreenL, text width=3cm, align=center, inner sep=2pt},
    arrow/.style={->, thick, >=Stealth}
]
\node[startstop] (start) {Début : $\sqrt{A}=B$};
\node[decision, below=1.2cm of start] (cond) {Condition $B\ge 0$ satisfaite ?};
\node[process, below=1.2cm of cond] (square) {Élever au carré : $A=B^2$};
\node[process, below=1.2cm of square] (solve) {Résoudre l'équation obtenue};
\node[decision, below=1.2cm of solve] (check) {Vérifier chaque solution dans $\sqrt{A}=B$};
\node[startstop, below=1.2cm of check] (end) {Solutions finales};
\node[startstop, left=3.5cm of cond] (nosol) {Aucune solution (ensemble vide)};
\draw[arrow] (start) -- (cond);
\draw[arrow] (cond) -- node[left] {Non} (nosol);
\draw[arrow] (cond) -- node[right] {Oui} (square);
\draw[arrow] (square) -- (solve);
\draw[arrow] (solve) -- (check);
\draw[arrow] (check) -- node[left] {Rejeter} ++(-3,0) |- (solve);
\draw[arrow] (check) -- node[right] {Valider} (end);
\end{tikzpicture}
\legende{Organigramme de la méthode de résolution de $\sqrt{A}=B$}
\end{popfigure}

\courssub{Résolution de $\sqrt{A}=\sqrt{B}$}

\begin{propriete}[Équivalence pour $\sqrt{A}=\sqrt{B}$]
\[
\sqrt{A}=\sqrt{B} \;\Longleftrightarrow\; \begin{cases} A = B \\ A \ge 0 \end{cases}
\]
\emph{Justification.} Les deux radicaux sont définis si et seulement si $A\ge 0$ et $B\ge 0$. Comme la fonction racine carrée est injective sur $[0,+\infty[$, l'égalité des radicaux entraîne l'égalité des radicandes, et la condition $A\ge 0$ (qui implique $B\ge 0$ par l'égalité $A=B$) suffit.
\end{propriete}

\courssub{Résolution de $\sqrt{ax+b}=cx+d$}

\begin{methode}[Résolution de $\sqrt{ax+b}=cx+d$ en 4 étapes]
\begin{enumerate}
\item \textbf{Conditions d'existence :} $ax+b \ge 0$ et $cx+d \ge 0$. On détermine l'intersection de ces conditions (domaine de validité).
\item \textbf{Élévation au carré :} $(ax+b) = (cx+d)^2$, ce qui donne une équation du second degré.
\item \textbf{Résolution :} Résoudre l'équation du second degré (factorisation, discriminant, formules).
\item \textbf{Vérification :} Tester chaque solution dans l'équation originale $\sqrt{ax+b}=cx+d$ ; ne conserver que celles qui la satisfont (elles appartiennent obligatoirement au domaine de validité).
\end{enumerate}
\end{methode}

\courssub{Exemple résolu : $\sqrt{x+2}=x$}

\begin{exoresolu}[Résolution de $\sqrt{x+2}=x$]
\textbf{Énoncé.} Résoudre dans $\mathbb{R}$ l'équation $\sqrt{x+2}=x$.

\textbf{Solution détaillée.}
\begin{enumerate}
\item \textbf{Conditions d'existence :}
\begin{align*}
x+2 &\ge 0 \quad\Rightarrow\quad x \ge -2 \\
x &\ge 0 \quad\text{(car le second membre doit être $\ge 0$)}
\end{align*}
L'intersection donne le domaine de validité : $x \ge 0$.

\item \textbf{Élévation au carré (équivalence grâce à $x\ge 0$) :}
\[
\sqrt{x+2}=x \;\Longleftrightarrow\; x+2 = x^2 \quad\text{avec } x\ge 0.
\]

\item \textbf{Résolution de l'équation du second degré :}
\[
x^2 - x - 2 = 0 \;\Longleftrightarrow\; (x-2)(x+1)=0
\]
Les solutions candidates sont $x=2$ et $x=-1$.

\item \textbf{Vérification dans l'équation originale :}
\begin{itemize}
\item Pour $x=2$ : $\sqrt{2+2}=\sqrt{4}=2$ \quad$\checkmark$
\item Pour $x=-1$ : $\sqrt{-1+2}=\sqrt{1}=1 \neq -1$ \quad$\times$ (rejetée car $-1\not\ge 0$)
\end{itemize}
\end{enumerate}
\textbf{Conclusion.} L'unique solution est $\boxed{x=2}$.
\end{exoresolu}

\begin{popfigure}
\begin{tikzpicture}
\begin{axis}[
    width=\linewidth,
    height=8cm,
    axis lines=middle,
    xlabel=$x$, ylabel=$y$,
    xmin=-2.5, xmax=4,
    ymin=-1, ymax=4,
    xtick={-2,-1,0,1,2,3},
    ytick={0,1,2,3},
    grid=major,
    grid style={popline},
    tick label style={font=\small},
    label style={font=\small},
    legend style={at={(0.98,0.02)}, anchor=south east, font=\small, fill=popcream, draw=popmuted}
]
\addplot[popcurve, domain=-2:4, samples=200] {sqrt(max(0,x+2))};
\addlegendentry{$y=\sqrt{x+2}$}
\addplot[popcurve, poporange, domain=-2:4, samples=100] {x};
\addlegendentry{$y=x$}
\addplot[only marks, mark=*, mark size=2.5, popgreen] coordinates {(2,2)};
\addlegendentry{Intersection $(2,2)$}
\end{axis}
\end{tikzpicture}
\legende{Représentation graphique de $y=\sqrt{x+2}$ et $y=x$ ; l'unique intersection valide est $(2,2)$}
\end{popfigure}

\courssub{Pièges fréquents et erreurs à éviter}

\begin{attention}[Solutions parasites et omission des conditions]
\begin{itemize}
\item \textbf{Oublier la condition $B\ge 0$ (ou $cx+d\ge 0$).} Élever au carré sans cette condition transforme $\sqrt{A}=B$ en $A=B^2$, qui admet aussi les solutions de $\sqrt{A}=-B$. Ces solutions sont \textbf{parasites} et doivent être rejetées.
\item \textbf{Ne pas vérifier les solutions dans l'équation originale.} Même si l'on a posé les conditions, une erreur de calcul peut produire une valeur qui ne satisfait pas l'équation de départ. La vérification est la seule garantie.
\item \textbf{Confondre $\sqrt{A^2}=A$ et $\sqrt{A^2}=|A|$.} La racine carrée renvoie toujours la valeur positive ou nulle.
\end{itemize}
\end{attention}

\courssub{Exercices d'application}

\begin{exercice}[Entraînement]
Résoudre dans $\mathbb{R}$ les équations suivantes :
\begin{enumerate}
\item $\sqrt{2x-3}=x-1$
\item $\sqrt{x^2-4}=x-2$
\item $\sqrt{3x+1}=2x-1$
\item $\sqrt{x+5}=\sqrt{2x-1}$
\end{enumerate}
\end{exercice}

\begin{corrige}[Exercice 1]
\begin{enumerate}
\item \textbf{Conditions :} $2x-3\ge0 \Rightarrow x\ge\frac32$ et $x-1\ge0 \Rightarrow x\ge1$ $\Rightarrow$ $x\ge\frac32$.\\
Équation au carré : $2x-3=(x-1)^2 \Rightarrow x^2-4x+4=0 \Rightarrow (x-2)^2=0 \Rightarrow x=2$.\\
Vérification : $\sqrt{2\cdot2-3}=\sqrt{1}=1$ et $2-1=1$ \quad$\checkmark$. Solution : $\boxed{x=2}$.
\item \textbf{Conditions :} $x^2-4\ge0 \Rightarrow x\le-2$ ou $x\ge2$ ; $x-2\ge0 \Rightarrow x\ge2$ $\Rightarrow$ $x\ge2$.\\
Équation au carré : $x^2-4=(x-2)^2 \Rightarrow x^2-4=x^2-4x+4 \Rightarrow 4x=8 \Rightarrow x=2$.\\
Vérification : $\sqrt{4-4}=0$ et $2-2=0$ \quad$\checkmark$. Solution : $\boxed{x=2}$.
\item \textbf{Conditions :} $3x+1\ge0 \Rightarrow x\ge-\frac13$ ; $2x-1\ge0 \Rightarrow x\ge\frac12$ $\Rightarrow$ $x\ge\frac12$.\\
Équation au carré : $3x+1=(2x-1)^2 \Rightarrow 4x^2-7x=0 \Rightarrow x(4x-7)=0 \Rightarrow x=0$ ou $x=\frac74$.\\
Vérification : $x=0$ rejeté ($0\not\ge\frac12$) ; $x=\frac74$ : $\sqrt{3\cdot\frac74+1}=\sqrt{\frac{25}{4}}=\frac52$ et $2\cdot\frac74-1=\frac52$ \quad$\checkmark$. Solution : $\boxed{x=\frac74}$.
\item \textbf{Conditions :} $x+5\ge0 \Rightarrow x\ge-5$ ; $2x-1\ge0 \Rightarrow x\ge\frac12$ $\Rightarrow$ $x\ge\frac12$.\\
Équivalence : $x+5=2x-1 \Rightarrow x=6$.\\
Vérification : $\sqrt{6+5}=\sqrt{11}$ et $\sqrt{2\cdot6-1}=\sqrt{11}$ \quad$\checkmark$. Solution : $\boxed{x=6}$.
\end{enumerate}
\end{corrige}

\courssub{Synthèse}

\begin{aretenir}[Points clés à retenir]
\begin{itemize}
\item Une équation irrationnelle contient l'inconnue sous un radical.
\item Le domaine de validité impose que tous les radicandes soient $\ge 0$ et, pour $\sqrt{A}=B$, que $B\ge 0$.
\item La résolution suit le schéma : \textbf{conditions} $\rightarrow$ \textbf{élévation au carré} $\rightarrow$ \textbf{résolution algébrique} $\rightarrow$ \textbf{vérification}.
\item L'élévation au carré n'est une équivalence que sous la condition de positivité du second membre ; sinon elle introduit des solutions parasites.
\item La vérification finale dans l'équation d'origine est \textbf{obligatoire} pour éliminer toute solution parasite.
\end{itemize}
\end{aretenir}

\begin{savaistu}[Applications techniques]
Les équations irrationnelles apparaissent naturellement en électricité (loi d'Ohm avec racine carrée pour la puissance), en mécanique (calcul de périodes d'oscillation, vitesses), en géométrie (longueurs, théorème de Pythagore) et en gestion (formules de rentabilité avec taux actuariels). Maîtriser leur résolution rigoureuse est indispensable pour modéliser et résoudre des problèmes concrets de l'atelier ou du chantier.
\end{savaistu}

\coursec{Inéquations irrationnelles et applications}

Dans la section précédente, nous avons appris à résoudre des \cle{équations irrationnelles} comme $\sqrt{x+3} = x+1$ : la clé était de se ramener, par élévation au carré sous conditions, à une équation polynomiale, puis de \emph{vérifier} les solutions. Nous abordons maintenant le cas des \cle{inéquations irrationnelles}, c'est-à-dire celles où l'inconnue apparaît sous un radical et où l'égalité est remplacée par une inégalité : $\sqrt{A} \leq B$, $\sqrt{A} \geq B$ ou encore $\sqrt{A} \leq \sqrt{B}$. L'idée reste la même — se débarrasser du radical — mais une difficulté nouvelle apparaît : \textbf{élever au carré une inégalité n'est pas toujours licite}. Toute cette section repose sur une question simple : \emph{le membre de droite est-il positif ou négatif ?}

\courssub{Le point de départ : le signe de la racine carrée}

Tout repose sur deux faits élémentaires mais décisifs.

\begin{aretenir}[Deux faits fondamentaux]
\begin{enumerate}
\item Une racine carrée n'est définie que si son contenu est positif ou nul : $\sqrt{A}$ existe si et seulement si $A \geq 0$.
\item Une racine carrée est toujours \textbf{positive ou nulle} : pour tout $A \geq 0$, on a $\sqrt{A} \geq 0$.
\end{enumerate}
\end{aretenir}

Conséquence immédiate : si l'on compare $\sqrt{A}$ à un nombre $B$, le signe de $B$ change tout. Par exemple, l'inéquation $\sqrt{x} \leq -3$ n'a \emph{aucune} solution, car une racine carrée ne peut pas être inférieure à un nombre strictement négatif. À l'inverse, $\sqrt{x} \geq -3$ est vraie pour \emph{tout} $x$ du domaine, car une racine carrée est toujours supérieure ou égale à un nombre négatif.

\begin{attention}[Erreur fréquente : élever au carré sans précaution]
L'élévation au carré \textbf{ne conserve l'ordre que pour des nombres positifs ou nuls}. Ainsi :
\begin{itemize}
\item $2 \leq 3$ entraîne bien $4 \leq 9$ (les deux membres sont positifs) ;
\item mais $-3 \leq 2$ alors que $(-3)^2 = 9 > 4 = 2^2$ : l'ordre est renversé !
\end{itemize}
On ne peut donc écrire $\sqrt{A} \leq B \Rightarrow A \leq B^2$ que si l'on a \textbf{d'abord vérifié que $B \geq 0$}. Oublier cette condition est l'erreur la plus fréquente au Probatoire.
\end{attention}

\courssub{Type 1 : inéquations de la forme $\sqrt{A} \leq B$}

Raisonnons d'abord intuitivement. Dire que $\sqrt{A}$ est \emph{plus petite} que $B$, c'est exiger deux choses : que $\sqrt{A}$ existe (donc $A \geq 0$), et que $B$ soit assez grand. Si $B$ est strictement négatif, c'est impossible. Si $B$ est positif ou nul, les deux membres sont positifs ou nuls et l'élévation au carré conserve l'ordre.

\begin{propriete}[Inéquation $\sqrt{A} \leq B$ (admise, illustrée graphiquement)]
Soient $A$ et $B$ deux expressions dépendant de $x$.
\begin{itemize}
\item Si $B < 0$ : l'inéquation $\sqrt{A} \leq B$ \textbf{n'a pas de solution}.
\item Si $B \geq 0$ : l'inéquation équivaut au système
\[ \begin{cases} A \geq 0 \\ A \leq B^2 \end{cases} \]
\end{itemize}
Cette propriété est admise ; elle se justifie par la croissance stricte de la fonction carré sur $[0\,;\,+\infty[$.
\end{propriete}

\begin{exemplebox}[Résolution complète de $\sqrt{x+3} \leq x+1$]
\textbf{Étape 1 — Domaine.} Il faut $x + 3 \geq 0$, c'est-à-dire $x \geq -3$.

\textbf{Étape 2 — Signe de $B = x+1$.} Si $x + 1 < 0$, c'est-à-dire $x < -1$, l'inéquation n'a pas de solution.

\textbf{Étape 3 — Cas $B \geq 0$, c'est-à-dire $x \geq -1$.} Le système équivalent est :
\[ \begin{cases} x + 3 \geq 0 \\ x + 3 \leq (x+1)^2 \end{cases} \]
La deuxième condition s'écrit :
\begin{align*}
x + 3 &\leq x^2 + 2x + 1 \\
0 &\leq x^2 + x - 2
\end{align*}
Le trinôme $x^2 + x - 2$ a pour discriminant $\Delta = 1 + 8 = 9$, ses racines sont $x_1 = -2$ et $x_2 = 1$. Étant positif à l'extérieur des racines, on obtient $x \leq -2$ ou $x \geq 1$.

\textbf{Étape 4 — Intersection.} On intersecte avec les conditions $x \geq -3$ et $x \geq -1$ : il reste $x \geq 1$.

\textbf{Conclusion :} $S = [1\,;\,+\infty[$.
\end{exemplebox}

\courssub{Type 2 : inéquations de la forme $\sqrt{A} \geq B$}

Ici, $\sqrt{A}$ doit être \emph{plus grande} que $B$. Si $B$ est strictement négatif, c'est automatiquement vrai dès que $\sqrt{A}$ existe : il suffit donc que $A \geq 0$. Si $B$ est positif ou nul, les deux membres sont positifs ou nuls et on peut élever au carré.

\begin{propriete}[Inéquation $\sqrt{A} \geq B$ (admise)]
Soient $A$ et $B$ deux expressions dépendant de $x$.
\begin{itemize}
\item Si $B < 0$ : l'inéquation $\sqrt{A} \geq B$ équivaut à $A \geq 0$.
\item Si $B \geq 0$ : l'inéquation équivaut à $A \geq B^2$ (la condition $A \geq 0$ est alors automatique, car $A \geq B^2 \geq 0$).
\end{itemize}
\end{propriete}

\begin{exemplebox}[Résolution de $\sqrt{2x+1} \geq x - 1$]
\textbf{Cas 1 : $x - 1 < 0$, c'est-à-dire $x < 1$.} Il suffit que $2x + 1 \geq 0$, soit $x \geq -\dfrac{1}{2}$. D'où $x \in \left[-\dfrac{1}{2}\,;\,1\right[$.

\textbf{Cas 2 : $x - 1 \geq 0$, c'est-à-dire $x \geq 1$.} On élève au carré :
\begin{align*}
2x + 1 &\geq (x-1)^2 \\
2x + 1 &\geq x^2 - 2x + 1 \\
0 &\geq x^2 - 4x
\end{align*}
Or $x^2 - 4x = x(x-4) \leq 0$ pour $x \in [0\,;\,4]$. Avec la condition $x \geq 1$, on obtient $x \in [1\,;\,4]$.

\textbf{Réunion des deux cas :} $S = \left[-\dfrac{1}{2}\,;\,4\right]$.
\end{exemplebox}

\begin{attention}[Cas 1 et cas 2 : réunion, pas intersection]
Pour le type $\sqrt{A} \geq B$, les deux cas ($B<0$ et $B \geq 0$) se \textbf{réunissent} : la solution finale est l'\emph{union} des deux ensembles. Pour le type $\sqrt{A} \leq B$, seul le cas $B \geq 0$ fournit des solutions, et on \emph{intersecte} les conditions du système. Bien distinguer les deux situations.
\end{attention}

\courssub{Tableau récapitulatif des quatre cas}

\begin{center}
\rowcolors{1}{popcream}{white}
\begin{tabularx}{\linewidth}{|l|X|X|}
\hline
\rowcolor{popblueL} \textbf{Inéquation} & \textbf{Si $B < 0$} & \textbf{Si $B \geq 0$} \\
\hline
$\sqrt{A} \leq B$ & Pas de solution & $\begin{cases} A \geq 0 \\ A \leq B^2 \end{cases}$ \\
\hline
$\sqrt{A} \geq B$ & $A \geq 0$ & $A \geq B^2$ \\
\hline
\end{tabularx}
\end{center}

\begin{popfigure}
\begin{tikzpicture}[scale=1.05, every node/.style={font=\small}]
% Noeud racine
\node[draw=popdark, fill=popblueL, rounded corners, thick, align=center] (racine) at (0,0) {Inéquation\\ irrationnelle};
% Branche gauche
\node[draw=popblue, fill=popblueL!50, rounded corners, thick, align=center] (g1) at (-4.2,-1.8) {$\sqrt{A} \leq B$};
\node[draw=popmuted, rounded corners, align=center] (g2) at (-6.2,-3.6) {$B<0$ :\\ $S=\emptyset$};
\node[draw=popblue, rounded corners, align=center] (g3) at (-2.4,-3.6) {$B\geq 0$ :\\ $A\geq 0$ et $A\leq B^2$};
% Branche droite
\node[draw=poporange, fill=poporangeL, rounded corners, thick, align=center] (d1) at (4.2,-1.8) {$\sqrt{A} \geq B$};
\node[draw=popmuted, rounded corners, align=center] (d2) at (2.4,-3.6) {$B<0$ :\\ $A \geq 0$};
\node[draw=poporange, rounded corners, align=center] (d3) at (6.2,-3.6) {$B\geq 0$ :\\ $A \geq B^2$};
% Arêtes
\draw[thick, popblue, -{Stealth}] (racine) -- (g1);
\draw[thick, poporange, -{Stealth}] (racine) -- (d1);
\draw[popmuted, -{Stealth}] (g1) -- (g2);
\draw[popblue, -{Stealth}] (g1) -- (g3);
\draw[popmuted, -{Stealth}] (d1) -- (d2);
\draw[poporange, -{Stealth}] (d1) -- (d3);
\end{tikzpicture}
\legende{Schéma d'équivalence : la discussion porte toujours sur le signe du membre de droite $B$.}
\end{popfigure}

\courssub{Inéquations de la forme $\sqrt{A} \leq \sqrt{B}$}

Lorsque les deux membres sont des racines carrées, les deux sont positifs ou nuls (dès qu'ils existent) et l'élévation au carré conserve l'ordre sans discussion.

\begin{propriete}[Inéquation $\sqrt{A} \leq \sqrt{B}$]
\[ \sqrt{A} \leq \sqrt{B} \iff 0 \leq A \leq B \]
Autrement dit, on résout le système $\begin{cases} A \geq 0 \\ A \leq B \end{cases}$ (la condition $B \geq 0$ est alors automatique).
\end{propriete}

\begin{exemplebox}[Résolution de $\sqrt{x+2} \leq \sqrt{3x-2}$]
Le système équivalent est :
\[ \begin{cases} x + 2 \geq 0 \\ x + 2 \leq 3x - 2 \end{cases} \iff \begin{cases} x \geq -2 \\ 4 \leq 2x \end{cases} \iff \begin{cases} x \geq -2 \\ x \geq 2 \end{cases} \]
L'intersection donne $S = [2\,;\,+\infty[$.
\end{exemplebox}

\courssub{Méthode générale et lecture graphique}

\begin{methode}[Résoudre une inéquation irrationnelle en 5 étapes]
\begin{enumerate}
\item \textbf{Domaine} : écrire la condition d'existence $A \geq 0$.
\item \textbf{Discussion} : déterminer le signe du membre de droite $B$ (résoudre $B < 0$ et $B \geq 0$).
\item \textbf{Système équivalent} : appliquer le cas du tableau récapitulatif adapté à la forme de l'inéquation.
\item \textbf{Résolution} : résoudre les équations et inéquations polynomiales obtenues (second degré, factorisation, tableau de signes).
\item \textbf{Intersection ou réunion} : croiser avec le domaine et les conditions de l'étape 2, puis conclure par un ensemble de solutions clairement écrit (intervalle ou réunion d'intervalles).
\end{enumerate}
\end{methode}

Illustrons graphiquement l'inéquation $\sqrt{x} \leq x - 2$. On trace la courbe de $y = \sqrt{x}$ et la droite $y = x - 2$. Les solutions sont les abscisses des points où la courbe est \emph{sous} la droite. La droite coupe l'axe des abscisses en $x = 2$ ; en résolvant $\sqrt{x} = x - 2$ (avec $x \geq 2$), on trouve $x = 4$ (la valeur $x = 1$, obtenue par élévation au carré, est rejetée car $1 < 2$). La zone solution est donc $[4\,;\,+\infty[$.

\begin{popfigure}
\begin{center}
\begin{tikzpicture}
\begin{axis}[
    width=0.85\linewidth, height=7cm,
    axis lines=middle,
    xlabel={$x$}, ylabel={$y$},
    xmin=-0.5, xmax=8, ymin=-1, ymax=3.5,
    xtick={1,2,3,4,5,6,7},
    ytick={1,2,3},
    grid=both, grid style={popline!40},
    legend style={at={(0.03,0.97)}, anchor=north west, font=\small}
]
% Courbe racine
\addplot[popcurve, popblue, domain=0:8, samples=200] {sqrt(max(0,x))};
\addlegendentry{$y=\sqrt{x}$}
% Droite
\addplot[popcurve, poporange, domain=-0.5:5.2, samples=2] {x-2};
\addlegendentry{$y=x-2$}
% Point d'intersection
\addplot[only marks, mark=*, popdark, mark size=2pt] coordinates {(4,2)};
\node[popdark, font=\small, anchor=south west] at (axis cs:4,2) {$(4\,;\,2)$};
% Axe solution
\draw[popgreen, line width=2.5pt] (axis cs:4,0) -- (axis cs:8,0);
\node[popgreen, font=\small, anchor=south] at (axis cs:6,-0.3) {$S=[4\,;\,+\infty[$};
\end{axis}
\end{tikzpicture}
\end{center}
\legende{L'inéquation $\sqrt{x} \leq x-2$ : la courbe bleue passe sous la droite orange à partir de $x=4$. Le segment vert sur l'axe des abscisses matérialise la solution $S=[4\,;\,+\infty[$.}
\end{popfigure}

\begin{exoresolu}[Inéquation complète : $\sqrt{x+3} \leq x+1$]
Énoncé. Résoudre dans $\mathbb{R}$ l'inéquation $\sqrt{x+3} \leq x+1$ et représenter les solutions sur une droite graduée.
\tcblower
\textbf{Solution détaillée.}

\textbf{1. Domaine.} $\sqrt{x+3}$ existe si et seulement si $x + 3 \geq 0$, soit $x \geq -3$.

\textbf{2. Discussion sur $B = x+1$.} Si $x + 1 < 0$ (c'est-à-dire $x < -1$), le membre de droite est négatif alors que le membre de gauche est positif ou nul : \textbf{pas de solution} dans ce cas.

\textbf{3. Cas $x \geq -1$.} Les deux membres sont positifs ou nuls, on élève au carré :
\[ x + 3 \leq (x+1)^2 \iff x^2 + x - 2 \geq 0. \]
$\Delta = 9$, racines $-2$ et $1$ ; le trinôme est positif à l'extérieur des racines : $x \in \left]-\infty\,;\,-2\right] \cup [1\,;\,+\infty[$.

\textbf{4. Intersection.} Avec $x \geq -3$ et $x \geq -1$ : seul $[1\,;\,+\infty[$ subsiste.

\textbf{5. Conclusion.} $\boxed{S = [1\,;\,+\infty[}$. Vérification : pour $x = 1$, $\sqrt{4} = 2 = 1+1$ ; pour $x = 0$, $\sqrt{3} \approx \num{1,73} > 1$, donc $0 \notin S$, cohérent.
\end{exoresolu}

\courssub{Applications : modéliser des situations concrètes}

Au Probatoire technique, les inéquations irrationnelles apparaissent souvent dans des problèmes de \cle{dimensions}, de \cle{distances} ou de \cle{coûts}. La démarche complète — mise en inéquation, résolution, interprétation — doit être rédigée et justifiée.

\begin{exoresolu}[Problème de chantier]
Énoncé. Sur un chantier à Douala, une échelle de longueur $\sqrt{x+5}$ mètres (où $x$ est un réglage en mètres) doit permettre d'atteindre une hauteur d'au plus $x - 1$ mètres pour des raisons de sécurité. Déterminer les valeurs de $x$ convenables.
\tcblower
\textbf{Mise en inéquation.} La contrainte s'écrit $\sqrt{x+5} \leq x - 1$.

\textbf{Domaine :} $x \geq -5$. \textbf{Signe de $B$ :} si $x < 1$, pas de solution.

\textbf{Cas $x \geq 1$ :} on élève au carré :
\[ x + 5 \leq (x-1)^2 \iff x^2 - 3x - 4 \geq 0. \]
$\Delta = 9 + 16 = 25$, racines $x_1 = -1$ et $x_2 = 4$. Le trinôme est positif pour $x \leq -1$ ou $x \geq 4$.

\textbf{Intersection} avec $x \geq 1$ : $x \geq 4$.

\textbf{Conclusion interprétée.} Le réglage doit être d'au moins $4$ mètres : $S = [4\,;\,+\infty[$. Pour $x = 4$ : échelle de $\sqrt{9} = 3$ m pour une hauteur de $3$ m — la contrainte est juste saturée.
\end{exoresolu}

\begin{exercice}[À faire]
\begin{enumerate}
\item Résoudre dans $\mathbb{R}$ : $\sqrt{3x+4} \leq x$ ; $\sqrt{x^2 - 3} \geq x - 2$ ; $\sqrt{2x-1} \leq \sqrt{x+4}$.
\item Un atelier découpe des plaques carrées dont la diagonale vaut $\sqrt{2x+8}$ cm. Pour passer dans la machine, la diagonale ne doit pas dépasser $x$ cm. Quelles valeurs de $x$ conviennent ?
\end{enumerate}
\end{exercice}

\begin{corrige}[Exercice 1 (question 1, première inéquation)]
Domaine : $x \geq -\dfrac{4}{3}$. Si $x < 0$, pas de solution (membre de droite négatif). Si $x \geq 0$ : $3x + 4 \leq x^2 \iff x^2 - 3x - 4 \geq 0$, racines $-1$ et $4$, positif à l'extérieur. Intersection avec $x \geq 0$ : $S = [4\,;\,+\infty[$.
\end{corrige}

\begin{savaistu}[Et en Terminale technique ?]
Les techniques de cette section seront réutilisées en classe de Terminale pour étudier les \cle{fonctions racine carrée} (domaine, sens de variation, courbe) et résoudre des problèmes d'\cle{optimisation} : minimiser une longueur de câble, maximiser une aire sous contrainte, dimensionner une pièce mécanique. La rigueur acquise ici — conditions d'existence, discussion sur les signes, vérification — est exactement celle attendue à l'épreuve de mathématiques du Baccalauréat technique.
\end{savaistu}

\begin{aretenir}[L'essentiel de la section]
\begin{itemize}
\item Avant toute résolution : \textbf{domaine} ($A \geq 0$) et \textbf{signe de $B$}.
\item $\sqrt{A} \leq B$ : impossible si $B<0$ ; sinon $\begin{cases} A \geq 0 \\ A \leq B^2 \end{cases}$.
\item $\sqrt{A} \geq B$ : si $B<0$, il suffit que $A \geq 0$ ; sinon $A \geq B^2$.
\item $\sqrt{A} \leq \sqrt{B} \iff 0 \leq A \leq B$.
\item On n'élève au carré une inégalité que lorsque \textbf{les deux membres sont positifs ou nuls}.
\item La réponse finale est toujours un ensemble d'intervalles, obtenu par intersection (ou réunion des cas), puis interprété dans le contexte du problème.
\end{itemize}
\end{aretenir}

\newpage

\coursec{Activité d'intégration}

\coursec{Équations, inéquations polynômiales et irrationnelles}

\courssub{Objectifs de l'activité}

À la fin de cette activité, tu seras capable de :
\begin{itemize}
\item traduire une situation concrète de dimensionnement par un \cle{système somme-produit} et le résoudre en se ramenant à une équation du second degré ;
\item factoriser un \cle{polynôme de degré 3} en utilisant une racine évidente ou un facteur commun, puis étudier son \cle{signe} à l'aide d'un tableau de signes ;
\item résoudre une \cle{équation irrationnelle} en vérifiant les conditions d'existence et en contrôlant les solutions obtenues ;
\item résoudre une \cle{inéquation irrationnelle} simple en utilisant l'équivalence $\sqrt{A} \geq B \Longleftrightarrow \begin{cases} A \geq 0 \\ B \geq 0 \text{ et } A \geq B^2 \end{cases}$ ou $\begin{cases} A \geq 0 \\ B < 0 \end{cases}$ ;
\item \cle{rédiger et justifier} une démarche complète, une réponse et une conclusion, comme à l'épreuve du Probatoire technique.
\end{itemize}

\courssub{Situation}

Maman Béa possède à Bafoussam, dans la région de l'Ouest, un champ rectangulaire qu'elle veut aménager pour y cultiver du maïs et des haricots. Le chef du quartier lui a remis un plan sommaire : le champ a un \textbf{périmètre de 100 m} et une \textbf{aire de 600 m$^2$}, mais le plan ne mentionne ni la longueur ni la largeur. Maman Béa doit connaître ces dimensions pour acheter le fil de clôture au marché et commander les plants.

En plus, elle souhaite aménager le champ de la manière suivante :
\begin{itemize}
\item creuser un \textbf{bassin cubique} de côté variable $x$ (en mètres) pour stocker l'eau de pluie ; le volume d'eau disponible, exprimé en m$^3$, est modélisé par l'expression $V(x) = x^3 - 6x^2 + 9x$. Le bassin ne sert à rien si le volume est nul, et il n'est rentable que si le volume est strictement positif ;
\item installer un \textbf{tuyau d'arrosage} reliant le bassin au coin opposé du champ. D'après les calculs du technicien de la coopérative agricole, la longueur $L$ du tuyau (en mètres) doit vérifier la relation $\sqrt{2L+1} = L - 1$.
\end{itemize}

Maman Béa te sollicite, toi l'élève de Première technique, pour l'aider à mener tous les calculs et à rédiger un rapport clair qu'elle présentera à la coopérative.

\courssub{Consignes}

\begin{enumerate}
\item \textbf{Dimensions du champ.} On note $x$ et $y$ les dimensions (en mètres) du champ rectangulaire.
\begin{enumerate}
\item Traduire les données de l'énoncé par un système de deux équations à deux inconnues faisant intervenir la somme et le produit de $x$ et $y$.
\item Montrer que $x$ et $y$ sont solutions d'une équation du second degré que l'on déterminera.
\item Résoudre cette équation et conclure sur les dimensions du champ.
\end{enumerate}
\item \textbf{Le bassin.} On considère le polynôme $V(x) = x^3 - 6x^2 + 9x$.
\begin{enumerate}
\item Factoriser complètement $V(x)$ (on commencera par mettre $x$ en facteur).
\item Déterminer les valeurs de $x$ pour lesquelles le volume est nul.
\item Dresser le tableau de signes de $V(x)$ sur $\mathbb{R}$ et en déduire l'ensemble des valeurs de $x$ pour lesquelles le volume est strictement positif. Interpréter le résultat pour Maman Béa.
\end{enumerate}
\item \textbf{Le tuyau d'arrosage.} Résoudre dans $\mathbb{R}$ l'équation $\sqrt{2L+1} = L - 1$, en précisant les conditions d'existence et en vérifiant chaque solution trouvée. Donner la longueur du tuyau à commander.
\item \textbf{Complément : inéquation irrationnelle.} On souhaite savoir pour quelles longueurs $L$ le tuyau est au moins aussi long que la racine carrée, c'est-à-dire résoudre $\sqrt{2L+1} \leq L - 1$. Résoudre cette inéquation en $\mathbb{R}$.
\item \textbf{Rédaction finale.} Rédiger une conclusion complète et justifiée pour chaque question, sous la forme d'un court rapport destiné à la coopérative, comme tu le ferais à l'épreuve du Probatoire technique.
\end{enumerate}

\courssub{Ressources à mobiliser}

\begin{aretenir}[Boîte à outils de la leçon]
\begin{itemize}
\item \textbf{Système somme-produit :} deux nombres de somme $S$ et de produit $P$ sont les solutions de l'équation $X^2 - SX + P = 0$.
\item \textbf{Équation du second degré :} pour $aX^2 + bX + c = 0$ ($a \neq 0$), on calcule le discriminant $\Delta = b^2 - 4ac$. Si $\Delta > 0$, deux solutions : $X_1 = \dfrac{-b - \sqrt{\Delta}}{2a}$ et $X_2 = \dfrac{-b + \sqrt{\Delta}}{2a}$.
\item \textbf{Factorisation d'un polynôme de degré 3 :} chercher un facteur commun ou une racine évidente $a$ (souvent $0$, $1$, $-1$, $2$) ; si $P(a) = 0$, alors $P(x)$ est factorisable par $(x - a)$.
\item \textbf{Signe d'un polynôme :} un polynôme factorisé a le signe du produit de ses facteurs ; on dresse un tableau de signes.
\item \textbf{Équation irrationnelle $\sqrt{A} = B$ :} elle exige $A \geq 0$ \emph{et} $B \geq 0$ ; on élève alors au carré : $A = B^2$, puis on \textbf{vérifie} chaque solution dans l'équation de départ.
\item \textbf{Inéquation irrationnelle $\sqrt{A} \leq B$ (ou $\geq$) :}
      \begin{itemize}
      \item Si $B < 0$ : $\sqrt{A} \leq B$ impossible (racine $\geq 0$) ; $\sqrt{A} \geq B$ toujours vrai si $A \geq 0$.
      \item Si $B \geq 0$ : équivalent à $\begin{cases} A \geq 0 \\ A \leq B^2 \end{cases}$ (pour $\leq$) ou $\begin{cases} A \geq 0 \\ A \geq B^2 \end{cases}$ (pour $\geq$).
      \end{itemize}
\end{itemize}
\end{aretenir}

\begin{attention}[Pièges fréquents]
\begin{itemize}
\item Dans une équation irrationnelle, élever au carré peut créer des \textbf{solutions étrangères} : la vérification finale est obligatoire, jamais facultative.
\item Ne pas oublier que le membre de droite de $\sqrt{A} = B$ doit être positif ou nul : une racine carrée est toujours \textbf{positive ou nulle}.
\item Dans le système somme-produit, le périmètre $100$ m donne $2(x+y) = 100$, donc la somme est $S = 50$ et non $100$.
\item Pour une inéquation irrationnelle, ne pas élever au carré sans discuter le signe du second membre : l'équivalence n'est valide que si les deux membres sont positifs.
\end{itemize}
\end{attention}

\courssub{Démarche et corrigé commenté}

\begin{exoresolu}[Le champ de Maman Béa : corrigé complet]
\textbf{1) Dimensions du champ.} Le champ a un périmètre de $100$ m et une aire de $600$ m$^2$. On pose $x$ et $y$ ses dimensions en mètres. Déterminer $x$ et $y$.
\tcblower
\textbf{Traduction.} Le périmètre d'un rectangle est $2(x+y)$ et son aire est $xy$. Le système s'écrit :
\[
\begin{cases} 2(x+y) = 100 \\ xy = 600 \end{cases}
\quad \Longleftrightarrow \quad
\begin{cases} x + y = 50 \\ xy = 600 \end{cases}
\]
C'est un système \cle{somme-produit} avec $S = 50$ et $P = 600$.

\textbf{Équation du second degré.} D'après la propriété somme-produit, $x$ et $y$ sont les solutions de :
\[ X^2 - SX + P = 0 \quad \Longleftrightarrow \quad X^2 - 50X + 600 = 0. \]

\textbf{Résolution.} Calculons le discriminant :
\[ \Delta = (-50)^2 - 4 \times 1 \times 600 = 2500 - 2400 = 100 > 0. \]
L'équation admet deux solutions :
\[ X_1 = \frac{50 - \sqrt{100}}{2} = \frac{50 - 10}{2} = 20 \quad ; \quad X_2 = \frac{50 + 10}{2} = \frac{50 + 10}{2} = 30. \]

\textbf{Conclusion.} Le champ de Maman Béa mesure \textbf{30 m de long} et \textbf{20 m de large}. Vérification : $2(30+20) = 100$ m et $30 \times 20 = 600$ m$^2$. Le résultat est conforme au plan du chef de quartier.
\end{exoresolu}

\begin{exoresolu}[Le bassin : volume nul et volume positif]
\textbf{2)} On considère $V(x) = x^3 - 6x^2 + 9x$. Factoriser $V(x)$, déterminer les valeurs annulant le volume, puis résoudre $V(x) > 0$.
\tcblower
\textbf{Factorisation.} On met d'abord $x$ en facteur (racine évidente $0$) :
\[ V(x) = x\left(x^2 - 6x + 9\right). \]
On reconnaît une identité remarquable : $x^2 - 6x + 9 = (x-3)^2$. D'où la factorisation complète :
\[ V(x) = x(x-3)^2. \]

\textbf{Volume nul.} $V(x) = 0 \Longleftrightarrow x = 0$ ou $(x-3)^2 = 0 \Longleftrightarrow x = 0$ ou $x = 3$. Le volume est nul pour $x = 0$ et pour $x = 3$ ($3$ est une racine double).

\textbf{Tableau de signes.} Le carré $(x-3)^2$ est toujours positif ou nul ; le signe de $V(x)$ est donc celui de $x$, sauf en $x = 3$ où $V$ s'annule.
\begin{center}
\begin{tabular}{|c|cccccccc|}
\hline
$x$ & $-\infty$ & & $0$ & & $3$ & & $+\infty$ \\
\hline
$x$ & $-$ & $-$ & $0$ & $+$ & $+$ & $+$ & $+$ \\
\hline
$(x-3)^2$ & $+$ & $+$ & $+$ & $+$ & $0$ & $+$ & $+$ \\
\hline
$V(x)$ & $-$ & $-$ & $0$ & $+$ & $0$ & $+$ & $+$ \\
\hline
\end{tabular}
\end{center}

\textbf{Conclusion.} $V(x) > 0 \Longleftrightarrow x \in ]0\,;\,3[ \cup ]3\,;\,+\infty[$. Comme $x$ est une longueur, on retient $x > 0$ avec $x \neq 3$. Pour Maman Béa : le bassin est rentable pour tout côté strictement positif, \textbf{sauf} pour $x = 3$ m où le volume calculé s'annule ; il faut donc éviter la dimension de 3 m.
\end{exoresolu}

\begin{exoresolu}[Le tuyau d'arrosage : équation irrationnelle]
\textbf{3)} Résoudre dans $\mathbb{R}$ l'équation $\sqrt{2L+1} = L - 1$.
\tcblower
\textbf{Conditions d'existence.} Il faut :
\begin{itemize}
\item $2L + 1 \geq 0$, c'est-à-dire $L \geq -\dfrac{1}{2}$ (la racine carrée existe) ;
\item $L - 1 \geq 0$, c'est-à-dire $L \geq 1$ (une racine carrée est positive ou nulle, donc le membre de droite doit être positif ou nul).
\end{itemize}
La condition la plus restrictive est $L \geq 1$.

\textbf{Résolution.} Pour $L \geq 1$, on élève au carré :
\begin{align*}
2L + 1 &= (L-1)^2 \\
2L + 1 &= L^2 - 2L + 1 \\
0 &= L^2 - 4L \\
0 &= L(L - 4).
\end{align*}
D'où $L = 0$ ou $L = 4$.

\textbf{Vérification (étape obligatoire).}
\begin{itemize}
\item $L = 0$ : ne vérifie pas la condition $L \geq 1$. De plus $\sqrt{2(0)+1} = 1$ et $0 - 1 = -1$ : $1 \neq -1$. C'est une \textbf{solution étrangère}, à rejeter.
\item $L = 4$ : $\sqrt{2(4)+1} = \sqrt{9} = 3$ et $4 - 1 = 3$. L'égalité est vérifiée.
\end{itemize}

\textbf{Conclusion.} L'équation admet une unique solution : $L = 4$. Maman Béa doit commander un tuyau d'arrosage de \textbf{4 mètres}.
\end{exoresolu}

\begin{exoresolu}[Complément : inéquation irrationnelle]
\textbf{4)} Résoudre dans $\mathbb{R}$ l'inéquation $\sqrt{2L+1} \leq L - 1$.
\tcblower
\textbf{Conditions d'existence :} $2L+1 \geq 0 \iff L \geq -0,5$.

\textbf{Résolution par discussion du signe du second membre.}
\begin{itemize}
\item \textbf{Cas 1 : $L - 1 < 0$ (c'est-à-dire $L < 1$).} \\
Le second membre est négatif. La racine carrée (positive ou nulle) ne peut pas être $\leq$ à un nombre négatif. \textbf{Aucune solution} dans ce cas.
\item \textbf{Cas 2 : $L - 1 \geq 0$ (c'est-à-dire $L \geq 1$).} \\
Les deux membres sont positifs ou nuls, on peut élever au carré en conservant l'ordre :
\begin{align*}
2L + 1 &\leq (L-1)^2 \\
2L + 1 &\leq L^2 - 2L + 1 \\
0 &\leq L^2 - 4L \\
0 &\leq L(L - 4).
\end{align*}
Le trinôme $L(L-4)$ est positif ou nul pour $L \leq 0$ ou $L \geq 4$.
En croisant avec la condition du cas ($L \geq 1$), on obtient $L \geq 4$.
\end{itemize}

\textbf{Conclusion.} L'ensemble des solutions est $S = [4\,;\,+\infty[$. Le tuyau doit mesurer \textbf{au moins 4 mètres}.
\end{exoresolu}

\begin{exoresolu}[Rédaction du rapport final (modèle de rédaction type Probatoire)]
\textbf{5)} Rédiger la conclusion complète destinée à la coopérative.
\tcblower
\textbf{Rapport technique pour la coopérative agricole de Bafoussam.}

\textbf{Objet :} Dimensionnement du champ, du bassin de rétention et du tuyau d'arrosage.

\textbf{1. Dimensions du champ.} \\
À partir des données du plan (périmètre 100 m, aire 600 m$^2$), nous avons traduit le problème par le système somme-produit $x + y = 50$, $xy = 600$, qui se ramène à l'équation $X^2 - 50X + 600 = 0$. Le discriminant $\Delta = 100$ étant strictement positif, nous obtenons deux dimensions : le champ mesure \textbf{30 m sur 20 m}.

\textbf{2. Bassin de rétention.} \\
La factorisation $V(x) = x(x-3)^2$ et l'étude de signes montrent que le volume est nul en $x = 0$ et $x = 3$, et strictement positif pour tout $x > 0$ différent de $3$ : toute dimension convient donc, \textbf{sauf 3 m} (volume nul). Maman Béa peut choisir toute cote positive différente de 3 m pour le côté du bassin cubique.

\textbf{3. Tuyau d'arrosage.} \\
L'équation irrationnelle $\sqrt{2L+1} = L - 1$, résolue sous la condition $L \geq 1$ avec vérification des candidats, fournit une unique solution $L = 4$ m ; la valeur $L = 0$ a été rejetée comme solution étrangère. \\
Par ailleurs, si l'on exige que le tuyau soit au moins aussi long que la distance calculée (inéquation $\sqrt{2L+1} \leq L - 1$), la résolution montre qu'il faut commander un tuyau de longueur \textbf{supérieure ou égale à 4 m}.

\textbf{4. Synthèse.} \\
En résumé : champ de $30 \times 20$ m, bassin de côté quelconque $>0$ et $\neq 3$ m, tuyau de 4 m (longueur exacte) ou $\geq 4$ m (longueur minimale). Toutes les réponses ont été vérifiées et sont cohérentes avec la réalité du terrain.

\textit{Fait à Bafoussam, le \dots \hfill Signature de l'élève}
\end{exoresolu}

\courssub{Bilan de l'activité}

Cette activité t'a permis de mettre en évidence les points suivants :
\begin{itemize}
\item un problème concret de dimensionnement (périmètre et aire connus) se traduit naturellement par un \cle{système somme-produit}, qui se résout grâce à une équation du second degré ;
\item la \cle{factorisation} d'un polynôme de degré 3 (facteur commun, racine évidente, identité remarquable) est la clé pour trouver ses racines et dresser son \cle{tableau de signes}, ce qui permet de répondre à des questions de rentabilité ou de faisabilité (inéquation polynomiale) ;
\item la résolution d'une \cle{équation irrationnelle} impose une démarche rigoureuse en trois temps : conditions d'existence, élévation au carré, puis \textbf{vérification} des solutions pour éliminer les solutions étrangères ;
\item la résolution d'une \cle{inéquation irrationnelle} impose une discussion sur le signe du second membre avant d'élever au carré, afin de préserver l'équivalence ;
\item enfin, la rédaction d'une conclusion justifiée, étape par étape, est une compétence évaluée au Probatoire technique : chaque réponse doit être annoncée, calculée, vérifiée et interprétée dans le contexte.
\end{itemize}

\begin{savaistu}[Et dans la vraie vie ?]
À Bafoussam comme dans toutes les zones agricoles du Cameroun, les coopératives calculent régulièrement surfaces, volumes de stockage et longueurs de canalisations. Les mêmes outils — équations du second degré, factorisation, équations et inéquations irrationnelles — servent aussi au maçon pour doser un chantier, à l'électricien pour dimensionner un câble (loi de Joule, chute de tension) et au mécanicien pour ajuster une pièce (tolérances, jeux). Les mathématiques de Première technique sont déjà des mathématiques de métier.
\end{savaistu}

\newpage

\coursec{Méthodes et exercices résolus}

\coursec{Équations, inéquations polynômiales et irrationnelles}

\coursec{Rappels : équations et inéquations du second degré}

\begin{propriete}[Discriminant et racines d'un trinôme]
Soit $P(x)=ax^2+bx+c$ avec $a\neq 0$ et $\Delta=b^2-4ac$ son \cle{discriminant}.
\begin{itemize}
\item Si $\Delta>0$ : $P(x)=0$ admet deux racines distinctes $x_1=\dfrac{-b-\sqrt{\Delta}}{2a}$ et $x_2=\dfrac{-b+\sqrt{\Delta}}{2a}$.
\item Si $\Delta=0$ : $P(x)=0$ admet une racine double $x_0=\dfrac{-b}{2a}$.
\item Si $\Delta<0$ : $P(x)=0$ n'admet aucune racine réelle.
\end{itemize}
\emph{Relations racines-coefficients :} $x_1+x_2=-\dfrac{b}{a}$ ; $x_1x_2=\dfrac{c}{a}$.
\emph{Forme factorisée :} $P(x)=a(x-x_1)(x-x_2)$ (si $\Delta\geq 0$).
\end{propriete}

\begin{methode}[Résoudre une équation du second degré dans $\mathbb{R}$]
\begin{enumerate}
\item Identifier $a$, $b$, $c$ et calculer $\Delta=b^2-4ac$.
\item \textbf{Si $\Delta>0$ :} $x_1=\dfrac{-b-\sqrt{\Delta}}{2a}$ ; $x_2=\dfrac{-b+\sqrt{\Delta}}{2a}$.
\item \textbf{Si $\Delta=0$ :} $x_0=\dfrac{-b}{2a}$.
\item \textbf{Si $\Delta<0$ :} $S=\varnothing$.
\item Conclure par une phrase : « L'ensemble des solutions est $S=\{\ldots\}$ » ou « $S=\varnothing$ ».
\end{enumerate}
\end{methode}

\begin{exoresolu}[Dalle de chantier -- Problème concret]
Un maçon coule une dalle rectangulaire. La longueur dépasse la largeur de \SI{2}{m} et l'aire est \SI{15}{m^2}. Déterminer les dimensions.
\tcblower
\textbf{Modélisation.} Largeur $x$ (en m), $x>0$. Longueur $x+2$. Aire : $x(x+2)=15$.
\textbf{Équation :} $x^2+2x-15=0$. Coefficients : $a=1$, $b=2$, $c=-15$.
\textbf{Discriminant :} $\Delta=2^2-4\times1\times(-15)=4+60=64>0$.
\textbf{Racines :}
\[x_1=\frac{-2-\sqrt{64}}{2}=\frac{-10}{2}=-5 \quad;\quad x_2=\frac{-2+\sqrt{64}}{2}=\frac{6}{2}=3.\]
\textbf{Interprétation :} Une longueur est positive, on rejette $x_1=-5$.
\textbf{Conclusion :} Largeur $=\SI{3}{m}$, Longueur $=\SI{5}{m}$. Vérification : $3\times5=15$.
\end{exoresolu}

\begin{propriete}[Signe d'un trinôme du second degré]
Soit $P(x)=ax^2+bx+c$ ($a\neq 0$) de racines $x_1\leq x_2$ (réelles ou non).
\begin{itemize}
\item Si $\Delta<0$ : $P(x)$ a le \textbf{signe de $a$} pour tout $x\in\mathbb{R}$.
\item Si $\Delta\geq 0$ : $P(x)$ a le \textbf{signe de $a$} à l'\emph{extérieur} de l'intervalle $[x_1\,;\,x_2]$ et le \textbf{signe contraire} à l'\emph{intérieur} (sur $]x_1\,;\,x_2[$).
\end{itemize}
\end{propriete}

\begin{methode}[Résoudre une inéquation du second degré]
Pour $ax^2+bx+c \geq 0$ (ou $\leq 0$, $>0$, $<0$) :
\begin{enumerate}
\item Résoudre l'équation associée $ax^2+bx+c=0$ (calculer $\Delta$, racines $x_1\leq x_2$).
\item Construire le \cle{tableau de signes} : une ligne pour $x$, une ligne pour le signe de $P(x)$.
\item Lire l'ensemble solution $S$ sur la ligne du signe, en respectant strict/large (crochets $]$ $[$ ou $[$ $]$).
\end{enumerate}
\end{methode}

\begin{exoresolu}[Inéquation du second degré]
Résoudre dans $\mathbb{R}$ : $-x^2+3x+4\geq 0$.
\tcblower
\textbf{1. Équation associée :} $-x^2+3x+4=0$. $a=-1$, $b=3$, $c=4$.
$\Delta=3^2-4(-1)(4)=9+16=25>0$.
$x_1=\dfrac{-3-5}{-2}=4$ ; $x_2=\dfrac{-3+5}{-2}=-1$. (Ordre : $-1 < 4$).
\textbf{2. Tableau de signes.} $a=-1<0$ $\Rightarrow$ signe $-$ à l'extérieur, $+$ entre les racines.
\begin{center}
\begin{tabular}{|c|ccccccc|}\hline
$x$ & $-\infty$ & & $-1$ & & $4$ & & $+\infty$ \\ \hline
$-x^2+3x+4$ & $-$ & $-$ & $0$ & $+$ & $+$ & $0$ & $-$ \\ \hline
\end{tabular}
\end{center}
\textbf{3. Conclusion.} On cherche $\geq 0$ (signe $+$ ou $0$), inégalité large $\Rightarrow$ bornes incluses.
\[S=[-1\,;\,4].\]
\end{exoresolu}

\begin{aretenir}[Rappels clés -- Second degré]
\begin{itemize}
\item $\Delta=b^2-4ac$ décide du nombre de racines.
\item Signe du trinôme = signe de $a$ \emph{hors} des racines, signe contraire \emph{entre} les racines.
\item Tableau de signes : valeurs racines $\to$ cases $0$ ; intervalles $\to$ signes $\pm$.
\item Inégalité large ($\geq,\leq$) $\to$ crochets fermés $[$ $]$ ; stricte ($>$,$<$) $\to$ crochets ouverts $]$ $[$.
\end{itemize}
\end{aretenir}

\coursec{Systèmes se ramenant à une équation du second degré}

\begin{propriete}[Système somme-produit]
Le système $\begin{cases} x+y=S \\ xy=P \end{cases}$ est équivalent à : $x$ et $y$ sont les deux racines (réelles) de l'équation $X^2-SX+P=0$.
Condition d'existence de solutions réelles : $\Delta' = S^2-4P \geq 0$.
Les solutions du système sont les \textbf{couples} $(x_1;y_1)$ et $(x_2;y_2)$ (ordre interchangeable).
\end{propriete}

\begin{methode}[Résolution par substitution (cas général)]
\begin{enumerate}
\item Exprimer une inconnue en fonction de l'autre grâce à l'équation la plus simple (souvent du 1er degré).
\item Substituer dans l'autre équation $\to$ équation du second degré en une inconnue.
\item Résoudre cette équation (discriminant).
\item Remonter à la seconde inconnue pour former les \textbf{couples solutions} $(x;y)$.
\item Vérifier les couples dans le système initial.
\end{enumerate}
\end{methode}

\begin{exoresolu}[Dimensions d'un atelier]
Un atelier rectangulaire a un périmètre de \SI{26}{m} et une aire de \SI{40}{m^2}. Trouver ses dimensions.
\tcblower
\textbf{Variables :} $x$, $y$ dimensions en mètres ($x>0, y>0$).
\textbf{Système :} $\begin{cases} 2(x+y)=26 \iff x+y=13 \\ xy=40 \end{cases}$
\textbf{Somme-Produit :} $S=13$, $P=40$. Équation : $X^2-13X+40=0$.
$\Delta'=(-13)^2-4\times40=169-160=9>0$.
$X_1=\dfrac{13-3}{2}=5$ ; $X_2=\dfrac{13+3}{2}=8$.
\textbf{Couples :} $(5;8)$ et $(8;5)$ (mêmes dimensions).
\textbf{Conclusion :} L'atelier mesure \SI{5}{m} sur \SI{8}{m}.
\end{exoresolu}

\begin{exercice}[Système symétrique]
Résoudre dans $\mathbb{R}^2$ : $\begin{cases} x+y=7 \\ x^2+y^2=25 \end{cases}$.
\end{exercice}
\begin{corrige}[Exercice : Système symétrique]
On utilise $(x+y)^2=x^2+y^2+2xy \Rightarrow 49=25+2xy \Rightarrow xy=12$.
Système équivalent : $\begin{cases} x+y=7 \\ xy=12 \end{cases}$.
Équation : $X^2-7X+12=0$. $\Delta=49-48=1$.
$X_1=\frac{7-1}{2}=3$ ; $X_2=\frac{7+1}{2}=4$.
Solutions : $(3;4)$ et $(4;3)$.
\end{corrige}

\coursec{Polynômes de degré 3 : racines et factorisation}

\begin{propriete}[Théorème du facteur -- Racine et divisibilité]
Soit $P$ un polynôme de degré 3 à coefficients réels et $a\in\mathbb{R}$.
\[P(a)=0 \quad\Longleftrightarrow\quad (x-a) \text{ divise } P(x).\]
Autrement dit : $a$ est une \cle{racine} de $P$ $\iff$ il existe un polynôme $Q$ de degré 2 tel que $P(x)=(x-a)Q(x)$.
\end{propriete}

\begin{methode}[Factoriser un polynôme de degré 3]
\begin{enumerate}
\item \textbf{Chercher une racine évidente} : tester $0$, $\pm1$, $\pm2$, $\pm3$... (diviseurs du terme constant). Si $P(a)=0$, $a$ est racine.
\item \textbf{Factoriser par $(x-a)$} : écrire $P(x)=(x-a)(ax^2+bx+c)$ et déterminer $a,b,c$ par \cle{identification} (développer et comparer les coefficients) ou par \cle{division euclidienne}.
\item \textbf{Factoriser le trinôme restant} $ax^2+bx+c$ (discriminant, racines) pour obtenir une factorisation complète en facteurs du 1er degré (si $\Delta\geq 0$).
\end{enumerate}
\end{methode}

\begin{experience}[Découverte : Racine et factorisation]
Soit $P(x)=x^3-2x^2-x+2$.
\begin{enumerate}
\item Calculer $P(1)$, $P(-1)$, $P(2)$. Que remarquez-vous ?
\item Pour une racine trouvée $a$, effectuer la division de $P(x)$ par $(x-a)$.
\item Écrire $P(x)$ sous forme factorisée. Combien de racines réelles $P$ admet-il ?
\end{enumerate}
\textit{Indication :} La division euclidienne ou l'identification donnent le quotient de degré 2.
\end{experience}

\begin{exoresolu}[Factorisation complète]
Soit $P(x)=x^3-4x^2+x+6$.
\begin{enumerate}
\item Vérifier que $2$ est une racine de $P$.
\item Factoriser $P(x)$ en produit de facteurs du premier degré.
\item Donner l'ensemble des racines de $P$.
\end{enumerate}
\tcblower
\textbf{1.} $P(2)=8-16+2+6=0$. Donc $2$ est racine, $P(x)$ divisible par $(x-2)$.
\textbf{2.} Identification : $P(x)=(x-2)(ax^2+bx+c)$.
Développement : $ax^3+(b-2a)x^2+(c-2b)x-2c$.
Système : $\begin{cases} a=1 \\ b-2a=-4 \\ c-2b=1 \\ -2c=6 \end{cases} \Rightarrow a=1,\; b=-2,\; c=-3$.
$P(x)=(x-2)(x^2-2x-3)$.
Trinôme : $\Delta=4+12=16$. Racines : $\frac{2\pm4}{2} \Rightarrow -1$ et $3$.
$x^2-2x-3=(x+1)(x-3)$.
\textbf{Factorisation finale :} $\boxed{P(x)=(x-2)(x+1)(x-3)}$.
\textbf{3.} Racines : $-1,\; 2,\; 3$.
\end{exoresolu}

\begin{savaistu}[Racines rationnelles]
Pour un polynôme $P(x)=a_nx^n+\dots+a_0$ à coefficients \textbf{entiers}, toute racine rationnelle $\frac{p}{q}$ (irréductible) vérifie : $p$ divise $a_0$ et $q$ divise $a_n$. En Première technique, on cherche surtout les racines \emph{entières} (diviseurs du terme constant).
\end{savaistu}

\coursec{Signe d'un polynôme de degré 3 et inéquations polynomiales}

\begin{propriete}[Signe d'un produit -- Règle des signes]
Le produit de facteurs du premier degré $(x-\alpha_1)(x-\alpha_2)\dots(x-\alpha_n)$ change de signe à chaque racine simple. Le tableau de signes se construit en croisant les signes de chaque facteur.
\end{propriete}

\begin{methode}[Résoudre $P(x) \geq 0$ (ou $\leq, >, <$) pour $P$ degré 3]
\begin{enumerate}
\item Factoriser $P(x)$ complètement (racine évidente + trinôme).
\item Pour chaque facteur $(x-\alpha)$, déterminer son signe : $>0$ si $x>\alpha$, $<0$ si $x<\alpha$, $=0$ si $x=\alpha$.
\item Dresser le \cle{tableau de signes global} : une ligne par facteur, une ligne produit. Colonnes : valeurs remarquables ($-\infty$, racines, $+\infty$) \textbf{et} intervalles entre elles.
\item Lire les solutions sur la ligne produit selon l'inégalité (strict/large).
\end{enumerate}
\end{methode}

\begin{exoresolu}[Inéquation de degré 3]
Résoudre dans $\mathbb{R}$ : $x^3-4x^2+x+6 > 0$.
\tcblower
\textbf{1. Factorisation (exercice précédent) :} $P(x)=(x+1)(x-2)(x-3)$.
Racines ordonnées : $-1 < 2 < 3$.
\textbf{2. Signes des facteurs :}
$x+1$ : $0$ en $-1$, $-$ avant, $+$ après.
$x-2$ : $0$ en $2$, $-$ avant, $+$ après.
$x-3$ : $0$ en $3$, $-$ avant, $+$ après.
\textbf{3. Tableau de signes global (10 colonnes : 1 var + 5 valeurs + 4 intervalles) :}
\begin{center}
\begin{tabular}{|c|ccccccccccc|}\hline
$x$ & $-\infty$ & & $-1$ & & $2$ & & $3$ & & $+\infty$ \\ \hline
$x+1$ & $-$ & $-$ & $0$ & $+$ & $+$ & $+$ & $+$ & $+$ & $+$ \\
$x-2$ & $-$ & $-$ & $-$ & $-$ & $0$ & $+$ & $+$ & $+$ & $+$ \\
$x-3$ & $-$ & $-$ & $-$ & $-$ & $-$ & $-$ & $0$ & $+$ & $+$ \\ \hline
$P(x)$ & $-$ & $-$ & $0$ & $+$ & $+$ & $0$ & $-$ & $-$ & $0$ & $+$ & $+$ \\ \hline
\end{tabular}
\end{center}
\textit{Note :} La ligne $P(x)$ se lit : signe sur $]-\infty;-1[$ ($-$), valeur en $-1$ ($0$), signe sur $]-1;2[$ ($+$), valeur en $2$ ($0$), signe sur $]2;3[$ ($-$), valeur en $3$ ($0$), signe sur $]3;+\infty[$ ($+$).
\textbf{4. Conclusion.} Inégalité stricte $>$ $\Rightarrow$ racines exclues.
\[S=]-1\,;\,2[\;\cup\;]3\,;\,+\infty[.\]
\end{exoresolu}

\begin{exercice}[Inéquation degré 3 -- Application]
Résoudre : $(x-1)(x^2-4)\leq 0$.
\end{exercice}
\begin{corrige}[Exercice : Inéquation degré 3]
Factorisation : $(x-1)(x-2)(x+2)\leq 0$. Racines : $-2 < 1 < 2$.
Tableau de signes (facteurs $x+2$, $x-1$, $x-2$, coefficient directeur $+$).
Ligne produit : $-$ sur $]-\infty;-2[$, $0$ en $-2$, $+$ sur $]-2;1[$, $0$ en $1$, $-$ sur $]1;2[$, $0$ en $2$, $+$ sur $]2;+\infty[$.
Inégalité large $\leq$ $\Rightarrow$ bornes incluses.
$S=]-\infty;-2]\;\cup\;[1;2]$.
\end{corrige}

\coursec{Équations irrationnelles}

\begin{propriete}[Équivalence fondamentale -- Racine carrée]
Pour $A(x)\geq 0$ et $B(x)\geq 0$ :
\[\sqrt{A(x)}=B(x) \quad\Longleftrightarrow\quad \begin{cases} A(x)\geq 0 \\ B(x)\geq 0 \\ A(x)=[B(x)]^2 \end{cases}\]
Pour $\sqrt{A(x)}=\sqrt{B(x)}$ : $\begin{cases} A(x)\geq 0 \\ B(x)\geq 0 \\ A(x)=B(x) \end{cases}$.
\emph{Rappel :} $\sqrt{u}\geq 0$ par définition. Toute solution ne vérifiant pas les conditions est \textbf{parasite}.
\end{propriete}

\begin{methode}[Résoudre $\sqrt{A(x)}=B(x)$]
\begin{enumerate}
\item Écrire les \textbf{conditions d'existence (C.E.)} : $A(x)\geq 0$ \textbf{et} $B(x)\geq 0$. Déterminer le domaine $\mathcal{D}$.
\item Élever au carré : résoudre $A(x)=[B(x)]^2$ (équation polynômiale degré 1 ou 2).
\item \textbf{Vérifier} chaque solution candidate dans $\mathcal{D}$ (ou dans l'équation initiale). Rejeter les solutions parasites.
\item Écrire l'ensemble solution $S$.
\end{enumerate}
\end{methode}

\begin{exoresolu}[Équation irrationnelle standard]
Résoudre dans $\mathbb{R}$ : $\sqrt{x^2-5}=x-1$.
\tcblower
\textbf{C.E. :} $\begin{cases} x^2-5\geq 0 \iff x\leq -\sqrt{5} \text{ ou } x\geq \sqrt{5} \\ x-1\geq 0 \iff x\geq 1 \end{cases} \Rightarrow \mathcal{D}=[\sqrt{5}\,;\,+\infty[$ (car $\sqrt{5}\approx 2,24 > 1$).
\textbf{Élevation au carré (sur $\mathcal{D}$) :} $x^2-5=(x-1)^2 \iff x^2-5=x^2-2x+1 \iff 2x=6 \iff x=3$.
\textbf{Vérification :} $3\in[\sqrt{5};+\infty[$ ? Oui. $\sqrt{9-5}=2$ et $3-1=2$. OK.
\textbf{Solution :} $S=\{3\}$.
\end{exoresolu}

\begin{exoresolu}[Équation $\sqrt{A}=\sqrt{B}$]
Résoudre : $\sqrt{2x+3}=\sqrt{x+5}$.
\tcblower
\textbf{C.E. :} $\begin{cases} 2x+3\geq 0 \iff x\geq -1,5 \\ x+5\geq 0 \iff x\geq -5 \end{cases} \Rightarrow \mathcal{D}=[-1,5\,;\,+\infty[$.
\textbf{Équation :} $2x+3=x+5 \iff x=2$.
\textbf{Vérification :} $2\in\mathcal{D}$ ? Oui. $\sqrt{7}=\sqrt{7}$.
$S=\{2\}$.
\end{exoresolu}

\begin{attention}[Piège : Solutions parasites]
Ne \textbf{jamais} écrire seulement « on élève au carré : $A=B^2$ » sans les C.E. $A\geq 0, B\geq 0$.
Exemple : $\sqrt{x}=-1$. Élever au carré donne $x=1$. Mais $\sqrt{1}=1\neq -1$. $S=\varnothing$.
La condition $B(x)\geq 0$ est \textbf{indispensable} car $\sqrt{A}\geq 0$ toujours.
\end{attention}

\coursec{Inéquations irrationnelles et applications}

\begin{propriete}[Équivalences pour inéquations irrationnelles]
Soit $A(x), B(x)$ des expressions réelles.
\begin{enumerate}
\item $\sqrt{A(x)} < B(x) \quad\Longleftrightarrow\quad \begin{cases} A(x)\geq 0 \\ B(x)>0 \\ A(x) < [B(x)]^2 \end{cases}$
\item $\sqrt{A(x)} > B(x) \quad\Longleftrightarrow\quad \begin{cases} B(x)<0 \\ A(x)\geq 0 \end{cases} \quad\textbf{OU}\quad \begin{cases} B(x)\geq 0 \\ A(x) > [B(x)]^2 \end{cases}$
\item $\sqrt{A(x)} \leq B(x)$ et $\sqrt{A(x)} \geq B(x)$ : adapter les inégalités larges/strictes ($B(x)\geq 0$ au lieu de $>0$ pour le 2e cas du $\geq$).
\end{enumerate}
\cle{Principe clé :} On n'élève au carré que des quantités \textbf{positives ou nulles} (d'où $B(x)>0$ ou $\geq 0$).
\end{propriete}

\begin{methode}[Résoudre une inéquation irrationnelle]
\begin{enumerate}
\item Identifier le type ($<$, $>$, $\leq$, $\geq$) et écrire le(s) système(s) équivalent(s) selon la propriété ci-dessus.
\item Résoudre chaque condition (inéquations du 1er ou 2nd degré) et faire les intersections (pour un « ET ») ou réunions (pour un « OU »).
\item Conclure par l'ensemble solution $S$.
\end{enumerate}
\end{methode}

\begin{exoresolu}[Inéquation $\sqrt{A}<B$]
Résoudre dans $\mathbb{R}$ : $\sqrt{x+3}<x+1$.
\tcblower
\textbf{Système équivalent :}
\[\begin{cases} x+3\geq 0 & (1) \\ x+1>0 & (2) \\ x+3 < (x+1)^2 & (3) \end{cases}\]
\textbf{(1) :} $x\geq -3$.
\textbf{(2) :} $x > -1$. Intersection (1)$\cap$(2) : $x > -1$.
\textbf{(3) :} $x+3 < x^2+2x+1 \iff 0 < x^2+x-2$.
Trinôme $x^2+x-2$ : $\Delta=9$, racines $-2$ et $1$. Signe $+$ (coeff dir $+$) à l'extérieur : $x<-2$ ou $x>1$.
\textbf{Intersection finale :} $(x>-1) \cap (x<-2 \cup x>1) = (x>1)$.
\textbf{Solution :} $S=]1\,;\,+\infty[$.
\end{exoresolu}

\begin{exoresolu}[Application technique -- Câblage électrique]
La tension $U$ (en volts) aux bornes d'un dipôle non linéaire est liée au courant $I$ (en ampères) par $U=\sqrt{I^2+4I+5}$. Pour quelle intensité $I$ la tension est-elle \textbf{strictement inférieure} à $I+3$ ?
\tcblower
\textbf{Inéquation :} $\sqrt{I^2+4I+5} < I+3$. Variable $I\in\mathbb{R}$.
\textbf{C.E. radicande :} $I^2+4I+5 = (I+2)^2+1 > 0$ pour tout $I$ (toujours vérifiée).
\textbf{Système :} $\begin{cases} I+3 > 0 \iff I > -3 \\ I^2+4I+5 < (I+3)^2 \end{cases}$
\textbf{2e inéquation :} $I^2+4I+5 < I^2+6I+9 \iff 5-9 < 2I \iff -4 < 2I \iff I > -2$.
\textbf{Intersection :} $I > -3$ et $I > -2 \Rightarrow I > -2$.
\textbf{Conclusion :} Pour tout courant $I > \SI{-2}{A}$ (physiquement $I\geq 0$), la tension reste inférieure à $I+3$.
\end{exoresolu}

\begin{exoresolu}[Inéquation $\sqrt{A}>B$ -- Cas disjonctif]
Résoudre : $\sqrt{x-1} > x-3$.
\tcblower
\textbf{C.E. :} $x-1\geq 0 \iff x\geq 1$.
\textbf{Type $>$ :} Union de deux cas.
\textbf{Cas 1 :} $x-3 < 0 \iff x < 3$. Avec C.E. $x\geq 1$ : $1\leq x < 3$. \textit{(Racine $\geq 0$ > nombre négatif : toujours vrai)}.
\textbf{Cas 2 :} $\begin{cases} x-3 \geq 0 \iff x\geq 3 \\ x-1 > (x-3)^2 \end{cases}$.
Inéquation : $x-1 > x^2-6x+9 \iff 0 > x^2-7x+10 \iff x^2-7x+10 < 0$.
Trinôme : $\Delta=49-40=9$, racines $2$ et $5$. Signe $-$ entre $2$ et $5$ : $2<x<5$.
Intersection avec $x\geq 3$ : $3\leq x < 5$.
\textbf{Réunion :} $[1;3[ \cup [3;5[ = [1;5[$.
\textbf{Solution :} $S=[1\,;\,5[$.
\end{exoresolu}

\begin{attention}[Les 5 erreurs fatales au Probatoire]
\begin{enumerate}
\item \textbf{Oublier les C.E.} (Conditions d'Existence) : $A(x)\geq 0$ et $B(x)\geq 0$ (équation) ou $B(x)>0$ (inéquation $<$). Sans C.E., solutions parasites garanties.
\item \textbf{Élever au carré sans signe} : $\sqrt{A}<B \not\Rightarrow A<B^2$ si $B\leq 0$. Toujours écrire le système complet.
\item \textbf{Signe du trinôme inversé} : Relire le coefficient $a$. $a>0 \Rightarrow +$ à l'extérieur ; $a<0 \Rightarrow +$ \emph{entre} les racines. Faire un mini-tableau si doute.
\item \textbf{Crochets ouverts/fermés} : $>$ ou $<$ $\to$ $]$ $[$ (exclus) ; $\geq$ ou $\leq$ $\to$ $[$ $]$ (inclus). Vérifier si les racines appartiennent au domaine.
\item \textbf{Systèmes : couples oubliés} ou 2e inconnue non calculée. Pour $x+y=S, xy=P$, les deux racines donnent \textbf{deux couples} $(x_1;y_2)$ et $(x_2;y_1)$. Pour degré 3 : ne pas oublier le facteur du 2nd degré après la 1ère factorisation.
\end{enumerate}
\end{attention}

\begin{savaistu}[Histoire \& Contexte Cameroun]
Les équations du 2nd degré étaient résolues géométriquement par les Babyloniens (vers 1800 av. J.-C.) pour des problèmes de partage de terres. Al-Khwarizmi (IXe s.) a donné les premières méthodes algébriques systématiques (« al-jabr »). Au Cameroun, ces outils modélisent le dimensionnement des structures (poutres, dalles), l'optimisation des coûts (aire/perimètre) et les caractéristiques électriques (loi des nœuds/mailles menant à des systèmes quadratiques). Le Probatoire attend une rédaction rigoureuse : \textbf{Énoncé $\to$ C.E. $\to$ Résolution $\to$ Vérification $\to$ Conclusion}.
\end{savaistu}

\begin{exercice}[Synthèse -- Préparation Probatoire]
\begin{enumerate}
\item Résoudre dans $\mathbb{R}$ : $\sqrt{3x+1}=x-1$.
\item Résoudre : $x^3-3x^2-4x+12\leq 0$. (Indication : tester $x=2$).
\item Un électricien doit poser un câble le long de l'hypoténuse d'un triangle rectangle de périmètre \SI{30}{m} et d'aire \SI{30}{m^2}. Déterminer la longueur du câble.
\end{enumerate}
\end{exercice}

\begin{corrige}[Exercice Synthèse]
\textbf{1.} C.E. : $3x+1\geq 0 \iff x\geq -1/3$ et $x-1\geq 0 \iff x\geq 1 \Rightarrow \mathcal{D}=[1;+\infty[$.
Équation : $3x+1=x^2-2x+1 \iff x^2-5x=0 \iff x(x-5)=0 \Rightarrow x=0$ ou $x=5$.
Vérification : $0\notin\mathcal{D}$ (parasite). $5\in\mathcal{D}$, $\sqrt{16}=4=5-1$. $S=\{5\}$.

\textbf{2.} $P(x)=x^3-3x^2-4x+12$. $P(2)=8-12-8+12=0$. Racine $2$.
$P(x)=(x-2)(x^2-x-6)$. Trinôme : $\Delta=25$, racines $-2$ et $3$.
$P(x)=(x-2)(x+2)(x-3)$. Racines : $-2 < 2 < 3$.
Tableau signes (coeff dir $+$) : $-$ sur $]-\infty;-2[$, $0$, $+$ sur $]-2;2[$, $0$, $-$ sur $]2;3[$, $0$, $+$ sur $]3;+\infty[$.
Inégalité $\leq 0$ (large) : $S=]-\infty;-2]\cup[2;3]$.

\textbf{3.} Soit $a,b$ cathètes, $c$ hypoténuse (câble).
$\begin{cases} a+b+c=30 \\ ab/2=30 \iff ab=60 \\ a^2+b^2=c^2 \end{cases}$.
$a+b=30-c$. $(a+b)^2=a^2+b^2+2ab \Rightarrow (30-c)^2=c^2+120$.
$900-60c+c^2=c^2+120 \Rightarrow 60c=780 \Rightarrow c=13$.
Vérification : $a+b=17, ab=60 \Rightarrow a,b$ racines de $X^2-17X+60=0 \Rightarrow 5$ et $12$. Triangle 5-12-13 valide.
\textbf{Longueur câble : \SI{13}{m}.}
\end{corrige}

\newpage

\coursec{Exercices}

\courssub{Je vérifie mes connaissances}

\begin{exercice}[QCM : une seule bonne réponse]
Pour chaque question, choisir la bonne réponse sans justification.
\begin{enumerate}
\item Deux nombres réels ont pour somme $S = 7$ et pour produit $P = 12$. Ils sont solutions de l'équation :
\begin{enumerate}
\item $X^2 - 7X + 12 = 0$ \qquad \textbf{b)} $X^2 + 7X + 12 = 0$ \qquad \textbf{c)} $X^2 - 12X + 7 = 0$
\end{enumerate}
\item Le nombre $2$ est racine du polynôme $P(x) = x^3 - 3x^2 + 4$ si et seulement si :
\begin{enumerate}
\item $P(2) = 0$ \qquad \textbf{b)} $P(0) = 2$ \qquad \textbf{c)} $P(x)$ est divisible par $x + 2$
\end{enumerate}
\item L'équation $\sqrt{x - 1} = -3$ admet dans $\mathbb{R}$ :
\begin{enumerate}
\item une solution \qquad \textbf{b)} deux solutions \qquad \textbf{c)} aucune solution
\end{enumerate}
\item L'ensemble de définition de l'équation $\sqrt{2x - 6} = x - 4$ est :
\begin{enumerate}
\item $[4\,;+\infty[$ \qquad \textbf{b)} $[3\,;+\infty[$ \qquad \textbf{c)} $\mathbb{R}$
\end{enumerate}
\end{enumerate}
\end{exercice}

\begin{exercice}[Vrai ou faux, en justifiant]
Répondre par VRAI ou FAUX et justifier chaque réponse.
\begin{enumerate}
\item Si $a$ est racine du polynôme $P$, alors $P(x)$ se factorise par $(x - a)$.
\item Tout polynôme de degré $3$ admet au moins une racine réelle évidente entière.
\item Pour résoudre $\sqrt{f(x)} = g(x)$, il suffit d'élever les deux membres au carré.
\item Un polynôme de degré $3$ ne peut pas avoir trois racines réelles distinctes.
\end{enumerate}
\end{exercice}

\begin{exercice}[Définitions et cours]
\begin{enumerate}
\item Donner la définition d'une racine (ou zéro) d'un polynôme $P$.
\item Énoncer la condition d'existence des solutions d'une équation de somme et produit : deux réels de somme $S$ et de produit $P$ existent si et seulement si \ldots
\item Rappeler la méthode de résolution d'une équation irrationnelle du type $\sqrt{f(x)} = g(x)$ (conditions à poser, étapes, vérification).
\end{enumerate}
\end{exercice}

\begin{exercice}[Calculs directs]
\begin{enumerate}
\item Résoudre dans $\mathbb{R}$ l'équation $x^2 - 5x + 6 = 0$.
\item Vérifier que $1$ est racine du polynôme $P(x) = x^3 - 2x^2 - x + 2$.
\item Déterminer le domaine de résolution de l'équation $\sqrt{5 - x} = x + 1$.
\item Trouver deux nombres réels dont la somme est $10$ et le produit $9$.
\end{enumerate}
\end{exercice}

\courssub{Je m'entraîne}

\begin{exercice}[Systèmes somme--produit]
\begin{enumerate}
\item Résoudre dans $\mathbb{R}^2$ le système $\begin{cases} x + y = 7 \\ xy = 12 \end{cases}$ et interpréter géométriquement les solutions (dimensions d'un rectangle).
\item Résoudre dans $\mathbb{R}^2$ le système $\begin{cases} x - y = 1 \\ xy = 20 \end{cases}$.
\item Deux nombres réels ont pour somme $10$ et pour produit $30$. Existent-ils ? Justifier à l'aide du discriminant.
\end{enumerate}
\end{exercice}

\begin{exercice}[Factorisation d'un polynôme de degré 3]
Soit $P(x) = 2x^3 - 3x^2 - 3x + 2$.
\begin{enumerate}
\item Vérifier que $-1$ est racine de $P$.
\item Déterminer les réels $a$, $b$ et $c$ tels que $P(x) = (x + 1)(ax^2 + bx + c)$.
\item Factoriser complètement $P(x)$ en produit de facteurs du premier degré.
\item Résoudre dans $\mathbb{R}$ l'équation $P(x) = 0$.
\end{enumerate}
\end{exercice}

\begin{exercice}[Signe d'un polynôme de degré 3]
Soit $Q(x) = x^3 - 2x^2 - 5x + 6$.
\begin{enumerate}
\item Vérifier que $1$ est racine de $Q$, puis factoriser complètement $Q(x)$.
\item Dresser le tableau de signes de $Q(x)$ sur $\mathbb{R}$.
\item En déduire l'ensemble des solutions de l'inéquation $Q(x) > 0$, puis de $Q(x) \leq 0$.
\end{enumerate}
\end{exercice}

\begin{exercice}[Équations et inéquations irrationnelles]
Résoudre dans $\mathbb{R}$, en vérifiant chaque solution obtenue :
\begin{enumerate}
\item $\sqrt{x + 7} = \sqrt{2x - 3}$ ;
\item $\sqrt{3x - 5} = x - 3$ ;
\item $\sqrt{x - 1} \leq 2$ ;
\item $\sqrt{2x + 3} \geq x$ (on discutera selon le signe du membre de droite).
\end{enumerate}
\end{exercice}

\courssub{Je me prépare au Probatoire}

\begin{exercice}[Le terrain de l'atelier]
Un lycée technique de Douala veut aménager un terrain rectangulaire pour ses travaux pratiques. Le demi-périmètre du terrain est de \SI{26}{m} et son aire de \SI{168}{m^2}. On note $x$ et $y$ les dimensions du terrain, avec $x \geq y > 0$.
\begin{enumerate}
\item Montrer que $x$ et $y$ sont solutions du système $\begin{cases} x + y = 26 \\ xy = 168 \end{cases}$.
\item En déduire que $x$ et $y$ sont solutions de l'équation $X^2 - 26X + 168 = 0$, puis déterminer les dimensions du terrain.
\item Calculer la longueur $d$ de la diagonale du terrain.
\item Vérifier que $d$ est solution de l'équation irrationnelle $\sqrt{X^2 - 144} = 14$.
\item Le chef des travaux veut clôturer le terrain avec une grille vendue \SI{2500}{F} le mètre. Calculer le coût total de la clôture.
\end{enumerate}
\end{exercice}

\begin{exercice}[Étude complète d'un polynôme de degré 3]
Un atelier fabrique des caisses. Le coût de fabrication, en milliers de francs, de $x$ caisses ($x \geq 0$) est modélisé par $C(x) = x^3 - 6x^2 + 9x$, et la recette par $R(x) = 4x$. On pose $B(x) = R(x) - C(x)$ le bénéfice.
\begin{enumerate}
\item Montrer que $B(x) = -x^3 + 6x^2 - 5x$, puis que $B(x) = -x(x^2 - 6x + 5)$.
\item Factoriser complètement le trinôme $x^2 - 6x + 5$.
\item Dresser le tableau de signes de $B(x)$ sur $[0\,;+\infty[$.
\item En déduire les quantités de caisses pour lesquelles l'atelier réalise un bénéfice positif ou nul.
\item Résoudre dans $\mathbb{R}$ l'équation $B(x) = 0$ et interpréter les solutions positives dans le contexte de l'atelier.
\end{enumerate}
\end{exercice}

\begin{exercice}[Démonstration guidée et équations irrationnelles]
\textbf{Partie A : démonstration de cours.} Soit $P$ un polynôme de degré $3$ et $a$ un réel tel que $P(a) = 0$. On veut montrer que $P(x)$ est divisible par $(x - a)$. On écrit $P(x) = (x - a)(px^2 + qx + r) + s$ où $p, q, r, s$ sont des réels.
\begin{enumerate}
\item En calculant $P(a)$, montrer que $s = 0$.
\item Conclure que $P(x) = (x - a)(px^2 + qx + r)$.
\end{enumerate}
\textbf{Partie B : application.} Soit $P(x) = x^3 - 4x^2 + x + 6$.
\begin{enumerate}[resume]
\item Trouver une racine évidente de $P$ parmi les entiers $-1$, $1$, $2$, $3$.
\item Factoriser complètement $P(x)$, puis résoudre l'inéquation $P(x) \geq 0$.
\end{enumerate}
\textbf{Partie C : équations irrationnelles.}
\begin{enumerate}[resume]
\item Résoudre dans $\mathbb{R}$ l'équation $\sqrt{x^2 - 5} = x - 1$ en précisant le domaine de résolution et en vérifiant les solutions.
\item Un technicien mesure que la diagonale $d$ d'une plaque carrée de côté $c$ vérifie $d = \sqrt{2c^2}$. Déterminer $c$ sachant que $d = \SI{6}{m}$.
\end{enumerate}
\end{exercice}

\newpage

\coursec{Corrigés des exercices}

\coursec{Corrigés détaillés — Leçon 07 : Équations, inéquations polynômiales et irrationnelles}

\begin{corrige}[Exercice 1 — QCM : une seule bonne réponse]
\begin{enumerate}
\item \textbf{a)} $X^2 - 7X + 12 = 0$.\\
Deux réels de somme $S$ et de produit $P$ sont les solutions de l'équation $X^2 - SX + P = 0$. Ici $S=7$, $P=12$, d'où $X^2 - 7X + 12 = 0$.
\item \textbf{a)} $P(2) = 0$.\\
Par définition, $a$ est \cle{racine} d'un polynôme $P$ si et seulement si $P(a)=0$.
\item \textbf{c)} aucune solution.\\
La fonction racine carrée ne prend que des valeurs positives ou nulles dans $\mathbb{R}$ ; elle ne peut pas être égale à $-3$.
\item \textbf{b)} $[3\,;+\infty[$.\\
L'équation $\sqrt{2x-6}=x-4$ n'a de sens que si le radicande est positif ou nul : $2x-6 \ge 0 \iff x \ge 3$. Le second membre $x-4$ est défini sur $\mathbb{R}$. L'\cle{ensemble de définition} est donc $[3\,;+\infty[$.
\end{enumerate}
\end{corrige}

\begin{corrige}[Exercice 2 — Vrai ou faux, en justifiant]
\begin{enumerate}
\item \textbf{VRAI}. Si $a$ est racine de $P$, alors $P(a)=0$. D'après le \cle{théorème du facteur}, $P(x)$ est divisible par $(x-a)$, c'est-à-dire qu'il existe un polynôme $Q$ tel que $P(x)=(x-a)Q(x)$.
\item \textbf{FAUX}. Un polynôme de degré 3 admet toujours au moins une racine réelle (théorème des valeurs intermédiaires), mais cette racine n'est pas nécessairement un entier évident. Exemple : $P(x)=x^3-2$ a pour unique racine réelle $\sqrt[3]{2}$, qui n'est pas un entier.
\item \textbf{FAUX}. Élever au carré ne suffit pas ; il faut au préalable poser les \cle{conditions d'existence} : $f(x)\ge 0$ (radicande) et $g(x)\ge 0$ (car $\sqrt{f(x)}\ge 0$). Après résolution, il faut \cle{vérifier} les solutions dans l'équation initiale.
\item \textbf{FAUX}. Un polynôme de degré 3 peut avoir trois racines réelles distinctes. Exemple : $P(x)=(x-1)(x-2)(x-3)=x^3-6x^2+11x-6$ admet les trois racines $1$, $2$, $3$.
\end{enumerate}
\end{corrige}

\begin{corrige}[Exercice 3 — Définitions et cours]
\begin{enumerate}
\item Un réel $a$ est une \textbf{racine} (ou \textbf{zéro}) d'un polynôme $P$ si et seulement si $P(a)=0$.
\item Deux réels de somme $S$ et de produit $P$ existent dans $\mathbb{R}$ si et seulement si le \cle{discriminant} $\Delta = S^2 - 4P$ est \textbf{positif ou nul} ($\Delta \ge 0$).
\item \textbf{Méthode de résolution de $\sqrt{f(x)} = g(x)$ :}
\begin{enumerate}
\item Déterminer l'ensemble de définition : $f(x) \ge 0$.
\item Poser la condition $g(x) \ge 0$ (car la racine carrée est toujours $\ge 0$).
\item Élever les deux membres au carré : $f(x) = [g(x)]^2$.
\item Résoudre l'équation obtenue.
\item \textbf{Vérifier} chaque solution candidate : elle doit appartenir à l'ensemble de définition et satisfaire $g(x) \ge 0$ et l'égalité initiale.
\end{enumerate}
\end{enumerate}
\end{corrige}

\begin{corrige}[Exercice 4 — Calculs directs]
\begin{enumerate}
\item $x^2 - 5x + 6 = 0$. $\Delta = (-5)^2 - 4\times 6 = 25-24=1$. $x = \frac{5 \pm 1}{2}$, donc $x_1 = 3$, $x_2 = 2$. Solutions : $\boxed{\{2\,;\,3\}}$.
\item $P(x)=x^3-2x^2-x+2$. $P(1)=1-2-1+2=0$. Donc $1$ est bien une racine de $P$.
\item $\sqrt{5-x}=x+1$. Conditions : $5-x \ge 0 \iff x \le 5$ et $x+1 \ge 0 \iff x \ge -1$. Domaine de résolution : $\boxed{[-1\,;\,5]}$.
\item Soient $x$ et $y$ les deux nombres. $x+y=10$, $xy=9$. Ils sont solutions de $X^2-10X+9=0$. $\Delta = 100-36=64$. $X = \frac{10 \pm 8}{2}$, soit $9$ et $1$. Les deux nombres sont $\boxed{9 \text{ et } 1}$.
\end{enumerate}
\end{corrige}

\begin{corrige}[Exercice 5 — Systèmes se ramenant à une équation du second degré]
\begin{enumerate}
\item Système $\begin{cases} x+y=7 \\ xy=12 \end{cases}$. $x$ et $y$ sont les racines de $X^2-7X+12=0$. $\Delta=49-48=1$. $X=\frac{7\pm1}{2}$ : $4$ et $3$. Solutions : $(x,y)=(4,3)$ ou $(3,4)$. Géométriquement, ce sont les dimensions (longueur et largeur) d'un rectangle de demi-périmètre $7$ et d'aire $12$.
\item Système $\begin{cases} x-y=1 \\ xy=20 \end{cases}$. Posons $x=y+1$. Alors $(y+1)y=20 \iff y^2+y-20=0$. $\Delta=1+80=81$. $y=\frac{-1\pm9}{2}$ : $y=4$ ou $y=-5$. Correspondamment $x=5$ ou $x=-4$. Solutions : $\boxed{(5,4) \text{ et } (-4,-5)}$.
\item Somme $S=10$, produit $P=30$. $\Delta = S^2-4P = 100-120 = -20 < 0$. Il n'existe \textbf{pas} de tels nombres réels.
\end{enumerate}
\end{corrige}

\begin{corrige}[Exercice 6 — Factorisation d'un polynôme de degré 3]
Soit $P(x)=2x^3-3x^2-3x+2$.
\begin{enumerate}
\item $P(-1)=2(-1)^3-3(-1)^2-3(-1)+2 = -2-3+3+2=0$. Donc $-1$ est racine.
\item On cherche $a,b,c$ tels que $P(x)=(x+1)(ax^2+bx+c)$. En développant : $ax^3+(a+b)x^2+(b+c)x+c$. Identification :
$\begin{cases} a=2 \\ a+b=-3 \Rightarrow b=-5 \\ b+c=-3 \Rightarrow -5+c=-3 \Rightarrow c=2 \\ c=2 \text{ (vérifié)} \end{cases}$
Donc $P(x)=(x+1)(2x^2-5x+2)$.
\item Factorisation du trinôme : $2x^2-5x+2$. $\Delta=25-16=9$. Racines : $\frac{5\pm3}{4}$ : $2$ et $\frac12$. $2x^2-5x+2 = 2(x-2)(x-\frac12) = (x-2)(2x-1)$.
Factorisation complète : $\boxed{P(x) = (x+1)(x-2)(2x-1)}$.
\item $P(x)=0 \iff (x+1)(x-2)(2x-1)=0$. Solutions : $\boxed{x=-1,\; x=2,\; x=\frac12}$.
\end{enumerate}
\end{corrige}

\begin{corrige}[Exercice 7 — Signe d'un polynôme de degré 3]
Soit $Q(x)=x^3-2x^2-5x+6$.
\begin{enumerate}
\item $Q(1)=1-2-5+6=0$, donc $1$ est racine. Division par $(x-1)$ (division synthétique par $1$) :
Coefficients : $1,\;-2,\;-5,\;6$.
$1 \rightarrow 1$ ; $1\times1=1 \rightarrow -2+1=-1$ ; $-1\times1=-1 \rightarrow -5-1=-6$ ; $-6\times1=-6 \rightarrow 6-6=0$.
Quotient : $x^2 - x - 6 = (x-3)(x+2)$.
Factorisation complète : $\boxed{Q(x) = (x-1)(x-3)(x+2)}$.
\item \textbf{Tableau de signes de $Q(x)$ sur $\mathbb{R}$ :}
\begin{center}
\begin{tabular}{|c|ccccccccc|}\hline
$x$ & $-\infty$ & & $-2$ & & $1$ & & $3$ & & $+\infty$ \\ \hline
$x+2$ & $-$ & $-$ & $0$ & $+$ & $+$ & $+$ & $+$ & $+$ & $+$ \\
$x-1$ & $-$ & $-$ & $-$ & $-$ & $0$ & $+$ & $+$ & $+$ & $+$ \\
$x-3$ & $-$ & $-$ & $-$ & $-$ & $-$ & $-$ & $0$ & $+$ & $+$ \\ \hline
$Q(x)$ & $-$ & $-$ & $0$ & $+$ & $0$ & $-$ & $0$ & $+$ & $+$ \\ \hline
\end{tabular}
\end{center}
\item $Q(x) > 0 \iff x \in \boxed{]-2\,;\,1[ \;\cup\; ]3\,;\,+\infty[}$.
$Q(x) \le 0 \iff x \in \boxed{]-\infty\,;\,-2] \;\cup\; [1\,;\,3]}$.
\end{enumerate}
\end{corrige}

\begin{corrige}[Exercice 8 — Équations et inéquations irrationnelles]
\begin{enumerate}
\item $\sqrt{x+7} = \sqrt{2x-3}$.
Domaine : $x+7\ge0 \Rightarrow x\ge-7$ ; $2x-3\ge0 \Rightarrow x\ge\frac32$. Donc $D=[\frac32,+\infty[$.
Élever au carré : $x+7 = 2x-3 \Rightarrow x=10$.
$10 \in D$, vérification : $\sqrt{17}=\sqrt{17}$. Solution : $\boxed{x=10}$.
\item $\sqrt{3x-5} = x-3$.
Domaine : $3x-5\ge0 \Rightarrow x\ge\frac53$ ; $x-3\ge0 \Rightarrow x\ge3$. Donc $D=[3,+\infty[$.
Élever au carré : $3x-5 = (x-3)^2 = x^2-6x+9 \Rightarrow x^2-9x+14=0$.
$\Delta=81-56=25$. $x=\frac{9\pm5}{2}$ : $x=7$ ou $x=2$.
$2 \notin D$ (rejeté). $7 \in D$, vérification : $\sqrt{16}=4$, $7-3=4$. Solution : $\boxed{x=7}$.
\item $\sqrt{x-1} \le 2$.
Domaine : $x-1\ge0 \Rightarrow x\ge1$.
Les deux membres sont positifs sur $D$, on élève au carré : $x-1 \le 4 \Rightarrow x \le 5$.
Intersection avec $D$ : $\boxed{[1\,;\,5]}$.
\item $\sqrt{2x+3} \ge x$.
Domaine : $2x+3\ge0 \Rightarrow x\ge-\frac32$.
\underline{Cas 1 :} $x < 0$. Alors $x$ est négatif, $\sqrt{2x+3}\ge0$, l'inéquation est vraie pour tout $x$ du domaine. Donc $-\frac32 \le x < 0$.
\underline{Cas 2 :} $x \ge 0$. Les deux membres sont $\ge0$, on élève au carré : $2x+3 \ge x^2 \iff x^2-2x-3 \le 0 \iff (x-3)(x+1)\le0 \iff x\in[-1,3]$.
Intersection avec $x\ge0$ : $[0,3]$.
Réunion des deux cas : $\boxed{[-\frac32\,;\,3]}$.
\end{enumerate}
\end{corrige}

\begin{corrige}[Exercice 9 — Le terrain de l'atelier]
\begin{enumerate}
\item Le demi-périmètre est $26$ m, donc $x+y=26$. L'aire est $168$ m$^2$, donc $xy=168$. Le système est bien $\begin{cases} x+y=26 \\ xy=168 \end{cases}$.
\item $x$ et $y$ sont les racines de $X^2-26X+168=0$. $\Delta = 26^2-4\times168 = 676-672=4$. $X=\frac{26\pm2}{2}$ : $14$ et $12$. Comme $x\ge y>0$, $x=14$ m, $y=12$ m.
\item Diagonale $d = \sqrt{x^2+y^2} = \sqrt{14^2+12^2} = \sqrt{196+144} = \sqrt{340} = 2\sqrt{85}$ m (environ $\num{18,44}$ m).
\item Vérification : $d^2 = 340$. $\sqrt{d^2-144} = \sqrt{340-144} = \sqrt{196} = 14$. Donc $d$ est bien solution de $\sqrt{X^2-144}=14$.
\item Périmètre $= 2(x+y)=52$ m. Coût $= 52 \times 2500 = \num{130000}$ F.
\end{enumerate}
\textbf{Barème indicatif (sur 6 points) :}
\begin{itemize}
\item Q1 : 1 pt (système correct).
\item Q2 : 2 pts (équation, discriminant, dimensions).
\item Q3 : 1 pt (diagonale sous forme exacte $2\sqrt{85}$).
\item Q4 : 1 pt (vérification de l'équation irrationnelle).
\item Q5 : 1 pt (coût avec unité).
\end{itemize}
\textbf{Points de vigilance :} Justifier le passage au système (demi-périmètre/aire), préciser $x\ge y$ pour choisir l'ordre, donner la diagonale sous forme exacte ($2\sqrt{85}$), vérifier l'équation irrationnelle en remplaçant $X$ par $d$, ne pas oublier les unités.
\end{corrige}

\begin{corrige}[Exercice 10 — Étude complète d'un polynôme de degré 3 : bénéfice]
\begin{enumerate}
\item $B(x)=R(x)-C(x)=4x-(x^3-6x^2+9x) = -x^3+6x^2-5x = -x(x^2-6x+5)$.
\item $x^2-6x+5$ : $\Delta=36-20=16$. Racines : $\frac{6\pm4}{2}=5$ et $1$. Donc $x^2-6x+5=(x-5)(x-1)$.
$B(x) = -x(x-5)(x-1)$.
\item \textbf{Tableau de signes de $B(x)$ sur $[0\,;+\infty[$ :}
\begin{center}
\begin{tabular}{|c|ccccccc|}\hline
$x$ & $0$ & & $1$ & & $5$ & & $+\infty$ \\ \hline
$-x$ & $0$ & $-$ & $-$ & $-$ & $-$ & $-$ & $-$ \\
$x-1$ & $-$ & $-$ & $0$ & $+$ & $+$ & $+$ & $+$ \\
$x-5$ & $-$ & $-$ & $-$ & $-$ & $0$ & $+$ & $+$ \\ \hline
$B(x)$ & $0$ & $-$ & $0$ & $+$ & $0$ & $-$ & $-$ \\ \hline
\end{tabular}
\end{center}
\item $B(x) \ge 0 \iff x \in \boxed{[1\,;\,5]}$. L'atelier réalise un bénéfice positif ou nul pour une production comprise entre $1$ et $5$ caisses (inclus).
\item $B(x)=0 \iff -x(x-1)(x-5)=0 \iff x=0,\;1,\;5$.
Solutions positives : $x=1$ et $x=5$. \textbf{Interprétation :} Ce sont les \cle{seuils de rentabilité} (points morts) : pour $1$ caisse et pour $5$ caisses, le bénéfice est nul ; entre ces deux valeurs, l'atelier est bénéficiaire.
\end{enumerate}
\textbf{Barème indicatif (sur 7 points) :}
\begin{itemize}
\item Q1 : 1 pt (expression de $B$, factorisation par $-x$).
\item Q2 : 1 pt (factorisation du trinôme).
\item Q3 : 2 pts (tableau de signes correct sur $[0,+\infty[$ avec bornes $0,1,5$).
\item Q4 : 1 pt (intervalle $[1,5]$).
\item Q5 : 2 pts (résolution $B(x)=0$, interprétation des seuils).
\end{itemize}
\textbf{Points de vigilance :} Ne pas oublier le signe moins devant $x$ dans la factorisation, bien prendre l'intervalle $[0,+\infty[$ pour le tableau de signes ($x$ nombre de caisses), mentionner que $x\ge0$, interpréter les racines positives en langage courant (seuils de rentabilité).
\end{corrige}

\begin{corrige}[Exercice 11 — Démonstration guidée et équations irrationnelles]
\textbf{Partie A : démonstration du théorème du facteur.}
\begin{enumerate}
\item Soit $P(x) = (x-a)(px^2+qx+r) + s$ (division euclidienne par $x-a$). En évaluant en $a$ : $P(a) = (a-a)(\cdots) + s = s$. Comme $P(a)=0$ par hypothèse, on a $s=0$.
\item Donc $P(x) = (x-a)(px^2+qx+r)$, ce qui prouve que $P(x)$ est divisible par $(x-a)$.
\end{enumerate}
\textbf{Partie B : application.}
$P(x)=x^3-4x^2+x+6$.
\begin{enumerate}[resume]
\item Test des diviseurs du terme constant ($\pm 1, \pm 2, \pm 3, \pm 6$) :
$P(-1)=-1-4-1+6=0$ ; $P(1)=1-4+1+6=4\neq0$ ; $P(2)=8-16+2+6=0$ ; $P(3)=27-36+3+6=0$.
Une racine évidente est $-1$ (ou $2$ ou $3$). On choisit $-1$.
\item Division de $P(x)$ par $(x+1)$ (division synthétique par $-1$) :
Coefficients : $1,\;-4,\;1,\;6$.
$1 \rightarrow 1$ ; $1\times(-1)=-1 \rightarrow -4-1=-5$ ; $-5\times(-1)=5 \rightarrow 1+5=6$ ; $6\times(-1)=-6 \rightarrow 6-6=0$.
Quotient : $x^2-5x+6 = (x-2)(x-3)$.
Factorisation complète : $\boxed{P(x) = (x+1)(x-2)(x-3)}$.
\item Inéquation $P(x) \ge 0$. Racines : $-1,\;2,\;3$. \textbf{Tableau de signes :}
\begin{center}
\begin{tabular}{|c|ccccccccc|}\hline
$x$ & $-\infty$ & & $-1$ & & $2$ & & $3$ & & $+\infty$ \\ \hline
$x+1$ & $-$ & $-$ & $0$ & $+$ & $+$ & $+$ & $+$ & $+$ & $+$ \\
$x-2$ & $-$ & $-$ & $-$ & $-$ & $0$ & $+$ & $+$ & $+$ & $+$ \\
$x-3$ & $-$ & $-$ & $-$ & $-$ & $-$ & $-$ & $0$ & $+$ & $+$ \\ \hline
$P(x)$ & $-$ & $-$ & $0$ & $+$ & $0$ & $-$ & $0$ & $+$ & $+$ \\ \hline
\end{tabular}
\end{center}
$P(x) \ge 0 \iff x \in \boxed{[-1\,;\,2] \;\cup\; [3\,;\,+\infty[}$.
\end{enumerate}
\textbf{Partie C : équations irrationnelles.}
\begin{enumerate}[resume]
\item $\sqrt{x^2-5} = x-1$.
Domaine : $x^2-5 \ge 0 \iff x \le -\sqrt5$ ou $x \ge \sqrt5$. Condition $x-1 \ge 0 \iff x \ge 1$. Intersection : $x \ge \sqrt5$ (car $\sqrt5 \approx 2,236 > 1$).
Élever au carré : $x^2-5 = (x-1)^2 = x^2-2x+1 \Rightarrow -5 = -2x+1 \Rightarrow 2x=6 \Rightarrow x=3$.
$3 \ge \sqrt5$, vérification : $\sqrt{9-5}=2$, $3-1=2$. Solution : $\boxed{x=3}$.
\item $d = \sqrt{2c^2} = c\sqrt2$. Donnée $d=6$ m. $c\sqrt2 = 6 \Rightarrow c = \frac{6}{\sqrt2} = 3\sqrt2$ m.
\end{enumerate}
\textbf{Barème indicatif (sur 8 points) :}
\begin{itemize}
\item Partie A : 2 pts (calcul de $P(a)$, conclusion $s=0$).
\item Partie B : 3 pts (recherche racine, factorisation complète, tableau de signes et résolution inéquation).
\item Partie C : 3 pts (équation irrationnelle avec domaine et vérification, application diagonale).
\end{itemize}
\textbf{Points de vigilance :} Dans la démonstration, bien écrire $P(a)=s$ et conclure $s=0$. Pour la factorisation, vérifier toutes les racines entières candidates. Pour le tableau de signes, aligner correctement les zéros et les signes sur les intervalles. Pour l'équation irrationnelle, ne pas oublier la condition $x-1\ge0$ et la vérification finale. Pour la diagonale, simplifier $\frac{6}{\sqrt2}$ en $3\sqrt2$.
\end{corrige}

\newpage

\coursec{Fiche bilan}

\begin{popfigure}
\begin{center}\tcbox[colback=popblue,colframe=popblue,coltext=white,arc=9pt,boxsep=4pt,left=6pt,right=6pt]{\bfseries\small Équations et inéquations}\end{center}
\vspace{-2pt}
\begin{tcbraster}[raster columns=3,raster equal height=rows,raster column skip=6pt,raster row skip=6pt,enhanced,blankest]
\begin{tcolorbox}[enhanced,colback=white,colframe=poporange,boxrule=0.9pt,arc=3pt,title={Second degré},fonttitle=\bfseries\scriptsize,coltitle=white,colbacktitle=poporange,left=4pt,right=4pt,top=2pt,bottom=2pt,boxsep=1pt,toptitle=1.5pt,bottomtitle=1.5pt,halign=left]
\begin{itemize}[nosep,leftmargin=*,itemsep=1pt,font=\scriptsize]
\item $\Delta=b^2-4ac$
\item Signe de $ax^2+bx+c$
\end{itemize}
\end{tcolorbox}
\begin{tcolorbox}[enhanced,colback=white,colframe=popgreen,boxrule=0.9pt,arc=3pt,title={Systèmes},fonttitle=\bfseries\scriptsize,coltitle=white,colbacktitle=popgreen,left=4pt,right=4pt,top=2pt,bottom=2pt,boxsep=1pt,toptitle=1.5pt,bottomtitle=1.5pt,halign=left]
\begin{itemize}[nosep,leftmargin=*,itemsep=1pt,font=\scriptsize]
\item $S=x+y$, $P=xy$
\item $t^2-St+P=0$
\end{itemize}
\end{tcolorbox}
\begin{tcolorbox}[enhanced,colback=white,colframe=poppurple,boxrule=0.9pt,arc=3pt,title={Polynômes degré 3},fonttitle=\bfseries\scriptsize,coltitle=white,colbacktitle=poppurple,left=4pt,right=4pt,top=2pt,bottom=2pt,boxsep=1pt,toptitle=1.5pt,bottomtitle=1.5pt,halign=left]
\begin{itemize}[nosep,leftmargin=*,itemsep=1pt,font=\scriptsize]
\item Racine évidente $a$
\item $P(x)=(x-a)Q(x)$
\end{itemize}
\end{tcolorbox}
\begin{tcolorbox}[enhanced,colback=white,colframe=poppink,boxrule=0.9pt,arc=3pt,title={Signe degré 3},fonttitle=\bfseries\scriptsize,coltitle=white,colbacktitle=poppink,left=4pt,right=4pt,top=2pt,bottom=2pt,boxsep=1pt,toptitle=1.5pt,bottomtitle=1.5pt,halign=left]
\begin{itemize}[nosep,leftmargin=*,itemsep=1pt,font=\scriptsize]
\item Tableau de signes
\item Inéquations
\end{itemize}
\end{tcolorbox}
\begin{tcolorbox}[enhanced,colback=white,colframe=popgold,boxrule=0.9pt,arc=3pt,title={Équations irrationnelles},fonttitle=\bfseries\scriptsize,coltitle=white,colbacktitle=popgold,left=4pt,right=4pt,top=2pt,bottom=2pt,boxsep=1pt,toptitle=1.5pt,bottomtitle=1.5pt,halign=left]
\begin{itemize}[nosep,leftmargin=*,itemsep=1pt,font=\scriptsize]
\item Condition $A\ge 0$
\item Élever au carré
\end{itemize}
\end{tcolorbox}
\begin{tcolorbox}[enhanced,colback=white,colframe=popblue!60,boxrule=0.9pt,arc=3pt,title={Inéquations irrationnelles},fonttitle=\bfseries\scriptsize,coltitle=white,colbacktitle=popblue!60,left=4pt,right=4pt,top=2pt,bottom=2pt,boxsep=1pt,toptitle=1.5pt,bottomtitle=1.5pt,halign=left]
\begin{itemize}[nosep,leftmargin=*,itemsep=1pt,font=\scriptsize]
\item Cas selon le signe
\item Vérifier les solutions
\end{itemize}
\end{tcolorbox}
\end{tcbraster}

\legende{Carte mentale de la leçon : équations, inéquations polynômiales et irrationnelles.}
\end{popfigure}

\begin{bilan}
\textbf{1. Définitions clés}
\begin{itemize}
  \item \cle{Discriminant} de $ax^2+bx+c$ ($a\neq 0$) : $\Delta=b^2-4ac$.
  \item \cle{Racine} d'un polynôme $P$ : un nombre $a$ tel que $P(a)=0$.
  \item \cle{Équation irrationnelle} : équation où l'inconnue apparaît sous un radical, par exemple $\sqrt{A(x)}=B(x)$.
\end{itemize}

\textbf{2. Formules et théorèmes à connaître par cœur}
\begin{center}
\begin{tabularx}{\linewidth}{|l|X|}
\hline
\rowcolor{popblueL} \textbf{Notion} & \textbf{Résultat} \\
\hline
$\Delta>0$ & Deux racines : $x_1=\dfrac{-b-\sqrt{\Delta}}{2a}$, $x_2=\dfrac{-b+\sqrt{\Delta}}{2a}$ \\
\hline
$\Delta=0$ & Racine double : $x_0=-\dfrac{b}{2a}$ \\
\hline
$\Delta<0$ & Pas de racine réelle ; $ax^2+bx+c$ du signe de $a$ sur $\mathbb{R}$ \\
\hline
Somme et produit & $x_1+x_2=-\dfrac{b}{a}$ ; $x_1x_2=\dfrac{c}{a}$ \\
\hline
Système somme--produit & $x$ et $y$ sont racines de $t^2-St+P=0$ \\
\hline
Factorisation & Si $P(a)=0$, alors $P(x)=(x-a)Q(x)$ avec $\deg Q=\deg P-1$ \\
\hline
Signe de $ax^2+bx+c$ & Signe de $a$ à l'extérieur des racines, signe de $-a$ entre les racines \\
\hline
\end{tabularx}
\end{center}

\textbf{3. Méthodes express}
\begin{itemize}
  \item \textbf{Factoriser un polynôme de degré 3} : chercher une racine évidente ($\pm1$, $\pm2$, $\pm3$\dots), écrire $P(x)=(x-a)(ax^2+bx+c)$ par identification, puis résoudre $Q(x)=0$ avec $\Delta$.
  \item \textbf{Signe d'un polynôme de degré 3} : factoriser complètement, dresser un tableau de signes avec chaque facteur.
  \item \textbf{Équation $\sqrt{A}=B$} : imposer $B\ge 0$, élever au carré $A=B^2$, résoudre, puis \textbf{vérifier} chaque solution dans l'équation de départ.
  \item \textbf{Inéquation $\sqrt{A}<B$} : équivalent à $A\ge 0$, $B>0$ et $A<B^2$. \textbf{Inéquation $\sqrt{A}>B$} : si $B<0$, il suffit que $A\ge 0$ ; si $B\ge 0$, alors $A>B^2$.
\end{itemize}
\end{bilan}

\begin{aretenir}[Checklist avant le Probatoire]
\begin{itemize}
  \item[$\square$] Je sais calculer $\Delta$ et conclure selon son signe.
  \item[$\square$] Je connais les formules $x_1+x_2=-\dfrac{b}{a}$ et $x_1x_2=\dfrac{c}{a}$.
  \item[$\square$] Je sais résoudre un système $\begin{cases} x+y=S \\ xy=P \end{cases}$ via $t^2-St+P=0$.
  \item[$\square$] Je sais trouver une racine évidente d'un polynôme de degré 3.
  \item[$\square$] Je sais factoriser $P(x)$ par $(x-a)$ par identification des coefficients.
  \item[$\square$] Je sais dresser le tableau de signes d'un polynôme de degré 3 factorisé.
  \item[$\square$] Je sais résoudre une inéquation polynômiale et donner l'ensemble des solutions en intervalles.
  \item[$\square$] Je pense à la condition d'existence $A(x)\ge 0$ pour $\sqrt{A(x)}$.
  \item[$\square$] Je n'élève au carré que lorsque les deux membres sont positifs.
  \item[$\square$] Je vérifie toujours les solutions d'une équation irrationnelle.
  \item[$\square$] Je distingue les deux cas pour $\sqrt{A}>B$ selon le signe de $B$.
  \item[$\square$] Je rédige ma conclusion avec des intervalles et je justifie chaque étape.
\end{itemize}
\end{aretenir}

\findecours

\end{document}
